\documentclass[11pt]{article}
\usepackage[english]{babel}
\usepackage[a4paper,margin=0.9in]{geometry}
\usepackage[T1]{fontenc}
\usepackage[utf8]{inputenc}
\usepackage{lmodern}
\usepackage{amsmath,amssymb,amsthm,mathtools,mathrsfs}
\usepackage{microtype}
\usepackage{fancyhdr}
\usepackage{array,longtable,enumitem,multicol}
\usepackage{tikz}

\providecommand{\dd}{\,d}
\providecommand{\tht}{\vartheta}
\providecommand{\vark}{\varkappa}

\providecommand{\sn}{\operatorname{sn}}
\providecommand{\cn}{\operatorname{cn}}
\providecommand{\dn}{\operatorname{dn}}
\providecommand{\am}{\operatorname{am}}
\providecommand{\ctg}{\operatorname{cotg}}
\providecommand{\eps}{\varepsilon}
\providecommand{\Jac}[2]{\left(\frac{#1}{#2}\right)}
\providecommand{\ii}{\mathrm{i}}
\providecommand{\cotg}{\operatorname{cotg}}

\setlength{\parindent}{1.5em}
\setlength{\parskip}{0.18em}
\emergencystretch=6em
\pagestyle{fancy}
\fancyhf{}
\cfoot{\thepage}
\lhead{Heinrich Weber, Lehrbuch der Algebra, Vol. II}
\rhead{English Translation}

\newcommand{\Om}{\Omega}
\newcommand{\Omm}{\Omega_m}
\newcommand{\Hm}{H_m}
\newcommand{\Lm}{\Omega_m}
\newcommand{\Norm}{\mathrm{N}}
\newcommand{\Normone}{\mathrm{N}_1}
\newcommand{\frp}{\mathfrak p}
\newcommand{\frP}{\mathfrak P}
\newcommand{\frG}{\mathfrak G}
\newcommand{\frA}{\mathfrak A}
\newcommand{\frB}{\mathfrak B}
\newcommand{\Gg}{\mathfrak G}
\newcommand{\Aa}{\mathfrak A}
\newcommand{\Bb}{\mathfrak B}
\newcommand{\pp}{\mathfrak p}
\providecommand{\eps}{\varepsilon}
\newcommand{\G}{\mathfrak G}
\newcommand{\B}{\mathfrak B}
\newcommand{\A}{\mathfrak A}
\newcommand{\frakA}{\mathfrak a}
\newcommand{\Na}{N_a}
\newcommand{\abs}[1]{\left|#1\right|}
\newcommand{\tg}{\operatorname{tg}}
\newcommand{\idxitem}[1]{\par\hangindent=1.4em\hangafter=1 #1}

\providecommand{\jac}[2]{\left(\frac{#1}{#2}\right)}
\providecommand{\Lzero}{\mathfrak L_0}
\providecommand{\Kgrp}{\mathfrak K}
\begin{document}
\begin{center}{\large Second Volume.}\\[0.4em]{\Large First Book. Groups.}\\[0.4em]\end{center}

\section*{First Section. General Group Theory.}

\section*{\S\,1. Definition of Groups.}

In the first volume, in connection with permutations, we became acquainted with the notion of a group and made important algebraic applications of it. Our next task must now be to formulate this concept, which is so exceedingly important in all more recent mathematics, in a more general way, and to become acquainted with the laws which govern it. We place the following definition at the beginning.

A system $P$ of things, or elements, of any kind becomes a group if the following assumptions are satisfied:

\begin{enumerate}[label=\arabic*.]
\item A rule is given by which, from a first and a second element of the system, a quite definite third element of the same system is derived.
\end{enumerate}
If $a$ is the first element, $b$ the second, and $c$ the third, one writes symbolically
\[
ab=c,\qquad c=ab,
\]
and says that $c$ is composed from $a$ and $b$, and that $a$ and $b$ are the components of $c$.

In this composition the commutative law is in general not assumed; that is, $ab$ may be different from $ba$. On the other hand,
\begin{enumerate}[resume,label=\arabic*.]
\item the associative law is assumed,
\end{enumerate}
that is, if $a,b,c$ are any three elements of $P$, then
\[
(ab)c=a(bc).
\]
From this it follows, by complete induction, that one always obtains the same result if, in any finite row of elements of $P$, $a,b,c,d\ldots$, one first composes two adjacent elements, then again two adjacent elements, and so on until the whole row is reduced to one element. This element is denoted by $abcd\ldots$. Thus, for example,
\[
\begin{aligned}
abcd&=(ab)cd=[(ab)c]d=(ab)(cd)\\
&=a(bc)d=[a(bc)]d=a[(bc)d]\\
&=ab(cd)=(ab)(cd)=a[b(cd)].
\end{aligned}
\]
\begin{enumerate}[resume,label=\arabic*.]
\item It is assumed that, if $ab=ab'$ or $ab=a'b$, then necessarily in the first case $b=b'$, and in the second case $a=a'$.
\end{enumerate}

If $P$ comprises a finite number of elements, then the group is called finite and the number of its elements its degree.

For finite groups 1., 2., 3. imply the following consequence:
\begin{enumerate}[resume,label=\arabic*.]
\item If two of the three elements $a,b,c$ of $P$ are given arbitrarily, then the third can always be determined, and in only one way, so that $ab=c$.
\end{enumerate}
Indeed, if $a,b$ are the given elements, the assertion 4. coincides with 1. If, however, $a$ and $c$ are given, let the second component $b$ in the composite $ab$ run through the whole system $P$, whose degree is $n$. Then by 1. and 3. the elements $ab$ so obtained are all distinct elements of $P$, and since their number is $n$, all elements of $P$, and hence also $c$, must occur among them. The same argument applies if $b$ and $c$ are given, by letting $a$ run through the whole system $P$.

For infinite groups this conclusion is no longer valid. Thus for infinite groups property 4. is also included as a requirement in the definition.

In the permutation groups considered in the first volume the marks of this general concept of a group are easily recognized.

We now draw from this definition some quite general consequences. By 4., for every given $b$ there is an element $e$ in $P$ which satisfies
\begin{equation}
eb=b,
\tag{1}
\end{equation}
and this $e$ is independent of $b$; for from (1) it follows, for every $c$, that
\[
ebc=bc,
\]
and by 4. the element $bc$ may mean any element of $P$. Similarly there is an element $e'$ satisfying, for every $b$,
\begin{equation}
be'=b.
\tag{2}
\end{equation}
But this element $e'$ is not distinct from $e$; for, putting $b=e'$ in (1) and $b=e$ in (2), one obtains
\[
ee'=e',\qquad ee'=e,
\]
and hence $e=e'$.

The element $e$ changes nothing when it is composed with arbitrary elements of $P$ and is called the identity of the group. In many cases it may, without misunderstanding, simply be denoted by ``1''.

For each element $a$ there is by 4. a definite element $a^{-1}$ satisfying
\begin{equation}
a^{-1}a=e.
\tag{3}
\end{equation}
From (3), (1), and (2) it follows that
\[
a^{-1}aa^{-1}=ea^{-1}=a^{-1}=a^{-1}e,
\]
and consequently, by 3.,
\begin{equation}
aa^{-1}=e.
\tag{4}
\end{equation}
The two elements $a,a^{-1}$ are called opposite or reciprocal to one another.

In special cases the commutative law may also hold for the composition of the elements of a group $P$; that is, for every two elements $a,b$ of the group one may have $ab=ba$.

\begin{enumerate}[resume,label=\arabic*.]
\item Groups which have this property are called commutative groups, or also Abelian groups.
\end{enumerate}

If the elements of two groups
\[
a,b,c,d\ldots \qquad\hbox{and}\qquad a',b',c',d'\ldots
\]
correspond to one another in a one-to-one manner in such a way that whenever $ab=c$ also $a'b'=c'$, then the groups are called isomorphic. The evident theorem holds that two groups isomorphic with a third group are isomorphic with one another. One may therefore collect all mutually isomorphic groups into one class of groups, which is itself again a group, whose elements are the generic concepts obtained by collecting the corresponding elements of the individual isomorphic groups into a common concept. The individual mutually isomorphic groups are then to be regarded as different representatives of one generic concept.

The properties of groups which may come into consideration are of different kinds. They may attach to particular groups and be derived from the nature of the elements of which the group consists, or from the nature of the law of composition. Or they may attach to groups as such and then have to be derived solely from the definition of the group concept. The latter properties are common to all isomorphic groups and may be called invariant properties of the group. If a comparison is allowed, one may recall the difference between metric and projective properties in geometry.

To the first kind of properties, derived from the special nature of the elements, belong for example, in permutation groups, the properties of transitivity and intransitivity, of primitivity and imprimitivity. To the invariant properties belong commutativity or noncommutativity, the degree, the divisors and their index, and the normal divisors. We shall first deal with these invariant properties, for which it is immaterial from which representative of a class of isomorphic groups they are derived.

\section*{\S\,2. The Divisors of Finite Groups.}

As we have already seen in the first volume, it is a particularly important question whether a part of the elements of a group is itself again a group. For this it is necessary and sufficient that every two elements of this part, when composed, again give an element belonging to the same part. This smaller group is then called a divisor of the whole group. The identity element ``1'' by itself forms a group in every group, and is therefore a divisor in every group.

We now restrict ourselves, as almost always in what follows, to finite groups, and suppose that
\[
P=1,a_1,a_2,\ldots,a_{n-1}
\]
is a group of degree $n$, with elements $a_0=1,a_1,a_2,\ldots,a_{n-1}$, and that
\[
Q=1,a_1,a_2,\ldots,a_{v-1}
\]
is a divisor of $P$ of degree $v$. Every such divisor $Q$ must have the fundamental properties of a group and therefore certainly contains the identity element ``1'', and with each of its elements $a_i$ also the opposite element $a_i^{-1}$. If $v<n$, we call $Q$ a proper divisor of $P$. Thus if $b$ is an element of $P$ not contained in $Q$, then the elements
\[
Q_1=b,a_1b,a_2b,\ldots,a_{v-1}b
\]
are all distinct from one another and are all not contained in $Q$; for if, say, $a_i b$ were contained in $Q$, then, since $Q$ is a group, $a_i^{-1}a_i b=b$ would also have to be contained in $Q$, contrary to the assumption.

If the whole group $P$ is not yet exhausted by $Q$ and $Q_1$, we take one of the remaining elements, which we may denote by $c$, and form
\[
Q_2=c,a_1c,a_2c,\ldots,a_{v-1}c.
\]
Then the elements of $Q_2$ are all distinct not only from one another, but also from the elements of $Q$ and $Q_1$. For if, for example, $a_i c=a_j b$, it would follow that $c=a_i^{-1}a_jb$; but $a_i^{-1}a_j$ is contained in $Q$ and is therefore identical with one of the elements $1,a_1,a_2,\ldots,a_{v-1}$, so that, contrary to the assumption, $c$ would be contained in $Q_1$. We continue in this way, forming the systems $Q,Q_1,Q_2,\ldots,Q_{\mu-1}$, the last of which is
\[
Q_{\mu-1}=g,a_1g,a_2g,\ldots,a_{v-1}g.
\]
Then $P$ must be exhausted by the systems $Q,Q_1,Q_2,\ldots,Q_{\mu-1}$, and since these systems contain only distinct elements, it follows that
\[
n=v\mu,
\]
or the theorem:
\begin{enumerate}[label=\arabic*.]
\item The degree of a group is divisible by the degree of each of its divisors. The quotient $\mu=n:v$ is called the index of the divisor $Q$.
\end{enumerate}
The systems $Q_1,Q_2,\ldots,Q_{\mu-1}$ we call, as in the special case of permutation groups, the side groups belonging to $Q$, and denote them by
\[
Q_1=Qb,
\quad Q_2=Qc,
\quad\ldots\quad
Q_{\mu-1}=Qg,
\]
so that we may also write symbolically
\begin{equation}
P=Q+Qb+Qc+\cdots+Qg.
\tag{1}
\end{equation}

At times we shall also call the whole system $Q,Q_1,\ldots,Q_{\mu-1}$ a system of side groups, and thus call the representation (1) the decomposition of $P$ into a system of side groups.

From the manner in which the side groups are formed it follows further that, if $b,c$ are any two elements of $P$, the two systems $Qb$ and $Qc$ are either wholly identical or have no element in common. For if $a_1b=a_2c$, where $a_1,a_2$ are two elements of $Q$, then for every other element $a$ of $Q$ one also has
\[
aa_1^{-1}b=aa_2^{-1}c,
\]
and when $a$ runs through the whole group $Q$, each of $aa_1^{-1}$ and $aa_2^{-1}$ also runs through the whole group $Q$. Hence $Qb$ and $Qc$ are identical.

As in Vol. I, \S\,154, 2., we may also decompose $P$ into side groups in the following way:
\[
P=Q+bQ+cQ+\cdots+gQ.
\]

The side groups do not have the marks of a group. For two elements of $Q_1$ to give, on composition, again an element of $Q_1$, one would have to have, say,
\[
a_1ba_2b=a_3b,
\]
and hence $b=a_2^{-1}a_1^{-1}a_3$. Thus, contrary to the assumption, $b$ would be contained in $Q$. The name side group is therefore to be understood only improperly.

Two divisors $Q,Q'$ of $P$ always have the element 1 in common. They may, however, also have other elements in common, and these common elements form a group, which we shall denote by $R$. For if the elements $a$ and $b$ belong both to $Q$ and to $Q'$, then, since both $Q$ and $Q'$ are groups, the same is true of the composite $ab$; hence $R$ is a group. This common group we call the greatest common divisor of $Q$ and $Q'$, or, with a geometrical resonance, the intersection of $Q$ and $Q'$. In the same way it follows that the elements common to any number of divisors of $P$ form a group, which we likewise call the greatest common divisor or the intersection of all these groups.

In every group we can form divisors by the following procedure. The identity element by itself is always a divisor.

Denote the repeated composition of an element with itself by powers, and denote the identity element by $a^0$. Then the row of elements
\[
1,a,a^2,a^3,\ldots
\]
all of which belong to $P$, cannot consist entirely of distinct elements. If the first element which occurs for the second time is
\[
a^n=a^{n+\alpha},
\]
then it follows that $a^\alpha=1$, and that $\alpha$ is the least positive number satisfying this condition. Then, if $m=q\alpha$ is any multiple of $\alpha$, one has $a^m=1$; conversely, whenever $a^m=1$, $m$ must be divisible by $\alpha$. For otherwise one could write $m=q'\alpha+\alpha'$, with $\alpha'$ positive and less than $\alpha$, and then $a^{\alpha'}=1$, contrary to the assumption that $\alpha$ is the least positive exponent for which $a^\alpha=1$.

One has
\[
a^\mu=a^{\mu'}
\]
always and only when $\mu\equiv\mu'\pmod{\alpha}$; here the exponents may also be negative, if $a^{-\nu}$ denotes the $\nu$-th power of the element $a^{-1}$, or the element opposite to $a^\nu$.

The row
\[
A=1,a,a^2,\ldots,a^{\alpha-1}
\]
then contains $\alpha$ mutually distinct elements of $P$, in which all powers of $a$ are represented. The elements $A$ plainly form a group of degree $\alpha$, since the exponents simply add under composition. This group is a divisor of $P$; hence $\alpha$ is a divisor of $n$. Thus for every element $a^n=1$.

The group $A$ is called the period of the element $a$, and $\alpha$ is also called the degree of the element $a$.

\section*{\S\,3. Normal Divisors of a Group.}

Let $P$ be a group as above, let $Q$ be a divisor of $P$, and let $b$ be an element of $P$ not contained in $Q$. Then the side group $Qb$ is not a group. On the other hand the system $b^{-1}Qb$, or written out,
\[
1,\;b^{-1}a_1b,\;b^{-1}a_2b,\ldots,\;b^{-1}a_{v-1}b
\]
certainly forms a group, since
\[
b^{-1}a_1b\cdot b^{-1}a_2b=b^{-1}a_1a_2b,
\]
and this group is isomorphic with $Q$. The group $b^{-1}Qb$ is called the group transformed by $b$. If $b$ itself belongs to $Q$, then $b^{-1}Qb$ is identical with $Q$. If $b$ is taken to be the different elements of $P$, one obtains a whole family of such groups, which we call the divisors of $P$ conjugate to $Q$, or briefly conjugate groups.

It may happen that all conjugate divisors are identical with one another. We have already seen, in the case of permutation groups, that this case is of special importance, and therefore introduce the following general definition.

If $Q$ is a divisor of $P$ which is identical with all its conjugate divisors, then $Q$ is called a normal divisor of $P$.

The group formed from the single element 1 is a normal divisor of every group. We obtain another normal divisor in the greatest common divisor $R$ of all divisors conjugate, in $P$, to any divisor $Q$.

Indeed, it was proved above that $R$, as the intersection of several divisors of $P$, is a group. If
\begin{equation}
Q,Q',Q'',\ldots
\tag{1}
\end{equation}
are the divisors of $P$ conjugate to $Q$, and if $b$ is any element of $P$, then the system of groups
\begin{equation}
b^{-1}Qb,
\quad b^{-1}Q'b,
\quad b^{-1}Q''b,
\ldots
\tag{2}
\end{equation}
is not different from the system (1). If $R$ is the intersection of the groups (1), then $b^{-1}Rb$ is the intersection of (2), and consequently $R$ is identical with $b^{-1}Rb$; that is, $R$ is a normal divisor of $P$.

If $N$ is a normal divisor of any group $P$ and $b$ is an arbitrary element of $P$, then
\begin{equation}
b^{-1}Nb=N\quad\hbox{or}\quad Nb=bN.
\tag{3}
\end{equation}
If $P$ has degree $n$, if $N$ has degree $v$ and index $\mu$, so that $n=\mu v$, then the elements $1,b_1,b_2,\ldots,b_{\mu-1}$ may be chosen so that the decomposition of $P$ into side groups gives
\[
P=N+Nb_1+Nb_2+\cdots+Nb_{\mu-1}=N+b_1N+b_2N+\cdots+b_{\mu-1}N.
\]
Thus
\begin{equation}
N_0=N,
\quad N_1=Nb_1=b_1N,
\quad N_2=Nb_2=b_2N,
\ldots,
\quad N_{\mu-1}=Nb_{\mu-1}=b_{\mu-1}N
\tag{4}
\end{equation}
is the system of side groups.

If $a$ is any element in $N$, then $Na$ is identical with $N$; hence $Nab$ is also identical with $Nb$. Since every element $c$ of $P$ occurs in $N$ or in one of its side groups, and so is contained in the form $ab_k$, it follows that $Nc=cN$ must coincide with one of the systems (4).

We also mention the following theorem, whose correctness follows immediately from the isomorphism of transformed groups. If $Q$ is a divisor of $P$ and $R$ is a divisor of $Q$, then, if $Q'$ and $R'$ are obtained from $Q$ and $R$ by transformation with the same element of $P$, $R'$ is also a divisor of $Q'$; and if $R$ is a normal divisor of $Q$, then $R'$ is also a normal divisor of $Q'$.

\section*{\S\,4. Composition of Parts.}

If $A$ and $B$ are any two rows of elements of a group $P$--whether groups or not--we shall understand by the symbolic product $AB$ the totality of all elements obtained by composing an element $a$ of $A$ with an element $b$ of $B$, according to the rule prevailing in $P$, to form $ab$. We call this kind of composition of $A$ and $B$ into $AB$ the composition of parts of $P$. We distinguish here between a part and a divisor of $P$, so that a divisor is always to be a group, whereas this is not necessary for a part.

In the same way one may compose three or more parts $A,B,C,\ldots$ of $P$, and the associative law holds for the composition of parts, as follows immediately from the fact that this law holds in $P$. Thus the meaning of a composite $ABC\ldots$ is uniquely determined.

In the composition $AB$ one of the parts $A,B$ may also consist of a single element, and then the notation $AB$ agrees with the notation used above for side groups.

For the composition of parts the following theorems hold:

\begin{enumerate}[label=\arabic*.]
\item The necessary and sufficient condition that a part $A$ of $P$ be a group is the equation
\begin{equation}
AA=A.
\tag{1}
\end{equation}
\end{enumerate}
For this equation says nothing other than that, if $a$ and $a'$ are contained in $A$, then $aa'$ also belongs to $A$, and hence that $A$ is a group.

\begin{enumerate}[resume,label=\arabic*.]
\item The two equations
\begin{equation}
AA=A,
\tag{2}
\end{equation}
\begin{equation}
AB=BA,
\tag{3}
\end{equation}
where $A$ is fixed and $B$ may be any arbitrary part of $P$, express that $A$ is a normal divisor of $P$.
\end{enumerate}
Indeed, by (2) $A$ is a divisor of $P$, and by (3), for each element $b$ of $P$, one has $b^{-1}Ab=A$; that is, $A$ is a normal divisor.

\begin{enumerate}[resume,label=\arabic*.]
\item If $A$ and $B$ are two divisors of $P$ (groups), then either of the two equations
\begin{equation}
AB=A,
\qquad BA=A,
\tag{4}
\end{equation}
is the necessary and sufficient condition that $B$ be a divisor of $A$.
\end{enumerate}
For if $B$ is a group and a divisor of the group $A$, and if $a$ is contained in $A$ and $b$ in $B$, hence also in $A$, then $ab$ is also contained in $A$. Thus $A$ contains every element of $AB$. Since secondly $B$, as a group, contains also the element 1, $AB$ contains every element of $A$; hence $AB$ is identical with $A$. In the same way one sees that $BA$ is identical with $A$. Conversely, from either of the equations (4), since $A$ as a group contains the identity element, it follows that every element of $B$ is contained in $A$, and hence that $B$ is a divisor of $A$.

\begin{enumerate}[resume,label=\arabic*.]
\item Under the composition of parts, the system of side groups belonging to a normal divisor of $P$ itself forms a group, in which the normal divisor $N$ is the identity, and $Nb_k,Nb_k^{-1}$ are opposite elements.
\end{enumerate}
For let $N$ be a normal divisor of $P$ and let
\begin{equation}
N,
\quad N_1=Nb_1,
\quad N_2=Nb_2,
\ldots
\tag{5}
\end{equation}
be the system of side groups. By the property of normal divisors $N$ is commutable with every $b_k$, hence $b_hb_kN=Nb_hb_k$, and consequently
\[
N_hN_k=NNb_hb_k.
\]
Since $b_hb_k$ is contained in $P$, $Nb_hb_k$ is identical with one of the side groups (5); and since $NN=N$ may be put, it follows, if we put $Nb_hb_k=N_i$, that
\begin{equation}
N_hN_k=N_i.
\tag{6}
\end{equation}
This proves the property \S\,1, 1. of groups for the composition of side groups. We have already emphasized that \S\,1, 2., the associative law, also holds. It remains to prove \S\,1, 3.; that is, if
\begin{equation}
N_hN_i=N_kN_l
\tag{7}
\end{equation}
is put, then from $N_h=N_k$ it follows that also $N_i=N_l$, and from $N_i=N_l$ it follows that also $N_h=N_k$.

According to the representation (5), however, we may write (7) in the two ways
\begin{equation}
Nb_hb_i=Nb_kb_l,
\qquad b_hb_iN=b_kb_lN.
\tag{8}
\end{equation}
If now $N_i=N_l$, then we may also assume $b_i=b_l$, and from (8), by composition with $b_i^{-1}$, we obtain
\[
Nb_h=Nb_k.
\]
If $Nb_h=Nb_k$, then we assume $b_h=b_k$, and obtain from (8), by composition with $b_h^{-1}$,
\[
b_iN=b_lN,
\]
and hence also $N_i=N_l$.

It is thereby proved that the system of side groups becomes a group under the composition of parts. That in this group the normal divisor $N$ is the identity element follows immediately from 1., since then $NN_k=NNb_k=N_k$. And since $Nb_kNb_k^{-1}=Nb_kb_k^{-1}=N$, the elements $Nb_k$ and $Nb_k^{-1}$ are opposite elements.

The group of the quantities $N_k$, whose existence has thereby been proved, is called the group of side groups, or the group complementary to $N$ with respect to $P$. Following Holder, we denote it by $P/N$. The degree of the complementary group is equal to the index of the group $N$.

At the end of these considerations we mention one more theorem on normal divisors.
\begin{enumerate}[resume,label=\arabic*.]
\item If $N$ is a divisor of $P$, $N'$ a divisor of $N$, and at the same time a normal divisor of $P$, then $N'$ is also a normal divisor of $N$.
\end{enumerate}
This is self-evident, since for every element $b$ of $P$ one has $b^{-1}N'b=N'$, and since every element of $N$ is also an element of $P$.

But one must not conversely conclude that, if $N$ is a normal divisor of $P$ and $N'$ is a normal divisor of $N$, then $N'$ must also be a normal divisor of $P$.
\clearpage
\section*{\S\,176. The real field $H_m$ contained in $\Omega_m$.}

Every number of the field $\Omega_m$ is transformed by the substitution $(r,r^{-1})$ into the conjugate-imaginary number which is also obtained by the interchange $(i,-i)$. The product of two such conjugate-imaginary numbers is the square of the absolute value of either of them.

A number in $\Omega_m$ is real if it remains unchanged under the interchange $(r,r^{-1})$, and only under this condition. Every such number can be expressed rationally by the two-term period $r+r^{-1}$, and therefore belongs to a real field of degree $\frac12\varphi(m)$ which is a divisor of $\Omega_m$. This field we shall call the real field contained in $\Omega_m$ and denote it by $H_m$ (although other real fields are also contained in $\Omega_m$, namely all divisors of $H_m$).

The $\frac12\varphi(m)=\frac12\mu$ numbers
\begin{equation}
1,\quad r+r^{-1},\quad r^2+r^{-2},\ldots,\quad r^{\frac12\mu-1}+r^{-\frac12\mu+1}
\tag{1}
\end{equation}
form a basis of the algebraic integers of the field $H_m$. For by \S\,170, (7),
\begin{equation}
\begin{aligned}
\omega={}&x_0+x_1r+x_2r^2+\cdots+x_{\frac12\mu-1}r^{\frac12\mu-1}+x_{\frac12\mu}r^{\frac12\mu} \\
&+x_{-1}r^{-1}+x_{-2}r^{-2}+\cdots+x_{-\frac12\mu+1}r^{-\frac12\mu+1}
\end{aligned}
\tag{2}
\end{equation}
is an algebraic integer of the field $\Omega_m$ if and only if the coefficients $x_0,x_1,
\ldots$ are rational integers.

The condition that this number should belong to the field $H_m$ is obtained by carrying out the substitution $(r,r^{-1})$ and setting the difference equal to zero. If one then also divides by $r-r^{-1}$, one obtains
\begin{align*}
0={}&(x_1-x_{-1})+(x_2-x_{-2})(r+r^{-1})+\cdots \\
&+(x_{\frac12\mu-1}-x_{-\frac12\mu+1})\frac{r^{\frac12\mu-1}-r^{-\frac12\mu+1}}{r-r^{-1}}
+x_{\frac12\mu}\frac{r^{\frac12\mu}-r^{-\frac12\mu}}{r-r^{-1}}.
\end{align*}
The divisions can be carried out here; and if one then multiplies by $r^{\frac12\mu-1}$, there results for $r$ a rational equation of degree lower than the $\mu$th, which can be satisfied only if all its coefficients vanish. This leads to the equations
\[
x_{\frac12\mu}=0,\qquad x_{\frac12\mu-1}=x_{-\frac12\mu+1},\ldots,\qquad x_1=x_{-1},
\]
and hence every algebraic integer $\omega$ of the field $H_m$ is contained in the form
\[
\omega=x_0+x_1(r+r^{-1})+x_2(r^2+r^{-2})+\cdots+x_{\frac12\mu-1}(r^{\frac12\mu-1}+r^{-\frac12\mu+1}).
\]
Instead of the basis (1), if one sets
\[
\rho=r+r^{-1},
\]
one may also choose the powers of $\rho$:
\begin{equation}
1,\quad \rho,\quad \rho^2,\ldots,\quad \rho^{\frac12\mu-1}
\tag{3}
\end{equation}
as a basis of the algebraic integers of $H_m$. For the quantities (3) can be expressed integrally in terms of the quantities (1), and conversely.

The number $\rho$ is the root of an irreducible equation of degree $\frac12\mu$, which we shall denote by $\psi(\rho)=0$. If $f(r)=0$ is the equation for $r$ (\S\,169), then the identical relation
\begin{equation}
f(x)=x^{\frac12\mu}\psi\left(x+\frac1x\right)
\tag{4}
\end{equation}
holds; differentiating gives
\begin{equation}
f'(r)=r^{\frac12\mu}\psi'(\rho)(1-r^{-2}).
\tag{5}
\end{equation}
From this we can form the fundamental number $\Delta_1$ of the field $H_m$, which, since $H_m$ contains only real numbers, must be positive. If we denote by $\Normone$ the norm with respect to $H_m$, this fundamental number is
\[
\Delta_1=\pm\Normone\psi'(\rho).
\]
But if $\Norm$ is the norm formed with respect to $\Omega_m$, then
\[
\Norm\psi'(\rho)=[\Normone\psi'(\rho)]^2=\Delta_1^2,
\]
and therefore formula (5), with regard to \S\,169, gives
\[
\Delta=\Norm(1-r^{-2})\Delta_1^2.
\]
By \S\,169, (7) and (10),
\begin{equation}
\begin{aligned}
\Norm(1-r^{-2})&=q &&\text{for odd }q,\\
&=4 &&\text{for }q=2,\\
\pm\Delta&=q\Delta_1^2 &&\text{for odd }q,\\
&=4\Delta_1^2 &&\text{for }q=2.
\end{aligned}
\tag{6}
\end{equation}
If one substitutes therein
\[
\Delta=\pm q^{q^{\kappa-1}[\kappa(q-1)-1]},
\]
it follows that
\begin{equation}
\begin{aligned}
\Delta_1&=q^{\,q^{\kappa-1}\frac{q-1}{2}-\frac{q^{\kappa-1}+1}{2}} &&\text{for odd }q,\\
&=2^{(\kappa-1)2^{\kappa-2}-1} &&\text{for }q=2,
\end{aligned}
\tag{7}
\end{equation}
so that $\Delta_1$ is always a power of $q$.

\section*{\S\,177. The prime ideals in the field $H_m$.}

We now have to examine the question of the relation between the prime ideals of the field $H_m$ and those of the field $\Omega_m$.

If $\frP$ denotes any prime ideal in $H_m$ of degree $f_1$, then $\frP$ must be a divisor of a natural prime number $p$, and therefore (by \S\,171) $\frP$ can be divisible by the second or a higher power of a prime factor $\frp$ in $\Omega_m$ only if $p=q$.

If, however, $\frP$ is a divisor of $q$, then it must be a power of $\sigma=1-r$. Put, for instance,
\[
\frP=\sigma^\lambda.
\]
Taking the norm of this with respect to $\Omega_m$, which is the square of the norm with respect to $H_m$, one obtains (by \S\,169)
\[
q^{2f_1}=q^\lambda,
\]
and hence $f_1=\frac12\lambda$. Since $f_1$ is an integer, $\lambda$ must be at least equal to $2$. But $\lambda$ also cannot be greater than $2$, since $\sigma^2$ is associated with the algebraic integer $(1-r)(1-r^{-1})$ contained in $H_m$, and therefore $\sigma^2$ must be divisible by $\frP$. Consequently $f_1=1$, and we obtain the theorem:

\begin{enumerate}[label=\arabic*.,leftmargin=2.2em]
\item The prime number $q$ is the $\frac12\mu$th power of a prime number of first degree existing in $H_m$.
\end{enumerate}

Now let $p$ be different from $q$, and let $\frp$ be a prime factor of $\frP$ in the field $\Omega_m$, of degree $f$, which is carried by the substitution $(r,r^{-1})$ into $\frp'$. Then $\frP$ is divisible both by $\frp$ and by $\frp'$; on the other hand, $\frp\frp'$ is contained in $H_m$ and is therefore divisible by $\frP$. Thus we have to distinguish whether $\frp$ and $\frp'$ are distinct or not.

If $\frp$ is distinct from $\frp'$, then $\frP$ is divisible by $\frp\frp'$, and it follows that $\frP=\frp\frp'$; taking norms gives $f_1=f$.

If, however, $\frp=\frp'$, then $\frP=\frp$, and taking norms gives
\[
2f_1=f.
\]
Which of these two cases occurs depends, therefore, on whether the prime ideal $\frp$ admits the substitution $(r,r^{-1})$ or not, that is, whether $(r,r^{-1})$ is contained in the group to which the prime ideal belongs or not. By \S\,172 this difference is determined by whether there is some exponent $h$ for which the congruence
\[
p^h\equiv -1\pmod m
\]
is satisfied, or whether there is no such exponent.

We summarize the result as follows:

\begin{enumerate}[label=\arabic*.,leftmargin=2.2em,start=2]
\item The natural prime numbers $p$ different from $q$ split into two kinds, $p_1,p_2$.

The first kind $p_1$ is characterized by the condition that for some exponent $h$
\[
p_1^h\equiv -1\pmod m,
\]
and these prime numbers, if they belong to the exponent $f$ modulo $m$ and if $\mu=ef$, split in the field $H_m$ into $e$ prime factors of degree $\frac12f$. Their prime factors are not further decomposable in $\Omega_m$ either.

The prime numbers of the second kind $p_2$ are determined by the condition that none of their powers is congruent to $-1$ modulo $m$; they split in the field $H_m$ into $\frac12e$ prime factors of degree $f$. Their prime factors are still decomposable in $\Omega_m$ into two factors each.
\end{enumerate}

For prime numbers of the first kind, $f$ must certainly be even. If $m$ is odd, it is also sufficient, in order that $p$ be a prime number of the first kind, that it belong to an even exponent; for then
\[
(p^{\frac12 f}-1)(p^{\frac12 f}+1)\equiv0\pmod m.
\]
The first of these factors, $p^{\frac12 f}-1$, is not divisible by $m$, since otherwise $p$ would belong to the exponent $\frac12f$. Consequently $p^{\frac12 f}+1$ must be divisible by $q$. But here both factors $p^{\frac12 f}-1$ and $p^{\frac12 f}+1$ cannot be divisible by $q$, since otherwise their difference, which is equal to $2$, would also be divisible by $q$. Hence $p^{\frac12 f}+1$ is divisible by $m$.

If, however, $m$ is a power of $2$, then the congruence
\[
p^h\equiv -1\pmod m
\]
is possible only if $p\equiv -1\pmod m$. For every even power of an odd number is $\equiv+1\pmod 8$ and therefore cannot be $\equiv-1\pmod m$. But if $h$ is odd, then
\[
\frac{p^h+1}{p+1}=p^{h-1}-p^{h-2}+\cdots+1
\]
is itself odd, and consequently $p^h+1$ is divisible by no higher power of $2$ than $p+1$. Thus we may add to Theorem 2 the following more precise determination:

\begin{enumerate}[label=\arabic*.,leftmargin=2.2em,start=3]
\item If $m$ is odd, then the prime numbers $p_1$ belong to an even exponent, and the prime numbers $p_2$ to an odd exponent.

If $m$ is even, then the prime numbers of the form $km-1$ belong to the first kind, and all other odd prime numbers to the second kind.
\end{enumerate}

\section*{\S\,178. The units of the field $H_m$.}

The units belonging to the field $H_m$ are of special importance for the applications. The units of the field $\Omega_m$ may also be reduced to them, according to the following theorem:

\begin{enumerate}[label=\arabic*.,leftmargin=2.2em]
\item Every unit of the field $\Omega_m$ is the product of a root of unity and a real unit of the field $H_m$.
\end{enumerate}

Denote any unit of the field $\Omega_m$ by $\mathscr E(r)$. Then $\mathscr E(r^{-1})$ is the conjugate-imaginary value, which likewise represents a unit, and the quotient $\mathscr E(r):\mathscr E(r^{-1})$ is again a unit whose absolute value
\begin{equation}
\sqrt{\frac{\mathscr E(r)}{\mathscr E(r^{-1})}\frac{\mathscr E(r^{-1})}{\mathscr E(r)}}
\tag{1}
\end{equation}
is equal to $1$. Consequently this quotient is a root of unity and hence is $=\pm r^h$, where $h$ is some integral exponent (\S\,175, Theorems 1 and 2). We therefore have
\begin{equation}
\mathscr E(r)=\pm r^h\mathscr E(r^{-1}).
\tag{2}
\end{equation}

In the rest of the proof a small distinction has now to be made according as $m$ is odd or even.

First suppose that $m$ is odd. Then in (2) we may take the exponent $h$ to be even, since, if necessary, we may replace it by $h+m$. Moreover, in formula (2) only the upper sign is admissible. For if the lower sign stood there, it would follow that
\[
\mathscr E(r)\equiv\mathscr E(1)\equiv-\mathscr E(1),\qquad 2\mathscr E(1)\equiv0\,[\mathrm{mod}\,(1-r)].
\]
But by \S\,169, $1-r$ is a divisor of $q$, and so is relatively prime to $2$; therefore $\mathscr E(1)$, and consequently also $\mathscr E(r)$, is divisible by $(1-r)$. This contradicts the supposition that $\mathscr E(r)$ is to be a unit.

Consequently (2) gives
\[
r^{-\frac12 h}\mathscr E(r)=r^{\frac12 h}\mathscr E(r^{-1}),
\]
whence it follows that $r^{-\frac12 h}\mathscr E(r)$ is a real unit. Denoting it by $e(r)$, we therefore have
\begin{equation}
\mathscr E(r)=r^{\frac12 h}e(r),
\tag{3}
\end{equation}
which proves Theorem 1 in this case.

If, secondly, $m$ is a power of $2$, then $\mu=\frac12m$ and $r^\mu=-1$; and along with $r$, $-r$ is also an $m$th root of unity. If necessary replacing $h$ by $h+\frac12m$, we may assume the upper sign in (2) and write
\begin{equation}
\mathscr E(r)=r^h\mathscr E(r^{-1}).
\tag{4}
\end{equation}
The point is now to prove that $h$ must be even.

If we represent $\mathscr E(r)$ in the basis $1,r,r^2,\ldots,r^{\mu-1}$, then
\begin{equation}
\mathscr E(r)=\sum_{s=0}^{\mu-1}x_sr^s,
\tag{5}
\end{equation}
where the $x_s$ are rational integers. But by (4) we then have
\[
\sum_{s=0}^{\mu-1}x_sr^s=\sum_{s=1}^{\mu-1}x_sr^{h-s},
\]
and from this it follows that $x_s=\pm x_{s'}$ whenever $s+s'\equiv h\pmod\mu$. If $h$ is odd, then of two numbers so connected one is even and the other odd (since $\mu$ is a power of $2$). Thus the expression (5) gives for $\mathscr E(r)$ an aggregate of terms of the form
\[
x_s(r^s\pm r^{s'}).
\]
But
\[
r^s\pm r^{s'}=r^s(1\pm r^{s'-s})=r^s(1\mp r^{\mu+s'-s})
\]
is always divisible by $1-r$, and from this it would follow that $\mathscr E(r)$ is divisible by $1-r$, which is not a unit, contrary to the concept of a unit. Therefore the assumption that the exponent $h$ in formula (4) is odd is inadmissible, and, as above, we may put
\[
\mathscr E(r)=r^{\frac12h}e(r),
\]
where $e(r)$ is a real unit. Theorem 1 is thereby proved.

We shall single out a system of real units. By \S\,169 the $\mu$ quantities $1-r^n$ are all associated with one another. Therefore the quotient
\[
\frac{1-r^n}{1-r^{n'}},
\]
where $n,n'$ are two numbers relatively prime to $m$, is a unit. From this, however, setting
\[
r=e^{\frac{2\pi i}{m}},
\]
one derives the real unit
\begin{equation}
\frac{\sin\frac{n\pi}{m}}{\sin\frac{n'\pi}{m}}
=\frac{r^{\frac n2}-r^{-\frac n2}}{r^{\frac{n'}2}-r^{-\frac{n'}2}}
=r^{\frac{n'-n}{2}}\frac{1-r^n}{1-r^{n'}}.
\tag{6}
\end{equation}
If $m$ is even, one may set $n'=n+\frac m2$, and obtains in this case the real units
\begin{equation}
\tau_n=\operatorname{tang}\frac{n\pi}{m},
\tag{7}
\end{equation}
where $n$ may denote any odd number.

Replacing $n$ by $\frac12m-n$ carries $\tau_n$ into its reciprocal value. If $n$ runs through the sequence of numbers
\[
1,3,5,\ldots,\frac m4-1,
\]
then all the values $\tau_n$ obtained are positive proper fractions.

\clearpage
\section*{Twenty-second Section. Abelian fields and cyclotomic fields.}

\section*{\S~179. Simple Abelian fields.}

We now turn to one of the most interesting applications of the theory of algebraic numbers.

The matter is the proof of the general theorem:\footnote{Kronecker first stated this theorem, but did not publish a complete proof of it. The first proof was made known by the author of the present work in vol.~8 of \textit{Acta Mathematica} (1886). Very recently a second proof by Hilbert has been published, which is based on the theorems developed in the nineteenth section (\textit{Nachrichten d. Gesellschaft d. Wissenschaften zu Göttingen}, 1896, issue 1).}

\textbf{I. All Abelian number fields in the absolute domain of rationality are cyclotomic fields.}

Once we have proved this theorem, the investigations of the third section on cyclotomic fields acquire an increased interest, because it is thereby shown that on the path indicated there one finds not only all cyclotomic fields, but all Abelian fields whatever.

Let
\[
  \Om=R(x)
\]
be an Abelian field of degree $m$, that is, a field consisting of all rational functions of a root $x$ of an irreducible Abelian equation of degree $m$, when the field $R$ of rational numbers is regarded as the domain of rationality. The Galois group $\Gg$ of this equation, which is therefore an Abelian group and has the same degree $m$ as the field $\Om$, is also the group of the field $\Om$.

If $x_1$ is a non-primitive number from $\Om$, then $\Om_1=R(x_1)$ is likewise an Abelian field, which we call a proper divisor of $\Om$, and whose degree is a proper divisor of the degree of $\Om$. For if $\Aa$ is the group to which the number $x_1$ belongs, then $\Gg\mid\Aa$ (\S~4 of this volume) is the group of the field $\Om_1$, and this is, together with $\Gg$, an Abelian group. A divisor of $\Om$ which has the same degree as $\Om$ is identical with $\Om$ (cf. Vol.~I, \S\S~144, 149, 156).

If there are now two non-primitive numbers $x_1,x_2$ in $\Om$ of such a kind that
\[
  \Om=R(x_1,x_2)
\]
can be put, then the fields $\Om_1=R(x_1)$, $\Om_2=R(x_2)$ are proper divisors of $\Om$, and $\Om$ is called \emph{composed} from the two fields $\Om_1,\Om_2$.

If the field $\Om$ permits no such representation, it is called \emph{simple}.

Thus, if it is assumed as proved that $\Om_1,\Om_2$ are cyclotomic fields, that is, that all their numbers are representable rationally by roots of unity, then the same follows for $\Om$.

If one of the fields $\Om_1,\Om_2$ is decomposable, then we can repeat the decomposition, and, since the degrees are decreasing positive integers, we must finally arrive at simple fields.

Theorem I therefore needs to be proved only for simple Abelian fields.

It is, however, to be recalled that the decomposability of a field $\Om$ by no means already follows from the existence of a divisor, and that in a decomposable field the decomposition into simple fields can occur in quite different ways, so that for fields the laws of divisibility do not hold as they do in the domain of numbers.

A criterion for the simplicity of a field can be derived from its group.

If the group $\Gg$ has two proper divisors $\Aa$ and $\Bb$ of such a kind that every element can be composed in one and only one way from an element of $\Aa$ and an element of $\Bb$, then we call the group $\Gg$ decomposable into the two components $\Aa,\Bb$; thus $\Gg$ is decomposable when, in the product formed according to the composition of parts (\S~4),
\begin{equation}
  \Gg=\Aa\Bb,
\tag{1}
\end{equation}
each element of $\Gg$ appears once and only once. The degree of $\Gg$ is then equal to the product of the degrees of $\Aa$ and $\Bb$.

If there are no such divisors, the group $\Gg$ is called indecomposable.

Then we can prove the theorem:

\textbf{1. If the group of an Abelian field is decomposable, then the field too is composed.}

To prove this, assume that the group $\Gg$ of the field $\Om$ is decomposable according to formula (1), and let $m,a,b$ be the degrees of $\Gg,\Aa,\Bb$. We take a number $\xi$ of the field $\Om$ belonging to the group $\Bb$, which is therefore transformed by the substitutions of $\Gg$ into $a$ different conjugate values, and likewise a $b$-valued number $\eta$ belonging to $\Aa$. We can then form a number
\[
  x=\alpha\xi+\beta\eta
\]
with rational numerical coefficients $\alpha,\beta$, which has $m$ different conjugate values, and then $\Om=R(x)$. Thus $\Om$ is composed from $R(\xi)$ and $R(\eta)$ (Vol.~I, \S~143).

If we now recall the representation of an Abelian group by a basis derived in \S\S~9, 10 of this volume, then repeated application of what has just been proved gives:

\textbf{2. Every Abelian field $\Om$ can be composed from such fields whose degree is a power of a prime number and whose group is cyclic. The degrees of these composing fields are the invariants of the group of $\Om$.}

These theorems are supplemented by the following:

\textbf{3. An Abelian field of prime-power degree with cyclic group is simple.}

If the group $\Gg$ is cyclic, then its elements can be represented by the powers of a basis, $G^\gamma$, when $\gamma$ runs through a complete system of residues modulo $m$. We obtain all divisors $\Aa$ of $\Gg$ represented by
\[
  \Aa=G^{b\gamma},
\]
when $b$ is a divisor of $m$ and $\gamma$ runs through a complete system of residues modulo $a=m:b$.

If
\[
  \Aa'=G^{b'\gamma}
\]
is a second divisor of $\Gg$ of degree $a'$ and if
\[
  b'\le b,
  \qquad a'\ge a,
\]
then, when we now assume that $m$ and hence also $a,b,a',b'$ are powers of a prime number $q$, $b'$ is a divisor of $b$, and consequently $\Aa$ is a divisor of $\Aa'$.

If therefore $\xi$ is a number of the field $\Om$ which belongs to the group $\Aa$, then $\xi$ is a primitive number of the field $R(\xi)$ of degree $b$, and in this field every number $\xi'$ is also contained which admits the substitutions of the group $\Aa'$. Thus if $\Aa,\Bb$ and $\Aa',\Bb'$ are proper divisors of $\Gg$, then $R(\xi)$ is different from $\Om$, and $R(\xi)$ is not enlarged by the adjunction of $\xi'$. Hence $\Om$ also cannot be composed from $R(\xi)$ and $R(\xi')$, whereby 3 is proved.

\section*{\S~180. The resolvents.}

According to what has now been proved, in the further considerations, whose aim is the proof of Theorem I, we may restrict ourselves to such Abelian fields whose degree is a prime power and whose group is cyclic. For this reason it also sufficed, for the intended application, that in the preceding section we restricted ourselves to the consideration of the cyclotomic fields $\Omega_m$ in which $m$ was a prime power. We retain as much as possible of the notation used there and put
\begin{equation}
  m=q^\kappa,
\tag{1}
\end{equation}
where $q$ is a prime number and $\kappa$ a positive exponent; and in the case where $q=2$ we assume $\kappa>1$. By $n$ we denote any number not divisible by $q$. Let $r$ be a primitive $m$-th root of unity, and $\Omega_m$ the field of rational functions of $r$.

Let now $\Om=R(x)$ be a simple Abelian field of degree $m$, so that $x$ is a root of a rational irreducible cyclic equation of degree $m$. We denote the roots of this equation, in a suitable order, by
\begin{equation}
  x_0,x_1,x_2,\ldots,x_{m-1},
\tag{2}
\end{equation}
and assume $x_m=x_0$, $x_{m+1}=x_1$, and so on. Then
\[
  x_1=\Phi(x_0),\quad x_2=\Phi(x_1),\ldots,\quad x_0=\Phi(x_{m-1}),
\]
where $\Phi$ denotes an integral function with rational coefficients, and the group of the field $\Om$ consists of the powers of the substitution $(x_0,x_1)$, or of the cyclic permutations

of the roots (2). Every cyclic function of these quantities is a rational number. What we have to prove is that the $x$ can be expressed rationally by roots of unity. If we adjoin to the field $\Om$ the $m$-th root of unity $r$, there arises a field $R(x,r)$ which contains the two fields $\Om=R(x)$ and $R(r)=\Omega_m$ as divisors.

If a number $\omega=\Phi(x,r)$ of the field $R(x,r)$ admits all substitutions $(x_0,x_1)$, $(x_0,x_2)$, \ldots, $(x_0,x_{m-1})$, then it is contained in the field $\Omega_m$; for then
\[
  m\omega=\Phi(x_0,r)+\Phi(x_1,r)+\cdots+\Phi(x_{m-1},r),
\]
and the coefficients of the powers of $r$ in this expression are not only cyclic but even symmetric functions of the roots (2), hence rational numbers. This remains true even in the case where the equation for $x$ is reduced by adjunction of $r$.

We apply this to the resolvents
\begin{equation}
  (r,x_0)=x_0+rx_1+r^2x_2+\cdots+r^{m-1}x_{m-1},
\tag{3}
\end{equation}
and, more generally, when $s$ denotes an arbitrary integral exponent,
\begin{equation}
  (r^s,x_0)=x_0+r^s x_1+r^{2s}x_2+\cdots+r^{(m-1)s}x_{m-1}.
\tag{4}
\end{equation}
By these resolvents $x_0$ itself can again be rationally expressed, namely
\begin{equation}
  m x_0=\sum_{s=0}^{m-1}(r^s,x_0),
\tag{5}
\end{equation}
and if therefore we can show that all these resolvents $(r^s,x_0)$ are cyclotomic numbers, then we have reached our goal. If we put
\begin{equation}
  (r^s,x_\lambda)=x_\lambda+r^s x_{\lambda+1}+r^{2s}x_{\lambda+2}+\cdots+r^{(m-1)s}x_{\lambda+m-1},
\tag{6}
\end{equation}
then $(r^s,x_\lambda)$ arises from $(r^s,x_0)$ by the substitution $(x_0,x_\lambda)$, and the fundamental relation holds:
\begin{equation}
  r^{\lambda s}(r^s,x_\lambda)=(r^s,x_0).
\tag{7}
\end{equation}
All these resolvents are numbers of the field $R(x,r)$; we consider preferably the following two combinations:
\begin{align}
  (r^s,x_0)^m &= \omega_s, \tag{8}\\
  (r^s,x_0)^{-1}(r,x_0)^s &= \alpha. \tag{9}
\end{align}

These two numbers admit, as follows immediately from (7), the substitutions $(x_0,x_\lambda)$ and are therefore numbers of the field $\Omega_m$.

If $n$ is not divisible by $q$, then, because of the irreducibility of the cyclotomic equation, none of the resolvents $(r^n,x_0)$ can vanish unless they all vanish; for the equation
\[
  (r,x_0)^m=0
\]
would be an equation for $r$ with rational coefficients, admitting all substitutions $(r,r^n)$.

Likewise, none of the resolvents $(r^s,x_0)$ can vanish unless the whole family vanishes in which $s$ has a fixed greatest common divisor with $m$.

If all $(r,x_0)$ vanish, formula (9) is not applicable. Then, however, we can directly consider the resolvents $(r^q,x_0)$, which belong to a field of degree $m q^{-1}$ arising from
\[
  y_0=x_0+x_q+x_{2q}+\cdots+x_{m-q},
\]
and if these quantities too vanish, we pass to the resolvents $(r^{q^2},x_0)$, and so on.

If, however, $(r,x_0)$ does not vanish, then by formula (9) all $(r^s,x_0)$ can be represented rationally by $(r,x_0)$ and by cyclotomic numbers.

Consequently it is sufficient for our purpose if we can prove that the non-vanishing resolvents $(r,x_0)$ are cyclotomic numbers.

Let $n$ run through a complete system of incongruent numbers not divisible by $q$. Then
\[
  \sum n\equiv 0\pmod m,
\]
for the numbers $n$ split into pairs which complement one another to $m$. From (7) it follows that the product
\begin{equation}
  \prod_n (r^n,x_0)=a
\tag{10}
\end{equation}
is a number in $\Omega_m$. But since this number moreover is not changed when any of the substitutions $(r,r^n)$ is performed in it, $a$ is a rational number.

\section*{\S~181. Preparation for the proof.}

After the developments of the preceding paragraph, for the proof of the great Theorem I it remains only to show that the number contained in $\Omega_m$,
\[
  \omega=(r,x_0)^m,
\]
is the $m$-th power of a cyclotomic number.

\textbf{1. The essential property of the numbers $\omega_n$ conjugate to $\omega$, on which the proof rests, is that}
\begin{equation}
  \omega_n^{-1}\omega_1^n=\alpha^m
\tag{1}
\end{equation}
\textbf{is the $m$-th power of a number in $\Omega_m$.}

Here $n$ may be any number not divisible by $q$. If we form, of the right- and left-hand sides of (1), the norms in the field $\Omega_m$, then, denoting by $a$ the norm of the number $\alpha$, which is a rational number, and observing that every $\omega_n$ has the same norm as $\omega$, we obtain
\[
  N(\omega)^{n-1}=a^m.
\]
If now $m$ is odd, then we can put $n=2$ in this, and obtain as a consequence of 1:

\textbf{2. The norm of $\omega$ is the $m$-th power of a rational number,}
\begin{equation}
  N(\omega)=a^m.
\tag{2}
\end{equation}
If $m$ is a power of $2$, this inference is no longer permissible. It then follows from 1 only that $N(\omega)^2$ is the $m$-th power of a rational number. Nevertheless, as is seen from (10) of the preceding paragraph, the numbers $\omega$ must, also in this case, still satisfy condition 2; but this condition is then no longer contained automatically in 1, and is added as a new property.

What we have to prove is that from 1 and 2 it follows that $\omega$ itself is the $m$-th power of a cyclotomic number, even if in a higher field.

For this it is necessary to investigate the decomposition of the numbers $\omega$ into their prime factors in the field $\Omega_m$.

The prime number $q$ is, as was proved in \S~169, associated with the $\varphi(m)$-th power of a prime number $\sigma=1-r$ existing in $\Omega_m$.

If therefore $\sigma^k$ is the power of $\sigma$ contained in $\omega$, we put
\[
  \omega=\sigma^k\omega',
\]
where $\omega'$ is relatively prime to the prime number $q$. Formation of the norm gives, by 2,
\[
  a^m=q^k b,
\]
where $b$ is a rational number which, in lowest terms, contains the prime number $q$ neither in the numerator nor in the denominator. From this it follows that $k$ must be divisible by $m$, and hence the first result:

\textbf{3. The exponent of the prime number $\sigma$ contained in $\omega$ is always divisible by $m$.}

Let now $p$ be a prime number different from $q$, and let $\pp_\xi$ be a prime factor of $p$ in the field $\Omega_m$. Let $\xi$ run through the sequence of numbers $\xi_1,\xi_2,\ldots,\xi_e$ determined more precisely in \S~172 in the different cases, so that
\begin{equation}
  p=\pp_{\xi_1}\pp_{\xi_2}\cdots\pp_{\xi_e},
\tag{3}
\end{equation}
and assume that $\pp_\xi^{k_\xi}$ is the power of $\pp_\xi$ contained in $\omega$. The power of $\pp_{n\xi}$ contained in $\omega_n$ is then $\pp_{n\xi}^{k_\xi}$; and if we determine $n'$ from the congruence
\begin{equation}
  nn'\equiv 1\pmod m,
\tag{4}
\end{equation}
then $\pp_\xi^{k_{n'\xi}}$ is the power of $\pp_\xi$ contained in $\omega_n$. Accordingly, in
\begin{equation}
  \omega_n^{-1}\omega_1^n
\tag{5}
\end{equation}
the power $\pp_\xi^{n k_\xi-k_{n'\xi}}$ is contained, and by 1 the exponent of this power must be divisible by $m$.

Thus for every $n$ not divisible by $q$ we obtain the congruence
\begin{equation}
  n k_\xi\equiv k_{n'\xi}\pmod m,
\tag{6}
\end{equation}
in which we can now also replace $\xi$ by $n\xi$, whereby
\begin{equation}
  n k_{n\xi}\equiv k_\xi\pmod m
\tag{7}
\end{equation}
is obtained. If in this we take $\xi=1$ and put $\xi,\xi'$ for $n,n'$ and $k$ for $k_1$, then from (7) follows
\begin{equation}
  k_\xi\equiv \xi' k\pmod m,
\tag{8}
\end{equation}
where $k$ now has the same value for all $\xi$, and $\xi,\xi'$ are connected by the congruence
\[
  \xi\xi'\equiv1\pmod m.
\]

If we take $n=p$, then $\pp_{p\xi}=\pp_\xi$ (\S~172), and so also $k_{p\xi}=k_\xi$; and from (7) and (8) follows
\begin{equation}
  k(p-1)\equiv0\pmod m.
\tag{9}
\end{equation}
Important consequences now follow from this. First, if $p-1$ is not divisible by $q$, then it follows that $k$ must be divisible by $m$, and we therefore have the theorem:

\textbf{4. The prime factors of a prime number $p$ which is not congruent to $1$ modulo $q$ are contained in $\omega$ only in such powers whose exponents are divisible by $m$.}

If, in the case where $m$ is a power of $2$, $p-1$ is divisible by $2$ but not by $4$, that is, if $p\equiv -1\pmod4$, then it follows from (9) only that $k$ is divisible by $\tfrac12m$, not that $k$ is divisible by $m$. The exponent of the power of $\pp_\xi$ contained in $\omega$ is, apart from multiples of $m$, equal to $k$; and if $k$ is divisible by $m$, everything is as in case 4. But if $k$ is divisible only by $\tfrac12m$, not by $m$, so that
\[
  k\equiv \tfrac12m\pmod m,
\]
then $k-\tfrac12m$ is divisible by $m$, and with reference to the decomposition (3) we can formulate the theorem as follows:

\textbf{5. If $m$ is a power of $2$ and $p$ is a prime number congruent to $-1$ modulo $4$, then the exponent $h$ can be so determined that all prime factors of $p$ in}
\[
  \omega\sqrt{p}^{\,-hm}
\]
\textbf{are contained only in such powers whose exponent is divisible by $m$.}

Finally let $p-1$ be divisible by $q$, or, if $m$ is even, by $4$. Let $m_1$ be the greatest common divisor of $m$ and $p-1$, and let
\[
  m=m_1m_2.
\]
Then by (9) $k$ must be divisible by $m_2$; and if we put
\[
  k=hm_2,
\]
then by (8) the product contained in $\omega$ is
\begin{equation}
  \left(\prod_\xi \pp_\xi^{\xi'}\right)^{hm_2}.
\tag{10}
\end{equation}
In addition, the prime factors $\pp_\xi$ occur in $\omega$ only in such powers whose exponents are divisible by $m$.

In this case $\xi$, by \S~172, III, runs through a complete system of residues not divisible by $q$ modulo $m_1$, and in (10) we may understand by $\xi'$ the least positive solution of the congruence
\[
  \xi\xi'\equiv1\pmod{m_1},
\]
because, if we change the exponents $\xi'$ in (10) modulo $m_1$, only $m$-th powers of the prime factors $\pp$ are added.

Then we can apply Kummer's theorem [\S~174, (16)] to the product (10), replacing $m$ there by $m_1$, and obtain, if we understand by $r_1=r^{m_2}$ an $m_1$-th root of unity and by the $\eta$ certain periods of the $p$-th roots of unity,
\[
  \prod_\xi \pp_\xi^{\xi'}=(r_1,\eta)^{m_1},
\]
and consequently
\begin{equation}
  \left(\prod_\xi \pp_\xi^{\xi'}\right)^{hm_2}=(r_1,\eta)^{hm}.
\tag{11}
\end{equation}
Consequently we have the theorem:

\textbf{6. If $p$ is a prime number and $p-1$ is divisible by $q$, or, when $q=2$, by $4$, then the positive exponent $h$ can be so determined that the number contained in $\Omega_m$}
\[
  \omega(r_1,\eta)^{-hm}
\]
\textbf{contains the prime factors of $p$ only in such powers whose exponent is divisible by $m$.}

It is now also no longer necessary to distinguish Theorem 5 from Theorem 6; for for $m_1=2$, 6 passes into 5, because then $r_1=-1$, and by Vol.~I, \S~171,
\[
  (r_1,\eta)^m=(-1,\eta)^m=\sqrt{p^{\,m}}
\]
holds.\footnote{In the author's paper in vol.~8 of \textit{Acta Mathematica} it was omitted to emphasize separately the case of Theorem 5.}

If we collect the result of Theorems 3 to 6 in one formula, we obtain, when $\eps$ denotes a unit functional, $\varphi$ any functional of the field $\Omega_m$, and $p$ runs through a sequence of prime numbers congruent to $1$ modulo $q$, among which the same prime number $p$ may occur several times,
\begin{equation}
  \omega=(r,x_0)^m=\eps\,\varphi^m\prod_p (r^{m_2},\eta)^m.
\tag{12}
\end{equation}

The resolvents of the cyclotomy occurring in these formulae, $(r^{m_2},\eta)$, are cyclotomic numbers in a higher field. But their $m$-th powers are numbers contained in the field $\Omega_m$ which are transformed by the substitution $(r,r^n)$ into
\[
  (r^{n m_2},\eta)^m.
\]
Likewise the combinations
\begin{equation}
  (r^{n m_2},\eta)^{-1}(r^{m_2},\eta)^n
\tag{13}
\end{equation}
are contained in the field $\Omega_m$ [\S~174, (9), (11)]. These resolvents are special cases of the general resolvent $(r,x_0)$. From (12), if we denote by $\varphi_n,\eps_n$ the functionals conjugate to $\varphi$ and $\eps$, we obtain
\begin{equation}
  \omega_n=\eps_n\varphi_n^m\prod_p (r^{m_2n},\eta)^m.
\tag{14}
\end{equation}
This formula teaches us the most important properties of the functionals $\varphi_n$ and $\eps_n$, first:

\textbf{7. The functionals $\varphi_n^m$ are associated with numbers of the field $\Omega_m$, and therefore belong to the principal class (\S~152).}

By 1, $\omega_n^{-1}\omega_1^n=\alpha^m$ is the $m$-th power of a number of $\Omega_m$. Since the combinations (13) are also numbers in $\Omega_m$, we can put
\[
  \prod_p (r^{n m_2},\eta)^{-1}(r^{m_2},\eta)^n=\frac{\alpha}{\beta},
\]
and obtain from (14)
\[
  \eps_n^{-1}\eps_1^n(\varphi_n^{-1}\varphi_1^n)^m=\beta^m,
\]
where $\beta$ is a number in $\Omega_m$.

It follows from this that the product
\[
  \beta\varphi_n\varphi_1^{-n}=\eps
\]
is a unit functional of the field $\Omega_m$, and from this arise, for the functionals $\varphi_n$ and the unit functionals $\eps_n$, the following two theorems:

\textbf{8. The conjugate functionals $\varphi_n$ have the property that all products}
\[
  \varphi_n\varphi_1^{-n}
\]
\textbf{belong to the principal class.}

\textbf{9. The unit functionals $\eps_n$ have the property that the products}
\[
  \eps_n\eps_1^{-n}=\eps^m
\]
\textbf{are $m$-th powers of unit functionals.}

For the proof of the principal theorem I, everything now depends on the proof of the two propositions which are to be inferred from the properties of the functionals $\varphi$ and $\eps$:

\textbf{A. The functionals $\varphi_n$ belong to the principal class.}

If we may presuppose this lemma, then $\varphi_n$ is associated with a number $\alpha_n$, and we can write formula (14), with a somewhat changed meaning of $\eps_n$, in the form
\begin{equation}
  \omega_n=\eps_n\alpha_n^m\prod_p(r^{m_2n},\eta)^m.
\tag{15}
\end{equation}
In this formula, however, $\eps_n$ is a numerical unit of the field $\Omega_m$ satisfying theorem 9.

If, therefore, from theorem 9 we can still prove the second lemma:

\textbf{B. A numerical unit $\eps_n$ of the field $\Omega_m$, which satisfies theorem 9, is the $m$-th power of a unit in $\Omega_m$, multiplied by a power of $r$;}

then in (15) we can extract the $m$-th root, and obtain, when by $\rho$ a root of unity of possibly higher degree than $m$ and by $\alpha$ a number of the field $\Omega_m$ is denoted,
\begin{equation}
  (r,x_0)=\rho\alpha\prod_p(r^{m_2},\eta),
\tag{16}
\end{equation}
where now on the right there stand nothing but cyclotomic numbers, and by this Theorem I is therefore proved.

\clearpage
\section*{\S\,182. Proof of the First Auxiliary Lemma for an Odd $m$.}

In the preceding paragraph we made the proof of Theorem I dependent on two auxiliary theorems, whose proof is now still to be sought.

We formulate the first of these auxiliary theorems as follows:

\begin{enumerate}[label=\alph*),leftmargin=2.2em]
\item If $\varphi_n$ is a system of non-zero conjugate functionals of the field $\Omega_m$, of which it is known that
\[
\varphi_n^m\quad\text{and}\quad\varphi_n^{-1}\varphi_1^n
\]
belong to the principal class, then the functionals $\varphi$ themselves belong to the principal class.
\end{enumerate}

The theorem will be proved if, for some power $\varphi^k$ of $\varphi$ whose exponent is not divisible by $q$, it can be proved that it belongs to the principal class. For if $\alpha,\beta$ are numbers in $\Omega_m$ and
\[
\varphi^m=\varepsilon'\alpha,
\qquad
\varphi^k=\varepsilon''\beta,
\]
then we can determine the rational integers $x,y$ so that
\[
mx+ky=1,
\]
and then
\[
\varphi=\varphi^{mx}\varphi^{ky}=\varepsilon^{\prime\prime\prime}\alpha^x\beta^y;
\]
thus, if $\varepsilon',\varepsilon'',\varepsilon^{\prime\prime\prime}$ are units, $\varphi$ itself is associated with a number.

According to the present state of the matter, the proof for an even $m$ is to be carried out by a quite different path from that for an odd $m$; we therefore occupy ourselves first with the easier case of an odd $m$\footnote{Compare, moreover, the new proof by Hilbert, in which the difference between even and odd $m$ recedes much more (p. 642).}.

The assumptions made in a) about $\varphi_n$ can, if $\alpha,\beta$ again denote arbitrary numbers of the field $\Omega_m$ and $\varepsilon$ unit functionals, be expressed by the two formulas
\begin{align}
\varphi_n^m&=\varepsilon\beta, \tag{1}\\
\varphi_n&=\varepsilon\alpha\varphi_1^n. \tag{2}
\end{align}
We shall now solve the problem step by step.

First let, as before, $\sigma$ be the prime factor contained in $q$, which is a number existing in $\Omega_m$, and let $\sigma^k$ be the power of $\sigma$ contained in $\varphi$; put
\begin{equation}
\varphi=\sigma^k\psi,
\tag{3}
\end{equation}
where $\psi$ denotes a functional prime to $q$, which satisfies the same conditions (1), (2) as $\varphi$, and if one of the two belongs to the principal class, the same is true of the other. Thus theorem a) needs to be proved only under the assumption that $\varphi$ is prime to $q$.

Now let $p$ be a prime number different from $q$, whose decomposition into prime factors has the expression
\begin{equation}
p=\prod_\xi \frp_\xi,
\tag{4}
\end{equation}
where $\xi$ runs through the number sequence more precisely determined in \S\,172. Let $\frp_\xi^{k_\xi}$ be the power of $\frp_\xi$ contained in $\varphi$, so that, if we set
\[
\varphi=\chi\prod_\xi \frp_\xi^{k_\xi},
\]
then $\chi$ is prime to $p$ and is related to no other prime numbers (see the end of \S\,173) than $\varphi$ itself. From this there results
\begin{equation}
\varphi_n=\chi_n\prod_\xi \frp_{n\xi}^{k_\xi},
\tag{5}
\end{equation}
and in this we let $n$ again run through the sequence of the numbers $\xi$. If we further set
\[
\psi=\prod_\xi \chi_\xi,
\]
then $\psi$ too is prime to $p$ and related to no other prime numbers than $\varphi$, and from (5), using (4), we obtain
\begin{equation}
\prod_\xi\varphi_\xi=\psi p^{\sum k_\xi},
\tag{6}
\end{equation}
and from this it follows that the functional $\psi$ satisfies the same two conditions (1), (2) as $\varphi$.

It is further shown that, if the $\varphi_\xi$ lie in the principal class, then $\psi$ is also contained in the principal class.

But if $p-1$ is not divisible by $q$, then we can also conclude the converse. For by (2), $\varphi_\xi$ is equivalent to $\varphi^\xi$, and consequently
\[
\prod_\xi\varphi_\xi\ \text{is equivalent to}\ \varphi^{\sum\xi},
\qquad
\text{is equivalent to}\ \psi,
\]
and if $\psi$ is associated with a number, the same holds for $\varphi^{\sum\xi}$. By the theorem \S\,172, however, if $p-1$ is not divisible by $q$, $\sum\xi$ is also not divisible by $q$, and consequently $\varphi$ itself is contained in the principal class.

This consideration can now be repeated if $\psi$ is still related to further prime numbers which are not congruent to 1 modulo $q$, and we arrive at the theorem:
\begin{center}
Theorem a) is proved in general if it can be proved under the assumption that $\varphi$ is related only to such prime numbers as are congruent to 1 modulo $q$.
\end{center}

If we make this assumption about the functionals $\varphi_n$, then we can apply Kummer's theorem in the form of \S\,174, 3.

In what follows we denote by $\varepsilon$ unit functionals, whose more precise knowledge is immaterial, and set
\begin{equation}
\Phi_n=\prod_t \varphi_{nt'}^{t}=\varepsilon\vartheta_n^m,
\tag{7}
\end{equation}
where $t$ runs through the residues of $m$ not divisible by $q$ which are smaller than $m$, and $t'$ is determined by $tt'\equiv1\pmod m$. Then $\vartheta_n$ is a cyclotomic number in a higher field, whose $m$th power is contained in $\Omega_m$, and which satisfies the second condition, namely that
\begin{equation}
\vartheta_n^{-1}\vartheta_1^n=\alpha
\tag{8}
\end{equation}
is a number of the field $\Omega_m$.

We now introduce a frequently used sign, understanding by $E(x)$ the greatest integer contained in the real number $x$, and by $(x)$ the remainder, which is always a proper fraction and, if $x$ itself should be an integer, is to be put equal to zero. Then
\begin{equation}
x=E(x)+(x).
\tag{9}
\end{equation}
If $x$ is an integer, then in this notation
\[
m\left(\frac xm\right)
\]
is also an integer, namely the least positive remainder of the division of $x$ by $m$.

In formula (7), the index of the functionals $\varphi$ is determined only modulo $m$. The exponent, however, is reduced to its least remainder modulo $m$. Thus if in (7) we replace the index $t'$ by $n't'$, then $t$ must be replaced by the remainder of the division of $nt$ by $m$, that is, [by (9)] by
\[
m\left(\frac{nt}{m}\right)=nt-mE\left(\frac{nt}{m}\right),
\]
and we can set
\[
\Phi_n=\prod_t \varphi_{t'}^{nt-mE(\frac{nt}{m})}=\varepsilon\vartheta_n^m,
\]
whereas
\[
\Phi_1=\prod_t \varphi_{t'}^t=\varepsilon\vartheta_1^m.
\]
Thus by (8)
\begin{equation}
\Phi_1^n\Phi_n^{-1}=\prod_t \varphi_{t'}^{mE(\frac{nt}{m})}=\varepsilon\alpha^m.
\tag{10}
\end{equation}
Accordingly the quotient
\[
\left(\frac{\prod_t \varphi_{t'}^{E(\frac{nt}{m})}}{\alpha}\right)^m
\]
is a unit functional, and since it is at the same time the $m$th power of a functional in $\Omega_m$, the base of this power is also a unit, that is,
\[
\prod_t \varphi_{t'}^{E(\frac{nt}{m})}=\varepsilon\alpha,
\]
and this product belongs to the principal class. But by our assumption
\[
\varphi_{t'}\ \text{is equivalent to}\ \varphi^{t'},
\]
and hence
\[
\prod_t \varphi_{t'}^{E(\frac{nt}{m})}\ \text{is equivalent to}\ \varphi^{\sum t'E(\frac{nt}{m})},
\]
and this holds for every arbitrary $n$ not divisible by $q$.

If therefore we can prove that $n$ can be so chosen that also the sum
\[
S_n=\sum_t t'E\left(\frac{nt}{m}\right)
\]
is not divisible by $q$, then theorem a) is proved.

We first want to make a transformation of the sum $S_n$, by which the question is led back to the much simpler case in which $m$ is a prime number.

If $m$ is a higher than first power of $q$, we set
\[
m=qm',\qquad t=t_1+qt_2,
\qquad
\begin{array}{l}
t_1=1,2,\ldots,q-1,\\
t_2=0,1,\ldots,m'-1.
\end{array}
\]
Here $m'$ is still a power of $q$; the summation letter $t_1$ runs through the number sequence $1,2,\ldots,q-1$, and $t_2$ through the number sequence $0,1,2,\ldots,m'-1$. We determine $t_1'$ from the congruence
\[
t_1t_1'\equiv1\pmod q
\]
and obtain
\[
t'\equiv t_1'\pmod q.
\]
Then
\begin{equation}
S_n=\sum_t t'E\left(\frac{tn}{m}\right)
\equiv
\sum_{t_1}t_1'\sum_{t_2}E\left(\frac{nt_2}{m'}+\frac{nt_1}{m}\right)
\pmod q.
\tag{11}
\end{equation}
We now assume $n<q$, hence also $nt_1<m$, so that $nt_1:m$ is a proper fraction. Then
\[
\begin{array}{ll}
\text{a)}& E\left(\dfrac{nt_2}{m'}+\dfrac{nt_1}{m}\right)=E\left(\dfrac{nt_2}{m'}\right),\\[1.2em]
\text{b)}& \text{or}=E\left(\dfrac{nt_2}{m'}\right)+1,
\end{array}
\]
and case b) occurs only if between
\[
\frac{nt_2}{m'}\quad\text{and}\quad \frac{nt_2}{m'}+\frac{nt_1}{m}
\]
there lies an integer. This is the case only if the difference between $nt_2:m'$ and the next greater integer is smaller than $nt_1:m$, that is, if
\begin{equation}
E\left(\frac{nt_2}{m'}\right)+1-\frac{nt_2}{m'}<\frac{nt_1}{m}.
\tag{12}
\end{equation}
If in the expression on the left side of (12)
\begin{equation}
E\left(\frac{nt_2}{m'}\right)+1-\frac{nt_2}{m'}
\tag{13}
\end{equation}
we let $t_2$ run through all its values, we obtain positive rational fractions not exceeding unity, with denominator $m'$, no two of which are equal to one another. For if two of the values arising from (13) by substitution of $t_2',t_2''$ for $t_2$ were equal, then $n(t_2'-t_2'')/m'$ would have to be an integer, which, since $n$ is relatively prime to $m'$ and $t_2'-t_2''$ is smaller than $m'$, cannot be. Thus the expressions (13) run in some order through the numerical values
\begin{equation}
\frac1{m'},\ \frac2{m'},\ldots,\frac{m'-1}{m'},\ 1,
\tag{14}
\end{equation}
and if $a$ denotes an arbitrary positive fraction, then $E(m'a)$ is the number of numbers of the sequence (14) which are not greater than $a$. By (12), therefore, the number of values of $t_2$ for which case b) occurs is equal to $E(nt_1/q)$, and there results
\[
\sum_{t_2}E\left(\frac{nt_2}{m'}+\frac{nt_1}{m}\right)
=
\sum_{t_2}E\left(\frac{nt_2}{m'}\right)+E\left(\frac{nt_1}{q}\right).
\]
To form $S_n$, we multiply this equation by $t_1'$ and sum. This gives
\[
S_n\equiv \sum_{t_1}t_1'\sum_{t_2}E\left(\frac{nt_2}{m'}\right)+\sum_{t_1}t_1'E\left(\frac{nt_1}{q}\right)\pmod q.
\]
Here now $\sum t_1'=\sum t_1\equiv0\pmod q$, because the $t_1$ pairwise give the sum $q$, and it follows that
\begin{equation}
S_n\equiv\sum_{t_1}t_1'E\left(\frac{nt_1}{q}\right)\pmod q.
\tag{15}
\end{equation}
The sum standing here on the right side is formed in the same way as $S_n$ itself in the case where $m$ is a prime number. For this case, however, there results
\[
\frac{nt_1}{q}=E\left(\frac{nt_1}{q}\right)+\left(\frac{nt_1}{q}\right),
\]
\[
\frac{(q-n)t_1}{q}=E\left(\frac{(q-n)t_1}{q}\right)+\left(\frac{(q-n)t_1}{q}\right),
\]
and from this, by addition,
\[
t_1=E\left(\frac{nt_1}{q}\right)+E\left(\frac{(q-n)t_1}{q}\right)+\left(\frac{nt_1}{q}\right)+\left(\frac{(q-n)t_1}{q}\right),
\]
and here the sum of the two positive proper fractions, which nevertheless must be an integer, can have only the value 1. It follows therefore that
\[
E\left(\frac{nt_1}{q}\right)+E\left(\frac{(q-n)t_1}{q}\right)=t_1-1,
\]
and from this, by multiplication with $t_1'$ and summation over all values of $t_1$ from 1 to $q-1$,
\[
\sum_{t_1}t_1'E\left(\frac{nt_1}{q}\right)+\sum_{t_1}t_1'E\left(\frac{(q-n)t_1}{q}\right)\equiv -1\pmod q.
\]
Consequently both sums on the left side certainly cannot be divisible by $q$. Then expression (15) shows that of the two sums $S_n$ and $S_{q-n}$ one certainly is indivisible by $q$.

This, however, was to be proved.

\section*{\S\,183. Proof of the Second Auxiliary Lemma for an Odd $m$.}

The second lemma, which by \S\,181 remains to be proved, is the following:

\begin{enumerate}[label=\arabic*.,leftmargin=2.2em,start=2]
\item If $\varepsilon_n$ is a system of conjugate numerical units of the field $\Omega_m$ which satisfies the condition that
\begin{equation}
\varepsilon_n\varepsilon_1^{-n}=\mathscr E^m
\tag{1}
\end{equation}
is the $m$th power of a unit in $\Omega_m$, then the $\varepsilon_n$ themselves are $m$th powers of units in $\Omega_m$, multiplied by a root of unity of the form $r^{nk}$.
\end{enumerate}

This proof too is quite elementary for the case of an odd $m$. We therefore first consider this case here.

By \S\,178, 1., we can represent every unit of the field $\Omega_m$ in the form
\begin{equation}
\varepsilon=\pm r^k\mathscr E(r),\qquad
\varepsilon_n=\pm r^{nk}\mathscr E(r^n),
\tag{2}
\end{equation}
where $\mathscr E(r)$ is a real unit of the field $\Omega_m$. Thus $\mathscr E(r)$ satisfies the condition
\begin{equation}
\mathscr E(r)=\mathscr E(r^{-1}),
\tag{3}
\end{equation}
from which, by (2), follows:
\begin{equation}
\varepsilon_1\varepsilon_{-1}=\mathscr E(r)^2;
\tag{4}
\end{equation}
but from (1), putting $n=-1$, one finds
\[
\varepsilon_1\varepsilon_{-1}=\mathscr E^m,
\]
that is, the left side of (4) is the $m$th power of a unit in $\Omega_m$, and hence also
\begin{equation}
\mathscr E(r)^2=\mathscr E^m
\tag{5}
\end{equation}
is the $m$th power of such a unit.

Since now $m$ is odd, the rational integer $x$ can be so determined that
\begin{equation}
2x+m=1,
\tag{6}
\end{equation}
and from this follows
\[
\mathscr E(r)=\mathscr E(r)^{2x}\mathscr E(r)^m.
\]
Thus if, according to (5), we set
\[
\mathscr E^x\mathscr E(r)=e(r),
\]
then follows
\begin{equation}
\mathscr E(r)=e(r)^m.
\tag{7}
\end{equation}
By (7), however, Theorem 2 is proved, and at the same time, for an odd $m$, the whole Theorem I.

This conclusion too fails for an even $m$, because then equation (6) is no longer soluble.

\section*{\S\,184. Preliminary Remarks on the Case of an Even $m$.}

For the case where $m$ is a power of 2, we take an entirely different path, about which some preliminary remarks may first find place here.

We have already seen in \S\,182 that the first lemma is proved if we can prove that some power of $\varphi$, whose exponent is not divisible by $q$, belongs to the principal class. In doing this, of the two properties of the functional $\varphi$ only the first is used, namely that the $m$th power of $\varphi$ is contained in the principal class.

Now by the general theorems in \S\,154 we know in every field an exponent $h$, namely the number of ideal classes, for which $\varphi^h$ belongs to the principal class whenever $\varphi$ is any functional in this field.

Thus if we can prove that the class number $h$ in the field $\Omega_m$, when $m$ is a power of 2, is an odd number, then by this alone, without anything further, Lemma 1 is proved.

This is to be the aim of our considerations in the next sections.

With respect to the second lemma the matter is similar. Here, from the assumptions in \S\,183, 2., one can derive, exactly as above, formula \S\,183, (5), from which, however, it can only be concluded that $\mathscr E(r)$ is the $\frac12m$th power of a real unit $e(r)$.

Thus, assuming the first lemma as proved, formula \S\,181, (14) gives:
\begin{equation}
(r^n,x_0)=\rho_n\sqrt{e(r^n)}\,\alpha_n\prod_p (r^{m_2n},\eta),
\tag{1}
\end{equation}
where $\rho_n$ is a root of unity of degree $m^2$, and $\alpha_n$ is a number in $\Omega_m$, and it would still have to be proved that $e(r^n)$ is the square of a unit in $\Omega_m$.

Since $e(r)$ is a real unit, we have
\begin{equation}
e(r^n)=e(r^{-n}),
\tag{2}
\end{equation}
and we shall further stipulate that the sign of the root is to be so determined (which, by a suitable assumption about $\rho_n$, we can always assume) that
\begin{equation}
\sqrt{e(r^n)}=\sqrt{e(r^{-n})}.
\tag{3}
\end{equation}
By the assumption, the product $(r^n,x_0)(r^{-n},x_0)$ is a number of the field $\Omega_m$ which is not changed by the substitution $(r,r^{-1})$, that is, a real number. Further, by \S\,174, (5),
\[
(r^{m_2n},\eta)(r^{-m_2n},\eta)=\pm p,
\]
so it is likewise real. Similarly $\alpha_n\alpha_{-n}$ is real, and from (1) it follows that
\[
\rho_n\rho_{-n}=\frac{(r^n,x_0)(r^{-n},x_0)}{e(r^n)\alpha_n\alpha_{-n}\prod_p (r^{m_2n},\eta)(r^{-m_2n},\eta)}
\]
is a real number, which, since it is a root of unity, can only be $=\pm1$.

Now in the equation resulting from this for $n=1$,
\[
(r,x_0)(r^{-1},x_0)=\pm e(r)\alpha_1\alpha_{-1}\prod_p(\pm p),
\]
we can carry out every substitution $(r,r^n)$, from which it follows that the sign of $\rho_n\rho_{-n}$ is independent of $n$, hence that
\begin{equation}
\rho_n\rho_{-n}=\rho_1\rho_{-1}=\pm1.
\tag{4}
\end{equation}
Now we further derive from (1):
\begin{equation}
\rho_n\rho_1^{-1}\sqrt{e(r^n)}\sqrt{e(r)}^{-n}=\frac{(r^n,x_0)(r,x_0)^{-n}}{\alpha_n\alpha_1^{-1}\prod_p(r^{m_2n},\eta)(r^{m_2},\eta)^{-n}},
\tag{5}
\end{equation}
and the right side is a number of the field $\Omega_m$. But the left side shows that it is a unit, and therefore, if $\mathscr E(r)$ denotes a real unit, we can set
\begin{equation}
\rho_n\rho_1^{-1}\sqrt{e(r^n)}\sqrt{e(r)}^{-n}=r^k\mathscr E(r).
\tag{6}
\end{equation}
If we substitute this value into formula (5) and make the substitution $(r,r^{-1})$, there results
\begin{equation}
\rho_{-n}\rho_{-1}^{-1}\sqrt{e(r^n)}\sqrt{e(r)}^{-n}=r^{-k}\mathscr E(r),
\tag{7}
\end{equation}
and by multiplication of (6) and (7), with regard to (4),
\begin{equation}
e(r^n)e(r)^{-n}=\mathscr E(r)^2.
\tag{8}
\end{equation}

In our considerations on the class number the following theorem will still result:
\begin{center}
A unit $e(r)$ of the real field $H_m$ which is positive with all its conjugates is the square of a unit in the field $H_m$.
\end{center}
Formula (8), however, shows that the unit $\pm e(r)$ has this property, and hence that $e(r)$ (since also $-1=i^2$ is the square of a unit) is the square of a unit of the field $\Omega_m$.

Thus the question has been reduced to the proof of the two mentioned theorems, which arises as the capstone of a long chain of considerations, but one extremely rich in interesting relations, to which the next sections are to be devoted.

\section*{Twenty-third Section. Class Number.}

\section*{\S~185. Dirichlet's theorem on the units.}

The investigations on the number of ideal classes in an algebraic field, which are now to give us the necessary supplements to the proofs of the preceding section, make use of the methods that Dirichlet introduced into number theory.\footnote{Dirichlet, \textit{Recherches sur diverses applications de l'analyse infinitésimale à la théorie des nombres}. Crelle's Journal, vols.~19, 21. Dirichlet's Works, vol.~I (1839, 1840). (The posthumous papers of Gauss, ``De nexu inter multitudinem classium etc.'' Gauss' Works, vol.~II, also belong here.) The application of these principles to cyclotomic numbers was first made by Kummer [Crelle's Journal, vols.~35, 40 (1847, 1850); Liouville's Journal, vol.~16 (1851)]. The general theory of units was likewise founded by Dirichlet [Dirichlet's Works, vol.~I, pp.~622, 633, 639 (1840, 1842, 1846)]. The general theory of units and the determination of the class numbers was given by Dedekind: Dirichlet-Dedekind, \textit{Vorlesungen über Zahlentheorie}, 4th ed., \S~183 ff. Minkowski, \textit{Geometrie der Zahlen}.}

They are no longer purely algebraic, insofar as transcendental functions and numbers, such as the natural logarithm and the number $\pi$, also occur in them. Use is also made of passages to the limit, as in integral calculus, from which we have already drawn a fine result in \S~156.

The investigations with which we first have to deal are quite general, that is, they hold for every algebraic number field, and it offers no advantage for simplicity to restrict them to special fields, such as the cyclotomic fields.

Thus let $\Omega$ be a field of the $n$th degree, and let
\begin{equation}
\Omega_1,\Omega_2,\ldots,\Omega_n
\tag{1}
\end{equation}
be the conjugate fields. The irreducible equation of the $n$th degree whose roots determine these conjugate fields for us will in general have both real and imaginary roots; the latter, however, always occur pairwise, so that to every imaginary field $\Omega^{(i)}$ there belongs another whose numbers are conjugate imaginary to those of $\Omega^{(i)}$, and which therefore arises from $\Omega^{(i)}$ by the substitution $(i,-i)$.

We combine such imaginary pairs into one unit and denote by $\nu$ the number of real fields and of pairs of imaginary fields in the sequence (1). If all conjugate fields are real, then $n=\nu$; if they are all imaginary, then $n=2\nu$.

Except in the case of the rational field, $\nu$ is equal to $1$ only in the case of an imaginary quadratic field. Otherwise $\nu$ is always greater than $1$.

The number of imaginary pairs is $n-\nu$, and the number of real fields is $2\nu-n$.

In the determinant by which, in \S~145, the square root of the ground number, $\sqrt{\Delta}$, is defined, the substitution $(i,-i)$ interchanges two conjugate imaginary rows at a time, and $\sqrt{\Delta}$ therefore changes its sign precisely when this number is odd, and otherwise not. In the first case $\sqrt{\Delta}$ is imaginary, in the second real. Hence it follows that the ground number $\Delta$ of the field $\Omega$ is positive or negative according as $n-\nu$ is even or odd.

If from two conjugate imaginary fields we retain only one, there results the sequence of conjugate fields
\[
\Omega_1,\Omega_2,\ldots,\Omega_\nu,
\]
to which we assign a system of signs
\[
\delta_1,\delta_2,\ldots,\delta_\nu
\]
with the convention that $\delta_s=1$ if $\Omega_s$ is real, and $=2$ if $\Omega_s$ is imaginary, so that $\sum \delta_s=n$.

Let $\eta$ denote a non-zero number of the field $\Omega$ and let
\[
\eta_1,\eta_2,\ldots,\eta_n
\]
be the conjugate numbers. To this number system we assign another number system
\[
\lambda_1,\lambda_2,\ldots,\lambda_\nu,
\]
which we define by the equation
\begin{equation}
\lambda_s=\delta_s\log |\eta_s|,
\tag{2}
\end{equation}
where $|\eta_s|$ denotes the absolute value of $\eta_s$, and $\log|\eta_s|$ denotes the real natural logarithm of the positive number $|\eta_s|$. If $\Omega_s$ is real and the sign $\pm$ is chosen so that $\pm\eta_s$ is positive, then
\[
\lambda_s=\log(\pm\eta_s).
\]
But if $\eta_s,\eta_s'$ is an imaginary pair, then
\[
\lambda_s=\log\eta_s\eta_s',
\]
and to $\eta_s$ and $\eta_s'$ belongs the same $\lambda_s$, so that the number of different $\lambda_s$ is equal to $\nu$. From this determination one further obtains
\begin{equation}
\lambda_1+\lambda_2+\cdots+\lambda_\nu=\log N_a(\eta),
\tag{3}
\end{equation}
where, as before, $N_a$ denotes the absolute norm.

We shall call the numbers $\lambda_1,\lambda_2,\ldots,\lambda_\nu$ the conjugate logarithms of the number $\eta$.

If $\eta$ is an integer of the field $\Omega$, then its absolute norm is a natural number, which is equal to $1$ only when $\eta$ is a unit; and from (3) it follows that:

\textbf{1. The sum of the conjugate logarithms of an integer $\eta$ is positive, and is equal to zero only when $\eta$ is a unit.}

Let $\omega_1,\omega_2,\ldots,\omega_n$ denote a basis of the integers in $\Omega$ (a minimal basis of $\Omega$, \S~145), and let $\omega_{1,s},\omega_{2,s},\ldots,\omega_{n,s}$, for $s=1,2,\ldots,n$, denote the conjugate numbers. Then all integers of the field $\Omega_s$ can be represented in the form
\begin{equation}
\eta_s=x_1\omega_{1,s}+x_2\omega_{2,s}+\cdots+x_n\omega_{n,s},
\tag{4}
\end{equation}
with rational integral coefficients $x_1,x_2,\ldots,x_n$. The determinant of this system,
\[
\sum \pm \omega_{1,1}\omega_{2,2}\cdots\omega_{n,n}=\sqrt{\Delta},
\]
is the square root of the discriminant $\Delta$ of the field, and hence is always different from zero. Thus the system (4) can be solved for the unknowns $x_1,x_2,\ldots,x_n$, and if we now require that the absolute values of the numbers $\eta_s$ are not to exceed a certain bound, then these solutions give bounds for the rational integers $x_i$.

This simple remark gives us the following important theorem:

\textbf{2. In an algebraic field $\Omega$ there is only a finite number of integers $\eta$ of the property that the absolute values of the conjugate numbers $\eta_s$ lie below a given bound.}

We now make use of the general theorem \S~155, (25), according to which, if $f(x_1,x_2,\ldots,x_n)$ is any positive quadratic form in $n$ variables with determinant $D$, the variables $x_i$ can be determined as rational integers, not all zero, so that
\[
f(x_1,x_2,\ldots,x_n)< n\sqrt[n]{0.404D},
\]
that is, becomes smaller than a number depending only on the determinant $D$.

We apply this theorem, as in \S~156, to the form
\[
f(x_1,x_2,\ldots,x_n)=\frac{|\eta_1|^2}{c_1^2}+\frac{|\eta_2|^2}{c_2^2}+\cdots+\frac{|\eta_n|^2}{c_n^2},
\]
where the $\eta_s$ are the linear forms (4), and in the case of a real $\eta_s$
\[
|\eta_s|^2=(\eta_s)^2,
\]
in the case of an imaginary pair $\eta_s,\eta_s'$
\[
|\eta_s|^2=|\eta_s'|^2=\eta_s\eta_s'.
\]
In the latter case $\eta_s\eta_s'$ is the sum of two real squares, and hence, if the $c_s$ are taken to be real, $f$ is a positive form. We further stipulate, for the hitherto arbitrary real quantities $c_s$, that if $\eta_s,\eta_s'$ form a conjugate imaginary pair, then
\begin{equation}
c_s=c_s'
\tag{5}
\end{equation}
is to hold; besides this we subject the $c_s$ to the condition
\begin{equation}
c_1c_2\cdots c_n=1.
\tag{6}
\end{equation}
We form the determinant of the function $f$, as in \S~156, by first regarding it as a form in the variables $\eta_1,\eta_2,\ldots,\eta_n$, and the determinant of this form is, because of (6), equal to $\pm1$. The determinant of $f$ with respect to the variables $x_i$ is therefore (Vol.~I, \S~56), apart from the sign, equal to the
square of the determinant of the linear forms $\eta_s$, that is, to the ground number of the field $\Omega$.

It follows that one can find a number $A$, independent of the $c_s$ and determined only by the field $\Omega$, of such a kind that, whatever the $c_s$ may be, $f(x_1,x_2,\ldots,x_n)$ for integral non-zero $x$ drops below $A$. Since $f$ is a sum of positive summands, the same holds for each of these summands, and the following theorem is thereby proved:

\textbf{3. If $c_s$ is any system of real positive numbers satisfying the conditions (5), (6), then there is a real number $A$, determined by the field $\Omega$ alone, such that one can always determine an integer $\eta$ of the field $\Omega$ which, together with its conjugate numbers $\eta_s$, satisfies the conditions}
\begin{equation}
|\eta_s|<c_sA.
\tag{7}
\end{equation}

Incidentally it follows from this that in every algebraic field $\Omega$, except the rational and the imaginary quadratic field, one can find non-zero integers whose absolute value falls below any given bound.

We shall transform the expression of theorem 3 somewhat by introducing the conjugate logarithms of $\eta$. Put $\log A=k$, so that $k$ is a real number determined by $\Omega$, and further
\begin{equation}
c_s=e^{\frac{\gamma_su}{\delta_s}},
\tag{8}
\end{equation}
where $u$ may be any real quantity, and where the numbers $\gamma_s$ satisfy, because of (5) and (6), the condition
\begin{equation}
\gamma_1+\gamma_2+\cdots+\gamma_\nu=\sum_{s=1}^{\nu}\gamma_s=0.
\tag{9}
\end{equation}
Since, moreover,
\[
|\eta_s|=e^{\frac{\lambda_s}{\delta_s}},
\]
it follows from (7) and (8) that
\begin{equation}
\lambda_s-\gamma_su<k\delta_s.
\tag{10}
\end{equation}
The sum of the $\nu$ expressions on the left side of this inequality is, by (9), equal to $\sum\lambda_s$, hence, by theorem 1, is not negative; and we find
\[
0\leq \sum_s(\lambda_s-\gamma_su)<kn.
\]

It follows from this, in view of (10), that a term of this sum, $\lambda_s-\gamma_su$, cannot be smaller than $-(n-\delta_s)k$, for otherwise the total sum would be negative; and if, for simplicity, we take the lower bound somewhat smaller,
\begin{equation}
-nk\delta_s<\lambda_s-\gamma_su<k\delta_s.
\tag{11}
\end{equation}

It is thereby proved that for every given system of numbers $u,\gamma_s$, provided only that condition (9) is fulfilled, there is an integer $\eta$ of the field $\Omega$ which, with its conjugate numbers, satisfies the limiting conditions (11).

Now let $g_s$ be a system of real quantities of which we demand only that
\[
\gamma_1g_1+\gamma_2g_2+\cdots+\gamma_\nu g_\nu=\alpha
\]
be different from zero. By (9) this excludes the case where all the $g_s$ are equal to one another; but if they are not, then for every given system of the $g_s$ there will be infinitely many systems $\gamma_s$ that satisfy the requirement (9).

If now $k\sum\delta_sg_s=g$ is put, then from (11), by multiplication with $g_s$ and summation, one obtains
\begin{equation}
-ng+\alpha u<\sum g_s\lambda_s<g+\alpha u.
\tag{12}
\end{equation}

According to this limiting condition, keeping $g_s$ and $\alpha$ fixed, we now determine a sequence of integers
\begin{equation}
\eta,\eta',\eta'',\eta''',\ldots,
\tag{13}
\end{equation}
by making, for $u$, a sequence of assumptions
\[
u,u',u'',u''',\ldots,
\]
which are more precisely determined as follows.

First choose $u$ so that
\[
-ng+\alpha u=a,
\qquad
 g+\alpha u=a'
\]
are positive; then put
\[
\delta=\frac{a'-a}{\alpha}=\frac{(n+1)g}{\alpha},
\]
and put
\[
u'=u+\delta,
\quad
u''=u'+\delta,
\quad
u'''=u''+\delta,
\quad\ldots,
\]
\[
a+\alpha\delta=a',
\quad a'+\alpha\delta=a'',
\quad a''+\alpha\delta=a''',
\quad\ldots,
\]
and from (12) we thus obtain, if $\lambda_s',\lambda_s'',\ldots$ are the conjugate logarithms of $\eta',\eta'',\ldots$,
\begin{equation}
\begin{array}{rcl}
a&<&\sum g_s\lambda_s<a',\\
a'&<&\sum g_s\lambda_s'<a'',\\
a''&<&\sum g_s\lambda_s''<a''',\\
&&\cdots\cdots,
\end{array}
\tag{14}
\end{equation}
and the sequence of these inequalities can be continued without bound.

Moreover all these numbers $\eta,\eta',\eta'',\ldots$ have, by (6) and (7), the property that their absolute norms $N,N',N'',\ldots$ remain below a fixed finite bound $e^{nk}$.

The number of possible values of $N,N',N'',\ldots$ is therefore finite, certainly not greater than the greatest integer contained in $e^{nk}$. Likewise the number of all possible residues of the numbers $x_1,x_2,\ldots,x_n$ [in (4)] modulo all the moduli $N,N',N'',\ldots$ is finite; and from this it follows that, if we continue the sequence of numbers (13) far enough, two distinct numbers $\eta^{(h)},\eta^{(k)}$ must occur in the sequence for which $N^{(h)}=N^{(k)}$ and the residues of $x_1,x_2,\ldots,x_n$ modulo $N^{(h)}$ are exactly the same.

But from (14), if $h>k$ is assumed, it follows that
\begin{equation}
\sum g_s\bigl(\lambda_s^{(h)}-\lambda_s^{(k)}\bigr)>0.
\tag{15}
\end{equation}
If $N$ is the absolute norm of these two numbers $\eta^{(h)},\eta^{(k)}$, then, because of the agreement of the residues of the $x$, the difference $\eta^{(h)}-\eta^{(k)}$ is divisible by $N$. On the other hand, by \S~138 (13), the number $N$ is divisible by $\eta^{(h)}$, and consequently $\eta^{(h)}-\eta^{(k)}$ is also divisible by $\eta^{(h)}$. It follows that $\eta^{(k)}$ is also divisible by $\eta^{(h)}$, and likewise $\eta^{(h)}$ by $\eta^{(k)}$. The two numbers are therefore associated, and if we put
\[
\frac{\eta^{(h)}}{\eta^{(k)}}=\varepsilon,
\]
then $\varepsilon$ is a unit of the field $\Omega$.

Moreover, if, denoting by $|\varepsilon_s|$ the absolute value of $\varepsilon_s$, we put
\begin{equation}
\delta_s\log|\varepsilon_s|=l_s,
\tag{16}
\end{equation}
then from (15) we obtain
\begin{equation}
g_1l_1+g_2l_2+\cdots+g_\nu l_\nu>0.
\tag{17}
\end{equation}

This proves the theorem:

\textbf{4. If $g_1,g_2,\ldots,g_\nu$ is any system of real numbers which are not all equal to one another, then in every algebraic field $\Omega$ one can determine a unit $\varepsilon$ such that the sum}
\[
g_1l_1+g_2l_2+\cdots+g_\nu l_\nu
\]
\textbf{is different from zero, where $l_s$ is the system of conjugate logarithms of $\varepsilon$.}

This important theorem, which is the foundation for what follows, is due, though in a somewhat changed form, to Dirichlet. This formulation was given by Minkowski.

\section*{\S~186. Systems of independent units and exponent systems of the units.}

If in the theorem just proved we put $g_1=1$, $g_2=0,\ldots,g_\nu=0$, which is allowed if from now on we assume $\nu>1$, it follows that in $\Omega$ there is a unit $\varepsilon$ for which one of the conjugate logarithms, $l_1$, is different from zero. Since every power of $\varepsilon$ has the same property, there are also infinitely many such units.

This theorem now permits a very important generalization:

\textbf{5. If $s\leq \nu-1$, then in the field $\Omega$ one can determine a system of units with the associated conjugate logarithms}
\begin{equation}
\begin{array}{cccc}
\varepsilon_1: & l_{1,1},&l_{1,2},&\ldots,l_{1,\nu},\\
\varepsilon_2: & l_{2,1},&l_{2,2},&\ldots,l_{2,\nu},\\
&\multicolumn{3}{c}{\cdots\cdots\cdots}\,\\
\varepsilon_s: & l_{s,1},&l_{s,2},&\ldots,l_{s,\nu},
\end{array}
\tag{1}
\end{equation}
\textbf{in such a way that any one of the $s$-rowed determinants of the matrix of the $l_{h,k}$, for instance}
\[
L_s=\sum \pm l_{1,1}l_{2,2}\cdots l_{s,s},
\]
\textbf{is different from zero.}

For $s=1$ the stated theorem is correct by the remark made at the beginning. We therefore assume that it
has been proved for $s-1$ and derive it generally by complete induction. For this purpose we arrange the determinant $L_s$ according to the elements of the last row, and write it thus:
\begin{equation}
L_s=g_1l_{s,1}+g_2l_{s,2}+\cdots+g_sl_{s,s},
\tag{2}
\end{equation}
where $g_s=L_{s-1}$ and hence, by the assumption made, can be taken to be different from zero. If we substitute these values of $g$ in the formula of theorem 4, \S~185, taking $g_{s+1}=0,\ldots,g_\nu=0$, while $g_s$ is not equal to zero, then, as long as $s<\nu$, certainly not all $g_1,g_2,\ldots,g_\nu$ are equal to one another; and it follows that $\varepsilon_s$ can be so chosen that $L_s$ is different from zero, as was to be proved.

This mode of inference cannot be extended to $s=\nu$, and the determinant $L_\nu$ must indeed always be zero, since for every unit $\sum_{1}^{\nu} l_s$ vanishes.

For the case $s=\nu-1$ we shall now denote the determinant $L_s$ by
\begin{equation}
L_{\nu-1}=L(\varepsilon_1,\varepsilon_2,\ldots,\varepsilon_{\nu-1}),
\tag{3}
\end{equation}
and call a system of units $\varepsilon_1,\varepsilon_2,\ldots,\varepsilon_{\nu-1}$ for which this determinant is different from zero a system of independent units.

The absolute value $L$ of the determinant $L(\varepsilon_1,\varepsilon_2,\ldots,\varepsilon_{\nu-1})$ is called the regulator of the system $\varepsilon_1,\varepsilon_2,\ldots,\varepsilon_{\nu-1}$.\footnote{Dirichlet-Dedekind, \textit{Vorlesungen über Zahlentheorie}, 4th ed., \S~183.}

If in the determinant
\begin{equation}
\left|\begin{array}{cccc}
l_{1,1},&l_{1,2},&\ldots,&l_{1,\nu}\\
\cdots&\cdots&\cdots&\cdots\\
l_{\nu-1,1},&l_{\nu-1,2},&\ldots,&l_{\nu-1,\nu}\\
x_1,&x_2,&\ldots,&x_\nu
\end{array}\right|,
\tag{4}
\end{equation}
in which $x_1,x_2,\ldots,x_\nu$ are arbitrary quantities, we add all columns to the last one, then, with regard to the relations $\sum_{1}^{\nu}l_{s,i}=0$, we obtain its value equal to
\[
\pm(x_1+x_2+\cdots+x_\nu)L.
\]
If, therefore, we take all the $x$ except for an arbitrary one
to be $=0$, and the remaining one to be $=1$, then from (3) there results one of the $(\nu-1)$-rowed determinants of the matrix
\[
\begin{array}{ccc}
l_{1,1},&\ldots,&l_{1,\nu}\\
\cdots&\cdots&\cdots\\
l_{\nu-1,1},&\ldots,&l_{\nu-1,\nu}
\end{array}
\]
and all these determinants therefore have, apart from the sign, the same value. It is therefore immaterial which one of the $\nu$ conjugate values is omitted in forming the regulator $L$.

If $\varepsilon_1,\varepsilon_2,\ldots,\varepsilon_{\nu-1}$ is an independent system of units, and if $\lambda_1,\lambda_2,\ldots,\lambda_\nu$ are the conjugate logarithms of any unit $\varepsilon$ in $\Omega$, then the numbers $\xi_1,\xi_2,\ldots,\xi_{\nu-1}$ can always and uniquely be determined so that
\begin{equation}
\begin{array}{rcl}
\xi_1l_{1,1}+\xi_2l_{2,1}+\cdots+\xi_{\nu-1}l_{\nu-1,1}&=&\lambda_1,\\
\xi_1l_{1,2}+\xi_2l_{2,2}+\cdots+\xi_{\nu-1}l_{\nu-1,2}&=&\lambda_2,\\
\multicolumn{3}{c}{\cdots\cdots\cdots}\,\\
\xi_1l_{1,\nu}+\xi_2l_{2,\nu}+\cdots+\xi_{\nu-1}l_{\nu-1,\nu}&=&\lambda_\nu
\end{array}
\tag{5}
\end{equation}
holds. For the first $\nu-1$ of these equations form a system of linear equations with non-vanishing determinant, and the $\nu$th equation follows from the others because the sum of the conjugate logarithms of every unit vanishes.

We call the system of numbers $\xi_1,\xi_2,\ldots,\xi_{\nu-1}$ the exponent system of the unit $\varepsilon$ with respect to the system $\varepsilon_1,\varepsilon_2,\ldots,\varepsilon_{\nu-1}$, and it remains to prove:

\textbf{6. That the exponents of every unit $\varepsilon$ are rational numbers whose denominators do not exceed a certain finite bound, and which therefore can all be assumed to have a common denominator depending only on the system $\varepsilon_i$.}

To see this, one must consider the following:

1) If we regard the quantities $\xi_1,\xi_2,\ldots,\xi_{\nu-1}$ as variables and let each of them, independently of the others, run from 0 to 1, then the left sides of (5) remain enclosed within finite bounds depending only on the units $\varepsilon_1,\varepsilon_2,\ldots,\varepsilon_{\nu-1}$. Within these bounds, therefore, the conjugate logarithms $\lambda$ of the unit $\varepsilon$ must also remain, as long as
the exponents $\xi_i$ remain restricted to non-negative proper fractional values; and hence, from the definition of the conjugate logarithms, one obtains an upper bound for $\varepsilon$ and its conjugates.

But by theorem 2, \S~185, there are in $\Omega$ only finitely many integers, and therefore still more only finitely many units, which in absolute value, with all their conjugates, lie below a finite bound.

If we now call a unit whose exponents lie between the bounds 0 and 1, including the lower and excluding the upper bound, a unit reduced with respect to the system $\varepsilon_1,\varepsilon_2,\ldots,\varepsilon_{\nu-1}$, then we have the theorem:

\textbf{There are only finitely many units in $\Omega$ which are reduced with respect to an independent system $\varepsilon_1,\varepsilon_2,\ldots,\varepsilon_{\nu-1}$.}

2) The units reproduce themselves by multiplication and division. If
\[
\begin{array}{c}
\xi_1',\xi_2',\ldots,\xi_{\nu-1}',\\
\xi_1'',\xi_2'',\ldots,\xi_{\nu-1}''
\end{array}
\]
are the exponents of two units $\varepsilon',\varepsilon''$, then the sums and the differences
\[
\xi_1'\pm\xi_1'',\ \xi_2'\pm\xi_2'',\ldots,\xi_{\nu-1}'\pm\xi_{\nu-1}''
\]
are the exponents of the units $\varepsilon'\varepsilon''$ and $\varepsilon':\varepsilon''$.

This follows immediately from the theorem that the logarithm of a product or of a quotient is equal to the sum or the difference of the logarithms of the two components.

3) The units $\varepsilon_1,\varepsilon_2,\ldots$ themselves have the exponents
\[
1,0,\ldots,0;\quad 0,1,\ldots,0;\quad\ldots,
\]
and from 2) it follows that every system of integers, substituted for $\xi_1,\xi_2,\ldots,\xi_{\nu-1}$, gives the exponent system of a unit which can be formed by multiplication of powers of the $\varepsilon_1,\varepsilon_2,\ldots$.

4) From 2) and 3) it follows that an exponent system $\xi_1,\xi_2,\ldots,\xi_{\nu-1}$ of a unit remains an exponent system if all its elements are multiplied by one and the same integer, and that it also retains this character when its elements $\xi_i$ are increased or decreased by arbitrary integers. Thus, if we use the sign $(x)$ in the sense of \S~182, that it denotes the excess
of the number $x$ over the greatest integer contained in $x$, the following is obtained:

If $\xi_1,\xi_2,\ldots,\xi_{\nu-1}$ is the exponent system of a unit, and $m$ is an integer, then also
\begin{equation}
(m\xi_1),(m\xi_2),\ldots,(m\xi_{\nu-1})
\tag{6}
\end{equation}
is the exponent system of a unit.

Now the quantities $(m\xi_i)$, when they are not zero, are positive proper fractions; and by 1), if the sequence of integers is continued sufficiently far, it must therefore occur that for two distinct values $m',m''$ of $m$ the number sequence (6) agrees. Since accordingly $m'\xi_i,m''\xi_i$ have the same excess over an integer, $(m'-m'')\xi_i$ is itself an integer; and so there is a rational integer $m$, at most equal to the number of exponent systems of reduced units, with the property that
\[
m\xi_1,m\xi_2,\ldots,m\xi_{\nu-1}
\]
become rational integers.

This number $m$ may change when the unit $\varepsilon$ is changed. But since all possible values of this number lie below a finite bound, there is a finite number in which all these values of $m$ are contained and which itself can be taken for $m$. Hence there is a certain positive integer $m$ which can be taken as denominator of all rational numbers which can occur as exponents of units. This proves theorem 6.

\section*{\S~187. Fundamental systems of units.}

We now choose, instead of the system of independent units $\varepsilon_i$, another system whose conjugate logarithms may be
\begin{equation}
\lambda_{i,1},\lambda_{i,2},\ldots,\lambda_{i,\nu}
\qquad i=1,2,\ldots,\nu-1.
\tag{1}
\end{equation}
By \S~186, (5),
\begin{equation}
\begin{array}{rl}
\lambda_{i,k}={}&\xi_{1,i}l_{1,k}+\xi_{2,i}l_{2,k}+\cdots+\xi_{\nu-1,i}l_{\nu-1,k},\\[2pt]
&i=1,2,\ldots,\nu-1;\quad k=1,2,\ldots,\nu,
\end{array}
\tag{2}
\end{equation}
where $\xi_{1,i},\xi_{2,i},\ldots,\xi_{\nu-1,i}$ denote the exponents of one of the new units with respect to the system $\varepsilon_1,\varepsilon_2,\ldots,\varepsilon_{\nu-1}$, which we may assume to be rational numbers with the common denominator $m$.

But by the multiplication theorem for determinants,
\begin{equation}
\left|\begin{array}{ccc}
\lambda_{1,1},&\ldots,&\lambda_{\nu-1,1}\\
\cdots&\cdots&\cdots\\
\lambda_{1,\nu-1},&\ldots,&\lambda_{\nu-1,\nu-1}
\end{array}\right|
=
\left|\begin{array}{ccc}
\xi_{1,1},&\ldots,&\xi_{1,\nu-1}\\
\cdots&\cdots&\cdots\\
\xi_{\nu-1,1},&\ldots,&\xi_{\nu-1,\nu-1}
\end{array}\right|
\left|\begin{array}{ccc}
l_{1,1},&\ldots,&l_{\nu-1,1}\\
\cdots&\cdots&\cdots\\
l_{1,\nu-1},&\ldots,&l_{\nu-1,\nu-1}
\end{array}\right|,
\tag{3}
\end{equation}
or, if we put
\[
\Lambda_{\nu-1}=\sum\pm\lambda_{1,1}\lambda_{2,2}\cdots\lambda_{\nu-1,\nu-1}
\]
and observe that the determinant
\[
\sum\pm\xi_{1,1}\xi_{2,2}\cdots\xi_{\nu-1,\nu-1}
\]
is a rational number with denominator $m^{\nu-1}$,
\begin{equation}
\Lambda_{\nu-1}=\frac{a}{m^{\nu-1}}L(\varepsilon_1,\varepsilon_2,\ldots,\varepsilon_{\nu-1}),
\tag{4}
\end{equation}
where $a$ is a rational integer.

It follows that the newly introduced units form an independent system if and only if the determinant of the $\xi_{i,k}$, that is, the number $a$, is different from zero.

Now among a finite or infinite number of non-zero rational fractions with the same denominator there is always one of least absolute value, and consequently, by (4), we can choose the new system of independent units so that its regulator $\pm\Lambda_{\nu-1}$ becomes as small as possible. Such a system we call a fundamental system of units; and the minimum value of the regulator itself, that is, the regulator of a fundamental system, we call the regulator of the field. We now assume that our system $\varepsilon_1,\varepsilon_2,\ldots,\varepsilon_{\nu-1}$ itself is a fundamental system.

Under this assumption it can be proved that all exponents of units must be integers.

Indeed, if we assume that there exists a unit $\varepsilon$ whose exponents determined by the formulae \S~186, (5) are not all integers, then by \S~186, 4) there is also a unit $\varepsilon_0$ whose exponents are proper fractions and do not all vanish.

If by $\lambda_1,\lambda_2,\ldots,\lambda_{\nu-1}$ in the formulae (5), \S~186, we understand the system of conjugate logarithms of $\varepsilon_0$, then
\[
L(\varepsilon_0,\varepsilon_2,\ldots,\varepsilon_{\nu-1})=\xi_1L(\varepsilon_1,\varepsilon_2,\ldots,\varepsilon_{\nu-1}),
\]
and if we therefore assume, which is plainly no restriction of generality, that $\xi_1$ is a non-zero proper fraction, then this formula contradicts the assumption that
\[
\varepsilon_1,\varepsilon_2,\ldots,\varepsilon_{\nu-1}
\]
is a fundamental system, because the determinant $L_{\nu-1}$ would be made smaller if $\varepsilon_1$ were replaced by $\varepsilon_0$.

Once we have a fundamental system, we can take any system of rational integers $\xi_1,\xi_2,\ldots,\xi_{\nu-1}$ as an exponent system and form a unit $\varepsilon$ with these exponents:
\[
\varepsilon=\varepsilon_1^{\xi_1}\varepsilon_2^{\xi_2}\cdots\varepsilon_{\nu-1}^{\xi_{\nu-1}}.
\]

It remains to answer the question how far a unit is determined by its exponent system.

Suppose that $\varepsilon',\varepsilon''$ are two units with the same exponent system. Then the exponent system of the unit $\varepsilon':\varepsilon''=\rho$ consists entirely of zeros, and consequently the conjugate logarithms of $\rho$ are all equal to zero. Therefore $\rho$ is an integer with the property that the absolute values of all numbers conjugate with $\rho$ are equal to $1$, and hence, as proved in \S~175, 2, it is a root of unity. Thus the following theorem is proved:

\textbf{1. There is in the field $\Omega$ a system of $\nu-1$ fundamental units}
\begin{equation}
\varepsilon_1,\varepsilon_2,\ldots,\varepsilon_{\nu-1},
\tag{5}
\end{equation}
\textbf{having the property that in the form}
\begin{equation}
\varepsilon=\rho\varepsilon_1^{\xi_1}\varepsilon_2^{\xi_2}\cdots\varepsilon_{\nu-1}^{\xi_{\nu-1}}
\tag{6}
\end{equation}
\textbf{all units of the field, each only once, are contained, when $\xi_1,\xi_2,\ldots,\xi_{\nu-1}$ run through all rational integers and $\rho$ runs through all roots of unity present in $\Omega$.}

It is now easy to derive all other fundamental systems from one fundamental system, by putting in (6), for the exponents, $\nu-1$ different systems of integers $\xi_{i,k}$ whose determinant is $=\pm1$. For if one puts
\[
\varepsilon_i'=\rho_i\varepsilon_1^{\xi_{1,i}}\varepsilon_2^{\xi_{2,i}}\cdots\varepsilon_{\nu-1}^{\xi_{\nu-1,i}},
\]
then from (3) it follows that
\[
L(\varepsilon_1',\varepsilon_2',\ldots,\varepsilon_{\nu-1}')=\pm L(\varepsilon_1,\varepsilon_2,\ldots,\varepsilon_{\nu-1}).
\]

Thus both determinants have the minimum value, and both systems are fundamental systems.

The case $\nu=1$, which we excluded above, occurs, apart from the case of the rational field, in which only the two units $\pm1$ exist, in the case of the imaginary quadratic field. In this case, in the twentieth section, we represented the integers of the field in the form
\[
\frac{x+y\sqrt{-m}}{2},
\]
where $-m$ denotes a fundamental discriminant, so that $m\equiv0$ or $\equiv3\pmod 4$, and $x,y$ are both even or both odd. We obtain the units if we solve the equation
\[
x^2+my^2=4
\]
in all possible ways in rational integers. But this equation, if $m>4$, has only the two solutions $x=\pm2,y=0$, and thus in these cases, as in the field of rational numbers, there are only the two units $\pm1$. If $m=4$, one finds the four units $\pm1,\pm i$, and finally if $m=3$, the six units
\[
\pm1,
\quad \pm\frac{-1+i\sqrt3}{2},
\quad \pm\frac{-1-i\sqrt3}{2}.
\]

In the next simplest case, $n=2$, $\nu=2$, that is, in the real quadratic field, the theory of units coincides with the theory of Pell's equation treated in \S~127 of the first volume.

\clearpage
\section*{\S\,188. Reduced numbers.}

If $\alpha$ is any non-zero integral or fractional number of the field $\Omega$, then, as with the units, we consider the system of conjugate logarithms of $\alpha$, by which we understand, as in \S~185, the real numbers
\begin{equation}
\lambda_1=\delta_1\log|\alpha_1|,\quad
\lambda_2=\delta_2\log|\alpha_2|,\ldots,\quad
\lambda_\nu=\delta_\nu\log|\alpha_\nu|,
\tag{1}
\end{equation}
which we also denote by $\lambda_s(\alpha)$. If we bring in the fundamental system $\varepsilon_1,\varepsilon_2,\ldots,\varepsilon_{\nu-1}$ of units, then

we can form the linear equations, analogously to \S~186, (5):
\begin{equation}
\begin{aligned}
\xi_1l_{1,1}+\xi_2l_{2,1}+\cdots+\xi_{\nu-1}l_{\nu-1,1}+\delta_1\xi_\nu&=\lambda_1,\\
\xi_1l_{1,2}+\xi_2l_{2,2}+\cdots+\xi_{\nu-1}l_{\nu-1,2}+\delta_2\xi_\nu&=\lambda_2,\\
&\cdots\cdots\cdots\\
\xi_1l_{1,\nu}+\xi_2l_{2,\nu}+\cdots+\xi_{\nu-1}l_{\nu-1,\nu}+\delta_\nu\xi_\nu&=\lambda_\nu,
\end{aligned}
\tag{2}
\end{equation}
whose determinant [\S~186, (4)] has the non-zero value
\begin{equation}
\left|
\begin{array}{ccccc}
l_{1,1}&l_{2,1}&\ldots&l_{\nu-1,1}&\delta_1\\
l_{1,2}&l_{2,2}&\ldots&l_{\nu-1,2}&\delta_2\\
\cdots&\cdots&\cdots&\cdots&\cdots\\
l_{1,\nu}&l_{2,\nu}&\ldots&l_{\nu-1,\nu}&\delta_\nu
\end{array}
\right|=nL(\varepsilon_1,\varepsilon_2,\ldots,\varepsilon_{\nu-1}).
\tag{3}
\end{equation}
Accordingly the unknowns $\xi_1,\xi_2,\ldots,\xi_{\nu-1},\xi_\nu$ are determined uniquely from (2).

By adding the equations (2) one first obtains very simply
\begin{equation}
n\xi_\nu=\log \Na(\alpha),
\tag{4}
\end{equation}
and therefore $\xi_\nu$ remains the same for all numbers $\alpha$ that are associated with one another.

The numbers $\xi_1,\xi_2,\ldots,\xi_{\nu-1}$ are called the exponents of the number $\alpha$ [with respect to the fundamental system $(\varepsilon_1,\varepsilon_2,\ldots,\varepsilon_{\nu-1})$].

From this definition it follows without further ado that the exponents of a product or of a quotient are equal to the sum or the difference of the corresponding exponents of the individual components. The exponents of a unit are rational integers, and every system of rational integers $\xi_1,\xi_2,\ldots,\xi_{\nu-1}$ is also the exponent system of a unit.

The exponents of associated numbers therefore differ from one another by integers, and one can find, for each number $\alpha$, an associated number $\alpha_0$ whose exponents lie between 0 and 1. Such numbers are called \textbf{reduced numbers} (with respect to the fundamental system $\varepsilon_1,\ldots,\varepsilon_{\nu-1}$). The reduced units are all, and only, the roots of unity present in $\Omega$, whose number we shall denote by $w$. Since the two roots of unity $\pm1$ are contained in every field, $w$ is at least $=2$ and is always an even number.

If two associated numbers have the same exponent system, then they differ only by a factor which

is a root of unity. The exponent system of the reduced number $\alpha_0$ derived from $\alpha$ is completely determined by $\alpha$ itself. This proves the theorem:

\begin{center}
\textbf{II. For every non-zero number $\alpha$ of the field $\Omega$ there exist $w$, and no more, reduced numbers $\rho\alpha_0$.}
\end{center}

If another fundamental system is used in place of $(\varepsilon_1,\varepsilon_2,\ldots,\varepsilon_{\nu-1})$, then the exponents $\xi_1,\xi_2,\ldots,\xi_{\nu-1}$ undergo an integral linear substitution of determinant $\pm1$. A reduced number can then cease to be reduced for the new fundamental system. The concept of a reduced number is therefore dependent on the choice of the fundamental system. But the number of reduced numbers derived from a given number is always the same.

By this theorem we have attained the purpose of selecting from the whole infinite system of mutually associated numbers a definite finite number of representatives.

\section*{\S\,189. Bounds for the number of integral numbers of the field $\Omega$ divisible by an ideal.}

Our earlier considerations have shown that every integral number of a field $\Omega$ has only a finite number of ideals as divisors. Since every integral ideal is a divisor of its norm, there are consequently also only finitely many integral numbers in $\Omega$ whose absolute norm does not exceed a given bound.

From these numbers we now again pick out a part, and ask for the number $T$ of all \textbf{integral numbers} of the field $\Omega$ whose absolute norm does not exceed a positive magnitude $t$, and which are divisible by a given ideal $\mathfrak a$.

If we take the rational field as a preliminary example, then there $T$ is the number of natural numbers which are smaller than $t$ and divisible by a given rational integer $m$, that is, in the notation used earlier (\S~182), $T=E\left(\frac{t}{m}\right)$.

The number $T$ is, in general, not a function of $t$ that can be determined more exactly; it grows to infinity together with $t$, but since it has only integral values it changes only step by step, and hence discontinuously. What matters here is the \textbf{limit to which the ratio $T:t$ tends when $t$ grows without bound} (which in the example cited above is $1:m$).

This limit can be found by considerations such as those which we have already used in \S~155.

We first form a basis form of the given ideal $\mathfrak a$:
\begin{equation}
\alpha=\alpha_1x_1+\alpha_2x_2+\cdots+\alpha_nx_n,
\tag{1}
\end{equation}
that is, a linear function of the $n$ variables $x_i$ having the property that, by substituting all rational integers for the $x_i$, from $\alpha$ there arise all integral numbers of the field $\Omega$ that are divisible by $\mathfrak a$ (\S~146). By $\alpha_{i,1},\alpha_{i,2},\ldots,\alpha_{i,n}$ we denote the values conjugate to $\alpha_i$, and form the linear substitution for the $x_i$:
\begin{equation}
\begin{aligned}
y_1&=\alpha_{1,1}x_1+\alpha_{2,1}x_2+\cdots+\alpha_{n,1}x_n,\\
y_2&=\alpha_{1,2}x_1+\alpha_{2,2}x_2+\cdots+\alpha_{n,2}x_n,\\
&\cdots\cdots\cdots\\
y_n&=\alpha_{1,n}x_1+\alpha_{2,n}x_2+\cdots+\alpha_{n,n}x_n.
\end{aligned}
\tag{2}
\end{equation}
If $A$ is the determinant of this system of equations,
\[
A=\sum\pm\alpha_{1,1}\alpha_{2,2}\cdots\alpha_{n,n},
\]
then $A^2$ is the discriminant of the system of functions $(\alpha_1,\alpha_2,\ldots,\alpha_n)$, and therefore, by \S~147, is equal to the product of the square of the absolute norm of $\alpha$ and the field discriminant $\Delta$. Since the absolute norm of $\alpha$ is at the same time the norm of the ideal $\mathfrak a$ (\S~151), we put
\begin{equation}
A^2=N(\mathfrak a)^2\Delta.
\tag{3}
\end{equation}

Now we can, as in \S~155, interpret $x_1,x_2,\ldots,x_n$ as coordinates of a point $(x)$ in a space of $n$ dimensions, and let there correspond to each point $(x)$ a point $(x')$ whose coordinates are determined by
\begin{equation}
x_1'=t^{\frac1n}x_1,
\quad x_2'=t^{\frac1n}x_2,
\ldots,
\quad x_n'=t^{\frac1n}x_n,
\tag{4}
\end{equation}
where $t$ denotes an arbitrary positive magnitude. To every region $S$ of points $(x)$ there corresponds a region $S'$ of points $(x')$. All figures in $S$ and the corresponding figures in $S'$ are similar, and the linear dimensions in $S$ are to those in $S'$ as $1:t^{\frac1n}$.

If we think of the region $S$ as given, then the corresponding region $S'$ will depend on $t$ and will increase with $t$. The points with integral coordinates $x_i$ we call, as before, lattice points. The number of lattice points lying in a finite region $S'$ is always finite, but grows with $t$, and shall be denoted by $Z_t$. Then by the theorems of \S~155, (14), the volume of the region $S$, that is, the $n$-fold integral extended over all points of $S$,
\[
\int\!\int\cdots\int dx_1\,dx_2\cdots dx_n,
\]
is equal to the limit of the ratio $Z_t:t$ for $t$ increasing without bound:
\begin{equation}
V=\lim_{t=\infty}\frac{Z_t}{t}.
\tag{5}
\end{equation}

This at the same time gives the more exact precision which we must impose on the region $S$ if our consideration is to be applicable. The region $S$ must be such that the $n$-fold integral $V$ has a definite meaning. Every lattice point is now, by (1), the image of an integral number of the field $\Omega$ divisible by $\mathfrak a$, and thus we can also say that $Z_t$ is the number of integral numbers in $\Omega$ divisible by $\mathfrak a$ whose images lie in the region $S'$.

In order to bound the region $S'$ in a suitable way, we choose a system of fundamental units of the field $\Omega$
\[
\varepsilon_1,\varepsilon_2,\ldots,\varepsilon_{\nu-1}
\]
with the system of conjugate logarithms
\[
l_{i,1},l_{i,2},\ldots,l_{i,\nu}
\]
and with regulator
\begin{equation}
\pm L(\varepsilon_1,\varepsilon_2,\ldots,\varepsilon_{\nu-1})=\pm\sum\pm l_{1,1}l_{2,2}\cdots l_{\nu-1,\nu-1},
\tag{6}
\end{equation}
which may be denoted by $L$.

We now define a system of functions $\xi_1,\xi_2,\ldots,\xi_\nu$ of the variables $x_i$ by the following conditions.

Let, when the $y_i$ have the meaning (2),
\begin{equation}
z_i=\delta_i\log|y_i|,
\tag{7}
\end{equation}
so that, when a system of rational integers is substituted for the $x_i$, $z_i$ passes, by \S~185, into the system of conjugate logarithms of the number $\alpha$.

Then, if $y_k$ belongs to a real field,
\begin{equation}
z_k=\frac12\log y_k^2,
\tag{8}
\end{equation}
and if $y_k,y_k'$ belong to an imaginary pair,
\begin{equation}
z_k=\log y_ky_k',
\tag{9}
\end{equation}
whereby the definition is completely unambiguous.

We now determine the variables $\xi_1,\xi_2,\ldots,\xi_{\nu-1},\xi_\nu$ from the linear equations
\begin{equation}
\begin{aligned}
\xi_1l_{1,1}+\cdots+\xi_{\nu-1}l_{\nu-1,1}+\delta_1\xi_\nu&=z_1,\\
\xi_1l_{1,2}+\cdots+\xi_{\nu-1}l_{\nu-1,2}+\delta_2\xi_\nu&=z_2,\\
&\cdots\cdots\cdots\\
\xi_1l_{1,\nu}+\cdots+\xi_{\nu-1}l_{\nu-1,\nu}+\delta_\nu\xi_\nu&=z_\nu,
\end{aligned}
\tag{10}
\end{equation}
whose determinant does not vanish.

If integral values are given to the variables $x_i$, then by \S~188, (2), (3) the $\xi_1,\xi_2,\ldots,\xi_{\nu-1}$ pass into the exponent system of the number $\alpha$, and by addition of the equations (10) [\S~185, (3)] one obtains
\begin{equation}
\xi_\nu=\frac1n\log \Na(\alpha).
\tag{11}
\end{equation}

We now pass from the point $(x)$ to the point $(x')$ by putting $x_i'=t^{\frac1n}x_i$. Thereby $y_i$ passes into
\[
y_i'=t^{\frac1n}y_i,
\]
and $z_i$ into
\[
z_i'=z_i+\frac1n\delta_i\log t.
\]
If, therefore, we put $z_i'$ in place of $z_i$, and denote by $\xi_1',\xi_2',\ldots,\xi_\nu'$ the values of $\xi_1,\xi_2,\ldots,\xi_\nu$ then obtained from (10), then
\begin{equation}
\xi_1'=\xi_1,
\quad \xi_2'=\xi_2,
\ldots,
\quad \xi_{\nu-1}'=\xi_{\nu-1},
\quad \xi_\nu'=\xi_\nu+\frac1n\log t.
\tag{12}
\end{equation}

We shall now bound the region $S'$ by requiring that
\begin{equation}
0\leq\xi_1'<1,\ldots,
0\leq\xi_{\nu-1}'<1,
\qquad
\xi_\nu'<\frac1n\log t
\tag{13}
\end{equation}
shall hold. In this way we achieve that the lattice points $(x')$ lying in the region $S'$ represent only reduced integral numbers $\alpha$, and all and only those whose absolute norm is smaller than $t$.

Among the reduced numbers, however, there are always $w$ associated with one another, if $w$ denotes the number of roots of unity contained in $\Omega$; and if

we therefore combine these $w$ associated numbers into one complex, then the number of these complexes is equal to the number $T$ of all non-associated integral numbers in $\Omega$ divisible by $\mathfrak a$ whose absolute norm is smaller than $t$. Accordingly the number of lattice points lying in $S'$ is
\begin{equation}
Z_t=wT.
\tag{14}
\end{equation}

For the bounding of the region $S$ one obtains from (12) and (13):
\begin{equation}
0\leq\xi_1<1,\ldots,
0\leq\xi_{\nu-1}<1,
\quad \xi_\nu<0,
\tag{15}
\end{equation}
and from (5) we obtain
\begin{equation}
\lim_{t=\infty}\frac{T}{t}
=\frac1w\int\!\int\cdots\int dx_1\,dx_2\cdots dx_n
=\frac1w V,
\tag{16}
\end{equation}
where the $n$-fold integral is to be extended over the region $S$ determined by (15).

\section*{\S\,190. Determination of the volume.}

The volume $V$ can be determined by the transformation formula for multiple integrals already adduced in \S~155, according to which, if $x_1,x_2,\ldots,x_n$ are functions of the variables $u_1,u_2,\ldots,u_n$,
\begin{equation}
\begin{aligned}
V&=\int\!\int\cdots\int dx_1\,dx_2\cdots dx_n\\
&=\int\!\int\cdots\int\left(\sum\pm
\frac{\partial x_1}{\partial u_1}
\frac{\partial x_2}{\partial u_2}\cdots
\frac{\partial x_n}{\partial u_n}\right)du_1\,du_2\cdots du_n,
\end{aligned}
\tag{1}
\end{equation}
where the functional determinant in the integral taken with respect to the variables $u$ is to be taken with its absolute value.

This functional determinant is also equal to the reciprocal value of the determinant
\[
\sum\pm
\frac{\partial u_1}{\partial x_1}
\frac{\partial u_2}{\partial x_2}\cdots
\frac{\partial u_n}{\partial x_n}.
\]

We first introduce the variables $u$ by a linear substitution, and observe that among the fields conjugate with $\Omega$, according to our assumption, there occur $n-\nu$ conjugate imaginary pairs

and $2\nu-n$ real ones. Accordingly, of the $n$ functions $y$ [\S~189, (2)]
\begin{equation}
y_k=\alpha_{1,k}x_1+\alpha_{2,k}x_2+\cdots+\alpha_{n,k}x_n
\tag{2}
\end{equation}
$2\nu-n$ are real and $2n-2\nu$ are conjugate in pairs. By the $n$ variables $u$ we shall understand real linear combinations of the $x$, namely the $2\nu-n$ real $y$ and the $2n-2\nu$ components of the conjugate imaginary pairs of the $y$, so that, if $y,y'$ is an imaginary pair,
\[
2u_1=y+y',\qquad 2iu_2=y-y',
\]
or
\begin{equation}
y=u_1+iu_2,
\qquad y'=u_1-iu_2
\tag{3}
\end{equation}
is put. If in the functional determinant of the $y$ one applies twice, to each pair of conjugate imaginary rows, the elementary theorem on determinants concerning the addition of rows (Vol.~I, \S~22), then there results
\[
\sum\pm
\frac{\partial u_1}{\partial x_1}
\frac{\partial u_2}{\partial x_2}\cdots
\frac{\partial u_n}{\partial x_n}
=
\frac1{(-2i)^{n-\nu}}
\sum\pm
\frac{\partial y_1}{\partial x_1}
\frac{\partial y_2}{\partial x_2}\cdots
\frac{\partial y_n}{\partial x_n}
=
\frac{A}{(-2i)^{n-\nu}},
\]
where $A$ has the meaning fixed in \S~189, (3).

The field discriminant is positive or negative according as the number of conjugate imaginary pairs, that is, $n-\nu$, is even or odd; consequently, if $\pm\Delta$ denotes the absolute value of the field discriminant, one obtains from \S~189, (4), (16):
\begin{equation}
\lim\frac{T}{t}
=
\frac{2^{n-\nu}}{wN(\mathfrak a)\sqrt{\pm\Delta}}
\int\!\int\cdots\int du_1\,du_2\cdots du_n.
\tag{4}
\end{equation}
It therefore remains only to treat the integral
\[
U=\int\!\int\cdots\int du_1\,du_2\cdots du_n.
\]

From the equations \S~189, (10), for every system of values of the variables $\xi_1,\xi_2,\ldots,\xi_\nu$ satisfying the conditions (15), one obtains a definite finite system of values for $z_1,z_2,\ldots,z_\nu$, and from these, by \S~189, (7), for every $y$ a non-zero absolute value. In the case of an imaginary pair (3), put
\begin{align}
u_1&=e^{\frac12 z}\cos\vartheta,
&u_2&=e^{\frac12 z}\sin\vartheta,
\tag{5}\\
y_1y_1'&=u_1^2+u_2^2=e^z.
\tag{6}
\end{align}
Then, when $\vartheta$ is taken within the bounds
\[
0\leq\vartheta<2\pi,
\]
there results, for fixed $z$, a corresponding pair of values $u_1,u_2$ for each value of $\vartheta$. Such an angle $\vartheta$ belongs to every imaginary pair.

Thus to every point of the region $S$ there belongs one and only one system of values of the variables $\xi_1,\xi_2,\ldots,\xi_\nu,\vartheta\ldots$ within the bounds
\begin{equation}
0\leq\xi_1<1,\ldots,
0\leq\xi_{\nu-1}<1,
\qquad
\xi_\nu<0,
\qquad
0\leq\vartheta<2\pi,\ldots
\tag{7}
\end{equation}

Conversely, if any system of values of the variables $\xi,\vartheta$ is given within the bounds (7), then from it one obtains uniquely the corresponding components $u_1,u_2$ of an imaginary $y$, but of the real $y$ only the absolute values; hence to each real $y$ either of the two signs can still be given. The number of these possible determinations is $2^{2\nu-n}$. If one of these determinations is selected, the variables $x$ are uniquely determined and are finite for every finite system of values of the $\xi,\vartheta$. For $\xi_\nu$ decreasing toward $-\infty$, the values of the $y$ will in part converge to zero.

Accordingly the region $S$ decomposes into $2^{2\nu-n}$ subregions $S_1,S_2,\ldots$, each of which is characterized by a definite combination of signs of the real $y_1,y_2,\ldots$; and accordingly the integral $U$ decomposes into just as many partial integrals $U_1,U_2,\ldots$,
\[
U=U_1+U_2+\cdots.
\]
In carrying out the integration in one of these components, which all turn out to be of the same value, we begin with an imaginary pair, if one is present.

If we regard $u_1,u_2$ as rectangular coordinates in a plane, then by (5) we can regard $e^{\frac12 z},\vartheta$ as polar coordinates, and obtain
\[
\int\!\int du_1\,du_2=\frac12\int\!\int e^z\,dz\,d\vartheta,
\]
where the integration with respect to $\vartheta$, which extends from 0 to $2\pi$, can be carried out. Thus one finds
\[
\int\!\int du_1\,du_2=\pi\int e^z\,dz.
\]

One now changes the order of integration and puts first any second imaginary pair that may still be present, which is again treated in the same way.

The different components $U_1,U_2,\ldots$ differ by the signs of the $u$ corresponding to the real $y$. We take the sign so that $\pm u$ becomes positive, and put
\begin{equation}
z=\log(\pm u),
\qquad \pm du=e^z\,dz.
\tag{8}
\end{equation}
Thus for the different $U_1,U_2,\ldots$ we obtain the same expression by a $\nu$-fold integral,
\begin{equation}
U_1=\pi^{n-\nu}\int\!\int\cdots\int e^{z_1+z_2+\cdots+z_\nu}\,dz_1\,dz_2\cdots dz_\nu,
\tag{9}
\end{equation}
if we take $dz_1,dz_2,\ldots,dz_\nu$ to be positive.

But by (6) and (8) the variables $z_k$ are here no others than those in \S~189, (8), (9), (10), and hence
\[
z_1+z_2+\cdots+z_\nu=n\xi_\nu.
\]
If we now apply the transformation formula (1) to the $\nu$-fold integral (9), in order to introduce the variables $\xi_1,\xi_2,\ldots,\xi_\nu$, then we have to form the determinant
\[
\sum\pm
\frac{\partial z_1}{\partial\xi_1}
\frac{\partial z_2}{\partial\xi_2}\cdots
\frac{\partial z_\nu}{\partial\xi_\nu},
\]
which by \S~188, (3), agrees in absolute value with $nL$, that is, with the $n$-fold value of the field regulator $L$; and therefore
\[
U_1=\pi^{n-\nu}nL\int_0^1d\xi_1\cdots\int_0^1d\xi_{\nu-1}\int_{-\infty}^{0}e^{n\xi_\nu}\,d\xi_\nu
=\pi^{n-\nu}L
\]
and
\[
U=2^{2\nu-n}\pi^{n-\nu}L.
\]
Finally, therefore, (4) gives the theorem:

\begin{center}
\textbf{1. If $T$ denotes the number of non-associated integral numbers divisible by $\mathfrak a$ whose absolute norm is smaller than $t$, then}
\end{center}
\begin{equation}
\lim_{t=\infty}\frac{T}{t}
=
\frac{2^\nu\pi^{n-\nu}L}{wN(\mathfrak a)\sqrt{\pm\Delta}}
=
\frac{g}{N(\mathfrak a)},
\tag{10}
\end{equation}
\begin{center}
\textbf{where $g$ is a positive number completely determined by the nature of the field $\Omega$, namely:}
\end{center}
\begin{equation}
g=\frac{2^\nu\pi^{n-\nu}L}{w\sqrt{\pm\Delta}}.
\tag{11}
\end{equation}

\clearpage
\section*{\S\,191. Theorems from the theory of series.}

For the further applications of the preceding results, some theorems from the theory of infinite series are required; these shall first be derived here.

\begin{enumerate}[label=\textbf{\arabic*.},leftmargin=2.2em]
\item Let $a_1,a_2,a_3,\ldots,a_n,\ldots$ be an unbounded system of real quantities (positive, negative, or also vanishing), of which we suppose that the sum
\begin{equation}
\sigma_n=a_1+a_2+\cdots+a_n
\tag{1}
\end{equation}
remains, for every arbitrary $n$, below a finite bound $C$ in absolute value. Then the series
\begin{equation}
S=\frac{a_1}{1^s}+\frac{a_2}{2^s}+\frac{a_3}{3^s}+\cdots
\tag{2}
\end{equation}
is, for every positive $s$, not only convergent, but also a continuous function of $s$.
\end{enumerate}

To prove this theorem, we decompose $S$ in this way:
\[
S=S_m+R_m,
\]
where
\[
S_m=\frac{a_1}{1^s}+\frac{a_2}{2^s}+\cdots+\frac{a_{m-1}}{(m-1)^s},
\]
\[
R_m=\frac{a_m}{m^s}+\frac{a_{m+1}}{(m+1)^s}+\frac{a_{m+2}}{(m+2)^s}+\cdots .
\]
Now by (1) one has
\[
a_m=\sigma_m-\sigma_{m-1},\quad a_{m+1}=\sigma_{m+1}-\sigma_m,\quad
a_{m+2}=\sigma_{m+2}-\sigma_{m+1},\ldots
\]
and therefore
\[
\begin{aligned}
R_m+\frac{\sigma_{m-1}}{m^s}
&=\sigma_m\left(\frac1{m^s}-\frac1{(m+1)^s}\right)\\
&\quad+\sigma_{m+1}\left(\frac1{(m+1)^s}-\frac1{(m+2)^s}\right)+\cdots .
\end{aligned}
\]
Since, by the assumption, $\sigma_{m-1},\sigma_m,\sigma_{m+1},\ldots$ are enclosed between finite bounds $\pm C$, it follows that $R_m+\frac{\sigma_{m-1}}{m^s}$ is enclosed between the two bounds
\[
\pm C\left(\frac1{m^s}-\frac1{(m+1)^s}+\frac1{(m+1)^s}-\frac1{(m+2)^s}+\cdots\right)
=\pm\frac{C}{m^s}
\]

and hence, in absolute value,
\[
R_m<\frac{2C}{m^s},
\]
or, if $s>c$,
\[
R_m<\frac{2C}{m^c}.
\]

This upper bound for $R_m$, which is independent of $s$, can however, if $c$ is positive, be depressed below any arbitrarily small value by taking $m$ sufficiently large.

From this follows not only the convergence, but also the continuity of $S$. For if one takes $m$ sufficiently large, then not only $R_m$, but also the fluctuation of $R_m$ for variable $s$, becomes infinitely small, and $S_m$ is, for a fixed $m$, a continuous function of $s$. Thus $S$ too is continuous.

It is also to be observed that the assumption on $\sigma_n$ by no means presupposes the convergence of the infinite series $\sum a_k$. It is not even necessary that the $a_k$ be real; for if they are imaginary, one need only apply the theorem to the real and to the imaginary component. Finally it is also not necessary to assume $a_1,a_2,a_3,\ldots$ to be constants. Everything remains valid when they are continuous functions of $s$.

\begin{enumerate}[label=\textbf{\arabic*.},leftmargin=2.2em,start=2]
\item Let $\mu_1,\mu_2,\ldots,\mu_n,\ldots$ denote an infinite set of positive numbers of such a kind that two finite numbers $\alpha,\beta$ can be given so that for every however large $n$
\begin{equation}
\alpha<\frac{n}{\mu_n}<\beta,
\tag{3}
\end{equation}
then the infinite series
\[
S=\frac1{\mu_1^s}+\frac1{\mu_2^s}+\frac1{\mu_3^s}+\cdots
\]
is convergent so long as the exponent $s$ is greater than $1$, and the product $(s-1)S$, as $s$ approaches the value $1$ sufficiently closely, likewise remains between the bounds $\alpha$ and $\beta$, or at least approaches one of these values without bound.
\end{enumerate}

The proof is very simple to conduct by reduction to an integral. For plainly
\[
\frac1{(n+1)^s}<\int_n^{n+1}\frac{dx}{x^s}<\frac1{n^s},
\]
because the first or the second value, $(n+1)^{-s}$ or $n^{-s}$, is obtained when, under the integral sign, the too small or the too large value $(n+1)^{-s}$ or $n^{-s}$ is put for $x^{-s}$.

Thus from (3), for every positive $s$, one obtains
\[
\alpha^s\int_n^{n+1}\frac{dx}{x^s}<\frac1{\mu_n^s}<\beta^s\int_{n-1}^{n}\frac{dx}{x^s}.
\]
If we put
\[
R_m=\frac1{\mu_{m+1}^s}+\frac1{\mu_{m+2}^s}+\frac1{\mu_{m+3}^s}+\cdots,
\]
then this gives
\[
\alpha^s\int_{m+1}^{\infty}\frac{dx}{x^s}<R_m<\beta^s\int_m^{\infty}\frac{dx}{x^s}
\]
or
\begin{equation}
\frac{\alpha^s}{(s-1)(m+1)^{s-1}}<R_m<\frac{\beta^s}{(s-1)m^{s-1}} .
\tag{4}
\end{equation}

Thus $R_m$, however many terms may be summed, remains finite if $s>1$, and for series with positive terms this is a sufficient condition for convergence.

If one further puts (4) in the form
\[
\frac{\alpha^s}{(m+1)^{s-1}}<(s-1)R_m<\frac{\beta^s}{m^{s-1}},
\]
one sees that the two bounds can be brought arbitrarily close to $\alpha,\beta$ if one takes $s-1$ small enough, and hence it follows that $(s-1)R_m$, for sufficiently small $s-1$, cannot lie outside the bounds $\alpha,\beta$ unless, as $s-1$ decreases, it approaches one of these values without bound. Since now $R_m$ differs from $S$ only by a component which is finite for every positive $s$, the theorem is thereby proved.

The theorem remains valid also when the inequality (3) does not hold for all $n$, but only holds as soon as $n$ has exceeded an arbitrarily assigned bound $m$. On

the basis of this remark, if $n:\mu_n$ approaches a finite limiting value $\gamma$, one may take $\alpha$ and $\beta$ arbitrarily close to this limiting value, and in this way obtains the following somewhat more precise form of the theorem:

\begin{enumerate}[label=\textbf{\arabic*.},leftmargin=2.2em,start=3]
\item If the positive numbers $\mu_1,\mu_2,\mu_3,\ldots$ are so constituted that
\[
\lim_{n=\infty}\frac{n}{\mu_n}=\gamma
\]
is a finite limiting value, then the infinite series
\[
S=\frac1{\mu_1^s}+\frac1{\mu_2^s}+\frac1{\mu_3^s}+\cdots
\]
is convergent for every $s$ greater than $1$, and
\[
\lim_{s=1}(s-1)S=\gamma.
\]
\end{enumerate}

Let now any system $M$ of infinitely many positive numbers $\mu_n$ be given, and let us arrange them as
\begin{equation}
\mu_1\leqq\mu_2\leqq\mu_3\leqq\cdots\quad (M),
\tag{5}
\end{equation}
so that in the series (5) a smaller term never follows a larger one; and let $T$ denote the number of terms in $M$ which are not greater than an arbitrarily given positive quantity $t$. Then the following theorem holds:

\begin{enumerate}[label=\textbf{\arabic*.},leftmargin=2.2em,start=4]
\item If one of the two limiting values
\begin{equation}
\lim_{n=\infty}\frac{n}{\mu_n},\qquad \lim_{t=\infty}\frac{T}{t}
\tag{6}
\end{equation}
is finite, then the other has the same finite value.
\end{enumerate}

First let
\begin{equation}
\lim_{n=\infty}\frac{n}{\mu_n}=\gamma
\tag{7}
\end{equation}
be a finite limiting value. Then the $\mu_n$ must necessarily grow to infinity with infinitely increasing $n$, and for every positive $t$ there is a value $m$ such that
\begin{equation}
\mu_m\leqq t<\mu_{m+1};
\tag{8}
\end{equation}
$m$ grows to infinity at the same time as $t$, and $T=m$ [by (5)]. Therefore by (8)
\begin{equation}
\frac{m}{\mu_m}\geqq\frac{T}{t}>\frac{m}{\mu_{m+1}} .
\tag{9}
\end{equation}

If now $\gamma=0$, then from this, letting $t$ and $m$ grow to infinity, it follows that $T:t$ also has the limiting value $0$. But if $\gamma$ is different from $0$, then by (7)
\[
\lim\frac{\mu_m}{\mu_{m+1}}=\lim\frac{m}{m+1}=1,
\]
and therefore the two limiting values in (9) approach the value $\gamma$, and hence it follows that
\begin{equation}
\lim_{t=\infty}\frac{T}{t}=\gamma.
\tag{10}
\end{equation}

Conversely, suppose secondly that the limit equation (10) is given. Then here too $\mu_n$ must grow to infinity with $n$, for otherwise, since the total number of all $\mu$ is infinite, $T$ would already be infinite for a finite $t$, and $\gamma$ could not be finite.

Now take a value $\mu_n$ which occurs in the series of mutually equal values
\[
\mu_{m+1},\mu_{m+2},\ldots,\mu_{m+l}
\]
so that
\[
m+1\leqq n\leqq m+l.
\]
Then $T$, when $t$ passes through the value $\mu_n$, will suddenly increase by $l$ units, and the ratio $T:t$ increases by $l:\mu_n$. But since $T:t$ is to have a finite limiting value, one must have
\begin{equation}
\lim\frac{l}{\mu_n}=0.
\tag{11}
\end{equation}
If now $t=\mu_n$, then $T=m+l$, and consequently
\begin{equation}
\frac{T}{t}=\frac{m+l}{\mu_n},\qquad
\lim\frac{m+l}{\mu_n}=\gamma.
\tag{12}
\end{equation}
Further,
\[
\frac{m}{\mu_n}<\frac{n}{\mu_n}\leqq\frac{m+l}{\mu_n},
\]
and consequently, by (11) and (12),
\[
\lim\frac{n}{\mu_n}=\lim\frac{m}{\mu_n}=\gamma.
\]
Thus equation (7) is derived from equation (10)\footnote{This theorem is proved in essentially the same way in Dirichlet--Dedekind, Supplement II. Compare also Riemann's mathematical works, second edition, No. XXX.}.

Using theorem 4, we can now also give theorem 3 the form:

\begin{enumerate}[label=\textbf{\arabic*.},leftmargin=2.2em,start=5]
\item If $T$ is the number of quantities $\mu_n$ that are not greater than $t$, then
\begin{equation}
\lim_{s=1}\sum_{n=1}^{\infty}\frac{s-1}{\mu_n^s}=\lim_{t=\infty}\frac{T}{t},
\tag{13}
\end{equation}
provided that the limiting value of $T:t$ is finite.
\end{enumerate}

Assume that the quantities $\mu_1,\mu_2,\mu_3,\ldots$ satisfy not only the conditions of theorem 3, but also the further conditions that, when
\begin{equation}
\mu_n=\frac{n+c_n}{\gamma}
\tag{14}
\end{equation}
is put, the $c_n$ do not become infinite with infinitely increasing $n$ (that is, for every however large $n$ they are enclosed between finite bounds). Then theorem 3 can be replaced by the following sharper one:

\begin{enumerate}[label=\textbf{\arabic*.},leftmargin=2.2em,start=6]
\item If the numbers $\mu_n$ have the property that a finite $\gamma$ can be determined so that
\begin{equation}
\gamma\mu_n-n=c_n
\tag{15}
\end{equation}
does not become infinite with infinitely increasing $n$, then the function of $s$ defined for every $s>1$ by
\begin{equation}
S=\frac1{\mu_1^s}+\frac1{\mu_2^s}+\cdots
\tag{16}
\end{equation}
has the property that the difference
\begin{equation}
S-\frac{\gamma}{s-1}=C
\tag{17}
\end{equation}
approaches a finite limit as $s-1$ decreases without limit.
\end{enumerate}

To prove this, put, by (14) and (16),
\begin{equation}
\sum_{n=1}^{\infty}\frac{\gamma^s}{n^s}-S
=\gamma^s\sum_{n=1}^{\infty}\frac1{n^s}
\left[1-\left(1+\frac{c_n}{n}\right)^{-s}\right].
\tag{18}
\end{equation}
If now $\varepsilon$ is a positive proper fraction, then, putting
\[
a_n=\frac1{n^\varepsilon}\left[1-\left(1+\frac{c_n}{n}\right)^{-s}\right],
\]
the product $a_n n^{1+\varepsilon}$ does not become infinite for infinitely increasing $n$, and consequently, by a known elementary theorem of the theory of series, $a_1+a_2+\cdots$ is an unconditionally convergent series.

Now put
\[
s=s_1+\varepsilon.
\]
Then from (18) one obtains
\[
\sum\frac{\gamma^s}{n^s}-S=\gamma^s\sum\frac{a_n}{n^{s_1}},
\]
and this, by 1., is a finite and continuous function of $s_1$ for positive $s_1$. But for $s_1=1-\varepsilon$ one has $s=1$, and consequently
\begin{equation}
\sum\frac{\gamma^s}{n^s}-S=D_s
\tag{19}
\end{equation}
is finite for $s=1$, namely equal to
\[
D_1=\sum\frac{\gamma c_n}{n(n+c_n)}.
\]

By application of simple theorems from the theory of the $\Gamma$-function, however, the sum $\sum\frac1{n^s}$ can be expressed by a definite integral. One has
\[
\frac1{n^s}=\frac1{\Gamma(s)}\int_0^\infty e^{-nx}x^{s-1}\,dx,
\]
and consequently
\[
\sum\frac1{n^s}=\frac1{\Gamma(s)}\int_0^\infty \frac{x^{s-1}e^{-x}\,dx}{1-e^{-x}}.
\]
Further [by the theorem $(s-1)\Gamma(s-1)=\Gamma(s)$],
\[
\frac1{s-1}=\frac1{\Gamma(s)}\int_0^\infty e^{-x}x^{s-2}\,dx,
\]
and hence
\[
\sum\frac1{n^s}-\frac1{s-1}
=\frac1{\Gamma(s)}\int_0^\infty x^{s-1}e^{-x}
\left(\frac1{1-e^{-x}}-\frac1x\right)dx.
\]
Here the right-hand side is a continuous function for all positive values of $s$, and for $s=1$ passes into Euler's constant
\[
\int_0^\infty e^{-x}\left(\frac1{1-e^{-x}}-\frac1x\right)dx
=-\Gamma'(1)=0.57721566\ldots .
\]
Now by (19)
\[
D_s=\gamma^s\left(\sum\frac1{n^s}-\frac1{s-1}\right)
-\left(S-\frac{\gamma^s}{s-1}\right),
\]

and therefore
\begin{equation}
\lim_{s=1}\left(S-\frac{\gamma^s}{s-1}\right)
=-\gamma\Gamma'(1)+D_1=C_1,
\tag{20}
\end{equation}
which was to be proved, is finite\footnote{The theorems on $\Gamma$-functions used here are found in the more detailed textbooks of integral calculus, for example Serret--Harnack, \emph{Lehrbuch der Differential- und Integralrechnung}, vol. II, p. 170 f. (Leipzig 1885).}.
\section*{\S\,192. Application to the determination of the class number.}

We now apply these theorems to the theory of the field $\Omega$. By the numbers $\mu_n$ of theorem 5 we understand the norms of all ideals of the field $\Omega$, and therefore by $T$ the number of ideals in $\Omega$ whose norm is not greater than a given positive quantity $t$; and we obtain
\begin{equation}
\lim_{s=1}\sum\frac{s-1}{N(\mathfrak a)^s}=\lim_{t=\infty}\frac{T}{t},
\tag{1}
\end{equation}
where the sum on the left-hand side extends over all ideals $\mathfrak a$ of the field. Now by \S~153 the ideals split into a finite number of classes
\[
A_1,A_2,\ldots,A_h,
\]
and the great goal is the determination of this number $h$, the class number.

The ideals $\mathfrak a_1$ of one of these classes $A_1$ are characterized by the fact that their products with one and the same integral functional $\varphi_1$, belonging to the class $A_1^{-1}$, are principal ideals, that is,
\[
\varphi_1\mathfrak a_1=\alpha
\]
is an integral number of the field $\Omega$.

If the norm of $\mathfrak a_1$ is not greater than $t$, then
\[
N_a(\alpha)\leqq N_a(\varphi_1)t=t_1.
\]
If $T_1$ is the number of ideals of the class $A_1$ whose norms are not greater than $t$, then $T_1$ is at the same time the number of non-associated integral numbers divisible by $\varphi_1$ whose norms are not

greater than $t_1$; and by the theorem \S~190, 1,
\begin{equation}
\lim_{t_1=\infty}\frac{T_1}{t_1}=\frac{g}{N_a(\varphi_1)},
\tag{2}
\end{equation}
where $g$ has the meaning given at the cited place, and is therefore a non-zero number determined by the nature of the field $\Omega$.

If now $T_2,t_2,\ldots,T_h,t_h$ have the corresponding meaning for the classes $A_2,\ldots,A_h$, as $T_1,t_1$ have for $A_1$, then
\[
T=T_1+T_2+\cdots+T_h,
\]
\[
\frac{T}{t}=\frac{T_1}{t_1}N_a(\varphi_1)+\frac{T_2}{t_2}N_a(\varphi_2)+\cdots+
\frac{T_h}{t_h}N_a(\varphi_h),
\]
and from (1) and (2) there results the fundamental formula
\begin{equation}
\lim_{s=1}\sum\frac{s-1}{N(\mathfrak a)^s}=gh.
\tag{3}
\end{equation}

The number $g$ may be regarded as known by \S~185, (11), even though its actual calculation still rests on the assumption that a fundamental system of units is known, and even though the determination of such a system in higher fields has always been regarded as one of the greatest difficulties. Then the calculation of the class number still depends on determining the limiting value on the left-hand side of (3), which of course succeeds only in special cases, but often permits important conclusions concerning the nature of the class number. We shall still develop a transformation of the sum
\begin{equation}
\sum\frac1{N(\mathfrak a)^s}=\Phi(s)
\tag{4}
\end{equation}
into an infinite product, preparatory to the calculation.

Let $\mathfrak p$ be any prime ideal of $f$-th degree in the field $\Omega$, so that
\[
N(\mathfrak p)=p^f.
\]
Then the infinite geometric series is
\[
1+\frac1{N(\mathfrak p)^s}+\frac1{N(\mathfrak p)^{2s}}+\cdots
=\frac1{1-p^{-sf}};
\]
if these series are multiplied together for all distinct prime ideals $\mathfrak p$, then, by known theorems from

the theory of infinite series, one obtains a sum of terms of the form
\[
\left(\frac1{N(\mathfrak p)^x}\frac1{N(\mathfrak p')^{x'}}\frac1{N(\mathfrak p'')^{x''}}\cdots\right)^s
=\left(\frac1{N(\mathfrak p^x\mathfrak p'^{x'}\mathfrak p''^{x''}\cdots)}\right)^s,
\]
which are all contained in the form $N(\mathfrak a)^{-s}$, and every such term arises once and only once. Thus
\begin{equation}
\Phi(s)=\prod_{\mathfrak p}\frac1{1-N(\mathfrak p)^{-s}}.
\tag{5}
\end{equation}
If therefore $\mathfrak p_1,\mathfrak p_2,\ldots,\mathfrak p_e$ are the mutually distinct prime factors of $p$, and $f_1,f_2,\ldots,f_e$ their degrees (\S~143), then
\begin{equation}
\Phi(s)=\prod_p
\frac1{(1-p^{-sf_1})(1-p^{-sf_2})\cdots(1-p^{-sf_e})},
\tag{6}
\end{equation}
where the infinite product $\prod$ is to be extended over all natural prime numbers $p$, and by (3), (4)
\begin{equation}
gh=\lim_{s=1}(s-1)\Phi(s).
\tag{7}
\end{equation}

From this one can draw a conclusion important for many applications:

If from the development
\[
1+\frac1{p^s}+\frac1{p^{2s}}+\cdots=\frac1{1-p^{-s}}
\]
we form the product for all prime numbers $p$, then we get
\begin{equation}
\prod_p\frac1{1-p^{-s}}=\sum_{n=1}^{\infty}\frac1{n^s},
\tag{8}
\end{equation}
where the sum on the right-hand side extends over all natural numbers $n$; and thus this product has a finite value for all values of $s$ greater than $1$. The reciprocal value, namely the product
\begin{equation}
P=\prod_p(1-p^{-s}),
\tag{9}
\end{equation}
whose factors are all smaller than $1$, nevertheless has, so long as $s>1$, a value different from zero. This property is preserved if in forming the product $P$ we consider only a part of all prime numbers $p$, since the product is only increased by omitting arbitrary factors, without ever exceeding unity.

It follows that in the product (6), the partial products
\[
\prod\frac1{1-p^{-sf}},
\]
which extend over all prime numbers $p$ for which $f>1$, remain finite even for $s=1$. Since, on the other hand, $g,h$ have definite non-zero values, the following important theorem follows from (7):

\begin{center}
\textbf{I. If $\mathfrak p$ runs through the totality of prime ideals of first degree of any algebraic field $\Omega$, then the product}
\end{center}
\begin{equation}
(s-1)\prod\frac1{1-N(\mathfrak p)^{-s}}
\tag{10}
\end{equation}
\begin{center}
\textbf{has, for $s=1$, a finite limiting value different from zero.}
\end{center}

This property is preserved also if, in forming the product, an arbitrary finite number of prime ideals of first degree is omitted.

The norms $N(\mathfrak p)$ in formula (10) are equal to natural prime numbers $p$. But a prime number $p$ occurs in it as often as it contains prime factors of first degree in $\Omega$, hence at most $n$ times. If $\Omega$ is a normal field, then every prime number $p$ which is decomposable into prime factors of first degree actually occurs in it exactly $n$ times, and for the product (10) we may also write
\[
(s-1)\prod\frac1{(1-p^{-s})^n},
\]
where the product $\prod$ extends over all prime numbers $p$ decomposable into prime factors of first degree (\S~160).

An immediate consequence of theorem I is:
\begin{center}
\textbf{In every algebraic field there are infinitely many prime ideals of first degree.}
\end{center}

\section*{Twenty-fourth Section.\\Class number of the cyclotomic fields.}
\section*{\S\,193. Class-number representation in the cyclotomic field $\Omega_m$.}

The general results concerning the class number $h$ shall now be applied to the full cyclotomic field $\Omega_m$, under the standing assumption that $m$ is a power of a prime number $q$, and hence that
\begin{equation}
m=q^\alpha,\qquad \varphi(m)=q^{\alpha-1}(q-1)=\mu
\tag{1}
\end{equation}
has been put.

In the twenty-first section (\S~169, 171) we saw that in $q$ only one prime ideal of first degree occurs, and that a prime number $p$ different from $q$ and belonging to the exponent $f$ decomposes into $e$ distinct prime ideals of $f$-th degree. Here $f$ denotes the smallest positive exponent for which
\begin{equation}
p^f\equiv1\pmod m
\tag{2}
\end{equation}
and $e$ is determined by the equation
\begin{equation}
\mu=ef.
\tag{3}
\end{equation}

The function $\Phi(s)$ determined at the end of the preceding paragraph by formula (6) therefore here obtains the expression
\begin{equation}
\Phi(s)=\frac1{1-q^{-s}}\prod\frac1{(1-p^{-sf})^e},
\tag{4}
\end{equation}
where the product sign $\prod$ extends over all prime numbers $p$ different from $q$. Here $s$ is a variable which is always greater than $1$, and is finally to approach the bound $1$.

In order now to transform this expression for the function $\Phi$ further and prepare it for calculation, one must use the theorems on Abelian groups and group characters which we learned in paragraphs 14 to 16 of this volume.

The totality of the numbers $n$ not divisible by $q$, taken modulo $m$, forms an Abelian group $\mathfrak R$ of degree $\mu$, whose characters $\chi(n)$ were determined in \S~16. There was a small difference according as $q$ is odd or $q=2$.

1) If $q$ is odd, then there is a primitive root $c$ of $m$, and for every number $n$ there is an index $\gamma$ such that
\begin{equation}
n\equiv c^\gamma\pmod m.
\tag{5}
\end{equation}
If then $\Theta$ is a $\mu$-th root of unity, then
\begin{equation}
\chi(n)=\Theta^\gamma,
\tag{6}
\end{equation}
and all $\mu$ characters are obtained by putting for $\Theta$ the different $\mu$-th roots of unity.

If $n$ belongs to the exponent $f$, then $e$ is the greatest common divisor of $\gamma$ and $\mu$, and $\chi(n)$ is an $f$-th root of unity.

2) For $q=2$, $m>4$ (the case $m=4$ is not considered here), every odd number $n$ has two indices $\alpha,\beta$, determined modulo $2,\frac12\mu$, for which
\begin{equation}
n\equiv(-1)^\alpha 5^\beta\pmod m
\tag{7}
\end{equation}
holds. As characters one obtains, if $\varepsilon=\pm1$ and $\Theta$ is put equal to a $\frac12\mu$-th root of unity,
\begin{equation}
\chi(n)=\varepsilon^\alpha\Theta^\beta;
\tag{8}
\end{equation}
$f$ is the smallest positive solution of the two congruences
\[
\alpha f\equiv0\pmod2,\qquad \beta f\equiv0\pmod{\tfrac12\mu},
\]
and here also $\chi(n)$ is an $f$-th root of unity.

For both cases, however, the following remark holds:

3) If $n$ belongs to the exponent $f$, then, when $\chi$ is allowed to run through the whole group of the $\mu$ characters, every $f$-th root of unity is obtained $e$ times from $\chi(n)$.

For since the $\chi$ form an Abelian group, one first obtains each $f$-th root of unity which occurs at all equally often, namely as often as the value $1$, since if $\chi_1(n)=\chi_2(n)$,

then $\chi_1\chi_2^{-1}(n)=\chi_0(n)=1$. On the other hand, however, every $f$-th root of unity is obtained; for if in case 1) $\Theta$ is taken to be a primitive $\mu$-th root of unity, and in case 2) a primitive $\frac12\mu$-th root of unity, and in the latter case at the same time $\varepsilon=-1$, then $\chi(n)$ is a primitive $f$-th root of unity, and from its powers all $f$-th roots of unity can be derived.

If now $x$ denotes a variable, then, when $n$ belongs to the exponent $f$, the product extended over all characters $\chi$ is
\begin{equation}
\prod_\chi[1-\chi(n)x]=(1-x^f)^e,
\tag{9}
\end{equation}
and if one therefore puts $x=p^{-s}$, then from (4) one obtains
\begin{equation}
\Phi(s)=\frac1{1-q^{-s}}\prod_\chi\prod_p\frac1{1-\chi(p)p^{-s}},
\tag{10}
\end{equation}
where the first product sign extends over all characters $\chi$, and the second over all prime numbers $p$. Thus $\Phi(s)$ is a product of a finite number, $\mu$, of factors, each of which is an infinite product.

We now expand each factor of one of these infinite products in increasing powers of $p^{-s}$. This gives
\[
\frac1{1-\chi(p)p^{-s}}
=1+\chi(p)p^{-s}+\chi(p^2)p^{-2s}+\chi(p^3)p^{-3s}+\cdots .
\]
The product of all factors so obtained, if $\chi$ is kept fixed while $p$ runs through all primes different from $q$, is an aggregate of terms of the form
\[
\chi(p^k)\chi(p'^{\,k'})\chi(p''^{\,k''})\cdots p^{-sk}p'^{-sk'}p''^{-sk''}\cdots
=\chi(n)n^{-s},
\]
and every term of this form occurs in the product once and only once [compare \S~192, (8)]. Thus the transformation of $\Phi(s)$ into a finite product of infinite series is obtained:
\begin{equation}
\Phi(s)=\frac1{1-q^{-s}}\prod_\chi\sum_n\frac{\chi(n)}{n^s},
\tag{11}
\end{equation}
where the sum extends over all numbers $n$ not divisible by $q$, and the product over all $\mu$ characters $\chi$.

This expression is suited to determine the limiting value of $(s-1)\Phi(s)$ for $s=1$. Let us first consider the sum corresponding to the principal character $\chi=1$:
\[
\sum_n\frac1{n^s}.
\]

The numbers $n$ are obtained by removing from the totality of all natural numbers $k$ the multiples of $q$, that is, the totality of the numbers $qk$. Hence
\[
\sum_n\frac1{n^s}=\sum_k\frac1{k^s}-\sum_k\frac1{q^s k^s}
=(1-q^{-s})\sum_k\frac1{k^s}.
\]
But if in theorem 3, \S~191, we put for the $\mu_n$ the sequence of natural numbers $k$, then
\[
\lim_{s=1}(s-1)\sum_k\frac1{k^s}=1,
\]
and therefore we obtain
\begin{equation}
\lim_{s=1}\frac{s-1}{1-q^{-s}}\sum_n\frac1{n^s}=1.
\tag{12}
\end{equation}

The other factors of the product $\Phi(s)$, however, are, when the infinite series are ordered according to increasing values of $n$, continuous functions of $s$ by theorem 1, \S~191.

For the sum ordered according to increasing values of $n$,
\begin{equation}
\sum\chi(n),
\tag{13}
\end{equation}
is indeed not convergent, but nevertheless cannot go beyond a finite bound in absolute value. For as often as $n$ runs through a complete residue system modulo $m$, a contribution is added to the sum (13) which, by \S~11, 6, has the value $0$. Hence, if $\chi$ is not the principal character,
\begin{equation}
\lim_{s=1}\sum_n\frac{\chi(n)}{n^s}=\sum_n\frac{\chi(n)}{n},
\tag{14}
\end{equation}
and consequently, for the class number $h$ in the cyclotomic field $\Omega_m$, we now obtain the expression
\begin{equation}
gh=\prod_\chi\sum_n\frac{\chi(n)}{n},
\tag{15}
\end{equation}
in which the product sign $\prod$ refers only to the characters $\chi$ different from the principal character.

\section*{\S\,194. Determination of the sums $X$.}

The infinite series
\begin{equation}
X=\sum_n\frac{\chi(n)}{n},
\tag{1}
\end{equation}

on which, according to the preceding, the determination of the class number still depends, can be represented by finite expressions. Namely,
\[
\frac1n=\int_0^1x^{n-1}\,dx,
\]
and therefore
\[
X=\int_0^1\sum_n\chi(n)x^{n-1}\,dx.
\]
Now $\chi(n)$ remains unchanged when $n$ is increased by a multiple of $m$. If, however, $t$ denotes the least positive residue of $n$ and one puts
\[
n=t+ml,
\]
then $\chi(n)=\chi(t)$, and it follows that
\[
X=\int_0^1\sum_t\chi(t)x^{t-1}\sum_{l=0}^{\infty}x^{ml}\,dx.
\]
Put now
\begin{equation}
f(x)=\sum_t\chi(t)x^t.
\tag{2}
\end{equation}
Then $f(x)$ is an integral function of $x$ of degree $m-1$, and at the same time, by the relation $\sum\chi(t)=0$,
\begin{equation}
f(0)=0,\qquad f(1)=0,
\tag{3}
\end{equation}
so that $f(x)$ is divisible by $x(1-x)$. For $X$, by the summation formula for the geometric series $\sum x^{ml}$, one obtains
\begin{equation}
X=\int_0^1\frac{f(x)\,dx}{x(1-x^m)}.
\tag{4}
\end{equation}

To find such an integral, integral calculus prescribes decomposing the rational fraction under the integral sign into partial fractions. This partial-fraction decomposition gives, however, by (3), if
\[
r=e^{\frac{2\pi i}{m}}
\]
is put:
\[
\frac{f(x)}{x(1-x^m)}
=-\frac1m\sum_{s=1}^{m-1}\frac{f(r^s)}{x-r^s}
\]
[Vol. I, \S~29, (11)], and now we have further, by known

elementary theorems of integral calculus, if $\alpha$ denotes any angle between $0$ and $2\pi$,
\[
\int_0^1\frac{dx}{x-e^{i\alpha}}
=\frac12\log\left(4\sin^2\frac{\alpha}{2}\right)
+i\,\frac{\pi-\alpha}{2}.
\]
It follows from this, by (4), that
\begin{equation}
X=-\frac1m\sum_{s=1}^{m-1}f(r^s)
\left[\frac12\log\left(4\left(\sin\frac{s\pi}{m}\right)^2\right)-\frac{i\pi(m-2s)}{2m}\right].
\tag{5}
\end{equation}
Here by (2)
\begin{equation}
f(r^s)=\sum_t\chi(t)r^{st},
\tag{6}
\end{equation}
and hence [by (3)]
\begin{equation}
\sum_{s=1}^{m-1}f(r^s)=\sum_{s=0}^{m-1}f(r^s)=0,
\tag{7}
\end{equation}
since
\[
\sum_{s=0}^{m-1}r^{st}=0.
\]
Accordingly formula (5) simplifies somewhat further and gives
\begin{equation}
X=-\frac1m\sum_{s=1}^{m-1}f(r^s)
\left[\frac12\log\left(\sin\frac{s\pi}{m}\right)^2+\frac{i\pi s}{m}\right],
\tag{8}
\end{equation}
for which one may also, replacing $s$ by $m-s$ and again using (7), put
\begin{equation}
X=-\frac1m\sum_{s=1}^{m-1}f(r^{-s})
\left[\frac12\log\left(\sin\frac{s\pi}{m}\right)^2-\frac{i\pi s}{m}\right].
\tag{9}
\end{equation}

From (6), however, follows
\[
f(r^{-s})=\sum_t\chi(t)r^{-st},
\]
and the sum taken over $t$ is not changed if $t$ is replaced by $m-t$. Hence
\[
f(r^{-s})=\chi(-1)f(r^s).
\]
The factor $\chi(-1)$ has the value $+1$ or $-1$, since its square is $=+1$. We therefore now distinguish two kinds of characters $\chi_1(n),\chi_2(n)$, so that
\begin{equation}
\chi_1(-1)=1,\qquad \chi_2(-1)=-1
\tag{10}
\end{equation}

and accordingly the functions $f$ are also to be distinguished into $f_1$ and $f_2$, so that
\begin{equation}
f_1(r^{-s})=f_1(r^s),\qquad f_2(r^{-s})=-f_2(r^s).
\tag{11}
\end{equation}
Likewise we distinguish the expressions (8), (9) as $X_1$ and $X_2$, and, adding the two expressions, we get by (11)
\begin{equation}
\begin{aligned}
X_1&=-\frac1{2m}\sum_{s=1}^{m-1}f_1(r^s)\log\left(\sin\frac{s\pi}{m}\right)^2,\\
X_2&=-\frac{i\pi}{m^2}\sum_{s=1}^{m-1}s\,f_2(r^s).
\end{aligned}
\tag{12}
\end{equation}

\clearpage
\section*{\S\,195. On the class number in the real field contained in $\Omega_m$.}

Our general principles can also be applied in order to determine the class numbers in the divisors of the field $\Omega_m$.

In particular, the formulae (6), (7), \S~192 are applicable when we know the degrees of the prime ideals in such a divisor. The question concerning the ratio of the class numbers in the divisors of a field $\Omega_m$ to the class number in $\Omega_m$ itself is of great interest, and shall here be treated for the case in which the matter concerns the real field $H_m$ contained in $\Omega_m$, which we have already investigated in \S~176 ff.\footnote{Compare Kummer, ``Bestimmung der Anzahl u.s.f.'', Crelle's Journal, vol. 40 (1849).} It has there resulted that $q$ is, in the field $H_m$, the $\frac12\mu$-th power of a prime number of first degree, that the other prime numbers $p$ decompose into $e_1$ prime factors of $f_1$-th degree, so that
\begin{equation}
f_1e_1=\frac12\mu,
\tag{1}
\end{equation}
and that two kinds of prime numbers $p_1,p_2$ are to be distinguished, so that for the prime numbers of the first kind $f_1=\frac12 f$, $e_1=e$, while for those of the second kind $f_1=f$, $e_1=\frac12 e$.

For the field $H_m$ we accordingly obtain from \S~192, (6), (7) the class number $h_1$ from the formula
\begin{equation}
g_1h_1=\lim_{s=1}(s-1)\Phi_1(s),
\tag{2}
\end{equation}
\begin{equation}
\Phi_1(s)=\frac1{1-q^{-s}}\prod_{p_1}\frac1{(1-p_1^{-\frac12sf})^e}\prod_{p_2}\frac1{(1-p_2^{-sf})^{\frac12e}},
\tag{3}
\end{equation}
where the two product signs range over all prime numbers $p_1$ and $p_2$. Using the symbols $f_1,e_1$ explained above, one may also write this as
\begin{equation}
\Phi_1(s)=\frac1{1-q^{-s}}\prod_p\frac1{(1-p^{-sf_1})^{e_1}},
\tag{4}
\end{equation}
and extend the product sign over all prime numbers $p$.

The meaning of $g_1$ follows from \S~190, (11).

We shall now single out from the group $C$ of characters $\chi$ a divisor $C_1$, which is to consist of all those characters $\chi_1$ for which
\begin{equation}
\chi_1(-1)=+1.
\tag{5}
\end{equation}
It is first clear that these characters $\chi_1$ form a group $C_1$. This group is not identical with $C$, however, because there certainly is a character $\chi_0$ for which $\chi_0(-1)=-1$. To obtain such a one, we need only put, in \S~193, (6), a primitive $\mu$-th root of unity for $\Theta$, or put $\eps=-1$ in \S~193, (8). Then the totality of the characters $\chi_0\chi_1$, which do not belong to $C_1$, forms a coset $C_2$, whose elements may be denoted by $\chi_2$. Every character $\chi$ is therefore contained either in $C_1$ or in $C_2$; for if $\chi$ does not occur in $C_1$, then $\chi_0^{-1}\chi$ does occur in it, and consequently $\chi$ occurs in $C_2$. Hence $C_1$ is a divisor of $C$ of index $2$.

The characters $\chi_1$ may be formed in the following way.

1) If $m$ is odd, then, as in \S~193, we put
\[
n\equiv c^\gamma\pmod m,
\]
and in
\begin{equation}
\chi_1(n)=\Theta^\gamma
\tag{6}
\end{equation}
let $\Theta$ run through the totality of the $\frac12\mu$-th roots of unity.

2) If $m$ is even, then we put
\[
n\equiv (-1)^\alpha5^\beta
\]
and obtain
\begin{equation}
\chi_1(n)=\Theta^\beta,
\tag{7}
\end{equation}
where $\Theta$ again runs through all the $\frac12\mu$-th roots of unity.

Both statements result easily from the fact that for $n=-1$ one has, in the first case, $\gamma=\frac12\mu$, and in the second, $\alpha=1$, $\beta=0$, so that both times $\chi_1(-1)=1$, and from the fact that each of the formulae (6), (7) represents exactly $\frac12\mu$ different characters.

What is now essential is the following:
\begin{center}
\textbf{If $p$ is any prime number different from $q$, then in the form $\chi_1(p)$, as $\chi_1$ runs through the group $C_1$, one obtains every $f_1$-th root of unity, and every one equally often, namely $e_1$ times.}
\end{center}
That all $\chi_1(p)$ are roots of unity of degree $f_1$ is clear, since always (compare \S~177)
\[
p^{f_1}\equiv \pm 1\pmod m,
\]
and therefore
\[
[\chi_1(p)]^{f_1}=\chi_1(\pm1)=1.
\]
If however among the $\chi_1(p)$ a primitive $f_1$-th root of unity occurs, then all other $f_1$-th roots of unity also occur among them, and equally often, namely as often as the root of unity $1$.

This follows immediately from the group nature of the $\chi_1$. Thus it remains only to show that among the $\chi_1(p)$ a primitive $f_1$-th root of unity occurs.

This possibility follows very simply from the determination of $\chi_1$. First suppose that, in (6), with $m$ odd, $\Theta$ is a primitive root of unity of degree $\frac12\mu$, and that $\gamma$ is the index of $p$. Then a power of $\Theta^\gamma$, say $\Theta^{\gamma\lambda}$, is equal to $1$ if and only if
\[
2\gamma\lambda\equiv0\pmod\mu.
\]
Now $\mu=ef$, and $e$ is the greatest common divisor of $\gamma$ and $\mu$; therefore this condition becomes $2\lambda\equiv0\pmod f$. But by \S~177, 3, $f_1=f$ or $f_1=\frac12 f$ according as $f$ is odd or even, and consequently one must have $\lambda\equiv0\pmod {f_1}$. Thus $\Theta^\gamma$ is a primitive $f_1$-th root of unity.

If, however, $m$ is even, we likewise take in (7) a primitive $\frac12\mu$-th root of unity for $\Theta$, and put
\[
p\equiv(-1)^\alpha5^\beta\pmod m,
\]
so that $f$ is the smallest positive number satisfying the congruences
\[
\alpha f\equiv0\pmod2,
\qquad
\beta f\equiv0\pmod{\tfrac12\mu}.
\]
Now $\Theta^{\beta\lambda}$ is equal to $1$ only when $\beta\lambda\equiv0\pmod{\frac12\mu}$. If $\lambda\ne1$, so that $\beta\not\equiv0$, then the smallest value that $\lambda$ can have is even, and hence this smallest value of $\lambda$ is $f$.

But if $\lambda=1$, then $\beta$ must be $0$, and two further cases are possible:
\[
\begin{array}{ll}
\text{either} & f=1,\quad \alpha=0,\quad p\equiv1\pmod m,\quad \lambda=f,\\[2mm]
\text{or} & f=2,\quad \alpha=1,\quad p\equiv-1\pmod m,\quad \lambda=\frac12f.
\end{array}
\]
Only in this last case does $p$ belong to the first kind, and it follows therefore that here also, in general, $\Theta^\beta$ is a primitive $f_1$-th root of unity. The theorem is thereby proved.

This theorem permits us to apply formula \S~193, (9), which here gives
\begin{equation}
(1-p^{-sf_1})^{e_1}=\prod_{\chi_1}[1-\chi_1(p)p^{-s}],
\tag{8}
\end{equation}
where the product is extended over all characters $\chi_1$ of the group $C_1$. Now the function $\Phi_1(s)$ may be developed further in the same way as $\Phi(s)$ in \S~193. One obtains, on the same path as there, the formula (11),
\begin{equation}
\Phi_1(s)=\frac1{1-q^{-s}}\prod_{\chi_1}\sum_{n=1}^{\infty}\frac{\chi_1(n)}{n^s},
\tag{9}
\end{equation}
and when one passes to the limit $s=1$,
\begin{equation}
g_1h_1=\prod_{\chi_1}\sum_{n=1}^{\infty}\frac{\chi_1(n)}{n},
\tag{10}
\end{equation}
where now, however, the product is to be extended over all characters of the group $C_1$ with the exception of the principal character. This product is therefore a part of the product by which the number $gh$ is expressed in \S~193.

The numbers $g,g_1$ result from the general expression given in \S~190, (11):
\begin{equation}
\frac{2^\nu\pi^{n-\nu}L}{w\sqrt{\pm\Delta}}.
\tag{11}
\end{equation}
Since, as was shown in \S~178, the units in $\Omega_m$ and $H_m$, apart from the roots of unity, are the same, a fundamental system of units in $H_m$ is at the same time a fundamental system in $\Omega_m$.

Let therefore $\eps_1,\eps_2,\ldots,\eps_{\nu-1}$ denote a fundamental system of real units, and let $\eps_{i,1},\eps_{i,2},\ldots,\eps_{i,\nu}$ be the conjugate units to $\eps_i$. Then the conjugate logarithms in the field $H_m$ are
\[
\log|\eps_{i,1}|,\quad \log|\eps_{i,2}|,\quad\ldots,\quad \log|\eps_{i,\nu}|,
\]
whereas in the field $\Omega_m$, in which only conjugate-imaginary pairs occur (\S~185), they are
\[
2\log|\eps_{i,1}|,\quad 2\log|\eps_{i,2}|,\quad\ldots,\quad2\log|\eps_{i,\nu}|.
\]

If therefore we denote by
\begin{equation}
E=\pm\begin{vmatrix}
\log|\eps_{1,1}|&\log|\eps_{1,2}|&\cdots&\log|\eps_{1,\nu-1}|\\
\log|\eps_{2,1}|&\log|\eps_{2,2}|&\cdots&\log|\eps_{2,\nu-1}|\\
\cdots&\cdots&\cdots&\cdots\\
\log|\eps_{\nu-1,1}|&\log|\eps_{\nu-1,2}|&\cdots&\log|\eps_{\nu-1,\nu-1}|
\end{vmatrix}
\tag{12}
\end{equation}
the regulator of the field $H_m$ (\S~187), then, in applying the expression (11), one must put in the field $\Omega_m$
\[
L=2^{\nu-1}E,
\]
and in the field $H_m$
\[
L=E.
\]
Further,
\[
\begin{array}{lll}
\text{in }\Omega_m: & n=\mu, & \nu=\frac12\mu,\\
\text{in }H_m: & n=\frac12\mu, & \nu=\frac12\mu.
\end{array}
\]
In the field $H_m$ only the two roots of unity $\pm1$ are contained. In $\Omega_m$ the positive and negative powers of $r$ are contained, but no other roots of unity (\S~175), and in the enumeration it is still to be noted that, when $m$ is even, $-1$ is among the powers of $r$, while when $m$ is odd it is not. Hence we have
\[
\begin{array}{lll}
\text{in }\Omega_m: & w=2m, & m\text{ odd},\\
& w=m, & m\text{ even},\\
\text{in }H_m: & w=2.&
\end{array}
\]
Finally, if $\Delta,\Delta_1$ are the fundamental numbers of the fields $\Omega_m,H_m$ (\S~176), we still have
\[
\pm\Delta=q\Delta_1^2\quad(m\text{ odd}),
\qquad
\pm\Delta=4\Delta_1^2\quad(m\text{ even}),
\]
and it follows, if for simplicity we put $\nu$ for $\frac12\mu$, that
\begin{align}
g&=\frac{2^{2(\nu-1)}\pi^\nu E}{m\sqrt q\,\Delta_1}, &&m\text{ odd},\notag\\
g&=\frac{2^{2(\nu-1)}\pi^\nu E}{m\Delta_1}, &&m\text{ even},\notag\\
g_1&=\frac{2^{\nu-1}E}{\sqrt{\Delta_1}},\tag{13}
\end{align}
whence
\begin{align}
g&=\frac{2^{\nu-1}\pi^\nu g_1}{m\sqrt{q\Delta_1}}, &&m\text{ odd},\notag\\
g&=\frac{2^{\nu-1}\pi^\nu g_1}{m\sqrt{\Delta_1}}, &&m\text{ even}.\tag{14}
\end{align}

If we combine formula (10) with formula \S~193, (15), and use the notation introduced in \S~194,
\[
X_1=\sum\frac{\chi_1(n)}{n},\qquad X_2=\sum\frac{\chi_2(n)}{n},
\]
then we obtain
\[
gh=g_1h_1\prod X_2.
\]
Replacing $g$ by its value (14), we obtain for the class number
\begin{equation}
h=kh_1,
\tag{15}
\end{equation}
if
\begin{equation}
\begin{aligned}
k&=\frac{m\sqrt{q\Delta_1}}{2^{\nu-1}\pi^\nu}\prod X_2, &&m\text{ odd},\\
&=\frac{m\sqrt{\Delta_1}}{2^{\nu-1}\pi^\nu}\prod X_2, &&m\text{ even},
\end{aligned}
\tag{16}
\end{equation}
and where [by (10), (13)]
\begin{equation}
h_1=\frac{\sqrt{\Delta_1}}{2^{\nu-1}E}\prod X_1
\tag{17}
\end{equation}
is the class number of the field $H_m$. In the latter expression the product extends over all characters $\chi_1$ of the group $C_1$ with the exception of the principal character.

That $k$ is an integer, and therefore that $h$ is a multiple of $h_1$, will soon be shown. The numbers $k$ and $h_1$ are called the first and second factor of the class number.

In the case in which
\[
m=2^\varkappa
\]
is a power of $2$, we have by \S~176, (7),
\[
\Delta_1=2^{(\varkappa-1)2^{\varkappa-2}-1},
\]
and if we insert this value, then in this case we may write the preceding formulae as
\begin{equation}
\begin{aligned}
k&=\frac{\sqrt2\,2^{(\varkappa-3)2^{\varkappa-3}}2^\varkappa}{\pi^\nu}\prod X_2,\\
h_1&=\frac{\sqrt2\,2^{(\varkappa-3)2^{\varkappa-3}}}{E}\prod X_1.
\end{aligned}
\tag{18}
\end{equation}

\section*{\S\,196. Class number in the field of the eighth roots of unity.}

We shall first treat the case $m=8$, $\mu=4$, $\nu=2$, partly to illustrate the preceding results and partly to gain a basis for the complete induction to be used later.

Here, apart from the principal character, we have only three characters, namely, if
\[
n\equiv(-1)^\alpha5^\beta\pmod8,
\]
then
\[
\chi_1(n)=(-1)^\beta,
\qquad
\chi_2(n)=(-1)^\alpha,
\qquad
\chi_3(n)=(-1)^{\alpha+\beta}.
\]
For $n=-1$ one has $\alpha=1$, $\beta=0$, and consequently $\chi_1$ belongs to the first kind, $\chi_2,\chi_3$ to the second kind [\S~195, (5)]. For the four numbers $n=1,3,5,7$ the following values of these characters result:
\[
\begin{array}{c|rrrr}
&1&3&5&7\\\hline
\chi_1&+1&-1&-1&+1\\
\chi_2&+1&-1&+1&-1\\
\chi_3&+1&+1&-1&-1.
\end{array}
\]
From this one obtains the three functions $f_1,f_2,f_3$ [\S~194, (2)]:
\[
\begin{aligned}
f_1(x)&=x-x^3-x^5+x^7=x(1-x^2)(1-x^4),\\
f_2(x)&=x-x^3+x^5-x^7=x(1-x^2)(1+x^4),\\
f_3(x)&=x+x^3-x^5-x^7=x(1+x^2)(1-x^4).
\end{aligned}
\]
For $r$ we may put
\[
r=\frac{1+i}{\sqrt2},\quad r^2=i,
\quad r^3=-\frac{1-i}{\sqrt2},\quad r^4=-1,
\]
and from this it follows that
\[
\begin{array}{lll}
f_1(r)=2\sqrt2,& f_2(r)=0,& f_3(r)=2\sqrt2\,i,\\
f_1(r^2)=0,& f_2(r^2)=4i,& f_3(r^2)=0,\\
f_1(r^3)=-2\sqrt2,& f_2(r^3)=0,& f_3(r^3)=2\sqrt2\,i.
\end{array}
\]
For $x=r^4$ all three functions $f(x)$ vanish, and for $x=r^5,r^6,r^7$ the function $f_1(x)$ receives the same values, while $f_2(x)$ and $f_3(x)$ receive the opposite values, as for $x=r^3,r^2,r$.

Consequently, from \S~194, (12), with regard to the trigonometric formula
\[
\sin\frac{3\pi}{8}=\cos\frac\pi8,
\]
one obtains
\[
X_1=-\frac1{\sqrt2}\log\tan\frac\pi8,
\qquad
X_2=-\frac\pi4,
\qquad
X_3=-\frac\pi{2\sqrt2}.
\]
Now
\[
\tau=\frac{r-1}{r^2(r+1)}=\tan\frac\pi8=\sqrt2-1
\]
is a unit of the field $H_8$, which here is nothing other than the quadratic field arising from $\sqrt2$, and the units in it are therefore obtained from the solutions of Pell's equation
\[
u^2-2v^2=\pm1
\]
in the form $u+v\sqrt2$.

The smallest positive solution of this Pell equation is however $u=1$, $v=1$, and consequently, by Vol. I, \S~128, all units in $H_8$ are found in the form $\pm(1+\sqrt2)^a$, where $a$ is a positive or negative integral exponent. The unit $\tau$ results from this for $a=-1$, and hence all units are also contained in the form
\[
\pm\tau^a.
\]
Thus $\tau$ is a fundamental unit, and the determinant $E$ occurring in the formulae of the preceding paragraph [\S~195, (12)] here becomes, since $\tau$ is a proper fraction,
\[
E=-\log\tau.
\]
From \S~195, (15), (18), one therefore obtains
\[
h_1=1,\\ k=1,
\qquad h=1.
\]
The field $\Omega_8$ is therefore a one-class field, that is, one in which the decomposition of the integers into prime factors is possible according to the same laws as in the field of rational numbers.

At the same time it has also resulted that all units in the field $H_8$ are representable in the form $\pm\tau^a$. The conjugate values of the unit $\tau$ are however
\[
\tau_1=\tan\frac\pi8,\qquad
\tau_3=-\tan\frac{3\pi}{8}=-\tau_1^{-1},
\]
and therefore have different signs. We conclude from this for this case the theorem:
\begin{center}
\textbf{A unit in $H_8$ which, together with its conjugates, is of the same sign is, apart from the sign, the square of a unit.}
\end{center}

\section*{\S\,197. Recurrent calculation of the class number in the field $\Omega_m$, when $m$ is a power of $2$.}

For the general determination of the class number in the field $\Omega_m$, under the assumption that $m$ is any power of $2$ and greater than $8$, we apply a recurrent procedure. Thus let
\begin{equation}
m=2^\varkappa,
\qquad
\mu=2^{\varkappa-1},
\qquad
\nu=2^{\varkappa-2},
\qquad
\nu'=2^{\varkappa-3}.
\tag{1}
\end{equation}
Beside the field $\Omega_m$ we consider the field $\Omega_\mu$, which is a divisor of $\Omega_m$ and which we shall here also denote by $\Omega'_m$. Let $h'$ be the class number in the field $\Omega'_m$, that is, the number obtained from $h$ when $\varkappa$ is replaced by $\varkappa-1$, and put
\begin{equation}
h=h'H.
\tag{2}
\end{equation}
If we presuppose $h'$ as already known, then all that remains is the calculation of $H$, which, as will be shown, is an integer. Decomposing $h$ and $h'$ as in \S~195, we put
\begin{equation}
h=kh_1,
\qquad
h'=k'h'_1,
\qquad
H=AB,
\tag{3}
\end{equation}
\begin{equation}
k=k'A,
\qquad
h_1=h'_1B,
\tag{4}
\end{equation}
where $k'$ is obtained from $k$, and $h'_1$ from $h_1$, by replacing $\varkappa$ by $\varkappa-1$. In general, we shall now indicate by accents on letters that $\varkappa-1$ is to be put in place of $\varkappa$.

By \S~176, (7), one then has
\[
\Delta_1=2^{(\varkappa-1)2^{\varkappa-2}-1},
\qquad
\Delta'_1=2^{(\varkappa-2)2^{\varkappa-3}-1},
\]
and consequently
\begin{equation}
\Delta_1=2^{\varkappa2^{\varkappa-3}}\Delta'_1.
\tag{5}
\end{equation}
If $E$ has the meaning of \S~195, (12), and $E'$ is obtained from $E$ when $\varkappa$ is replaced by $\varkappa-1$, then we put further
\begin{equation}
E=E'D,
\tag{6}
\end{equation}

whereby $D$ is defined as a non-zero positive number. As regards the sums
\[
X=\sum\frac{\chi(n)}{n},
\]
one must recall that $\chi(n)$ [by \S~193, (8)] is defined by
\begin{equation}
n\equiv(-1)^\alpha5^\beta\pmod m,
\qquad
\chi(n)=(\pm1)^\alpha\Theta^\beta,
\tag{7}
\end{equation}
where $\Theta$ denotes a $\nu$-th root of unity.

Among the $\nu$-th roots of unity the $\nu'$-th roots of unity are also contained, and among the characters $\chi$ also the characters for the field $\Omega'_m$. We distinguish, accordingly, primitive and imprimitive characters of the field $\Omega_m$ by designating as primitive the characters $\chi$ proper to the field $\Omega_m$, that is, those in which $\Theta$ is a primitive $\nu$-th root of unity.

Among the sums $X$ all sums belonging to the field $\Omega'_m$ also occur:
\[
X'=\sum\frac{\chi'(n)}{n}.
\]
If now in (4) one inserts for $k,k'$ and $h_1,h'_1$ the expressions formed in \S~195, (16), (17), and then cancels the sums $X'$, one obtains
\begin{equation}
\begin{aligned}
\pi^{\nu'}A&=2\cdot2^{2^{\varkappa-4}(\varkappa-2)}\prod X_2,\\
DB&=2^{2^{\varkappa-4}(\varkappa-2)}\prod X_1,
\end{aligned}
\tag{8}
\end{equation}
where now the products $\prod$ extend only over the sums $X_1,X_2$ formed with the primitive characters $\chi_1,\chi_2$.

In the expressions for the $X$ given in \S~194, (12),
\begin{equation}
\begin{aligned}
X_1&=-\frac1{2m}\sum_{s=1}^{m-1} f_1(r^s)\log\left(\sin\frac{s\pi}{m}\right)^2,\\
X_2&=-\frac{i\pi}{m^2}\sum_{s=1}^{m-1} s f_2(r^s),
\end{aligned}
\tag{9}
\end{equation}
considerable simplifications now enter under the assumption that $m$ is a power of $2$ and that $\chi_1,\chi_2$ are primitive characters.

The number $1+\mu$ has the indices $\alpha=0$, $\beta=\nu'$, and consequently by (7)
\[
\chi(1+\mu)=\Theta^{\nu'}=-1,
\]
when $\chi$ is a primitive character. Further, for an odd $n$,
\[
n(1+\mu)\equiv n+\mu\pmod m,
\]
and consequently
\[
\chi(n+\mu)=-\chi(n).
\]
If we now consider the function $f(r^s)$,
\begin{equation}
f(r^s)=\sum_n\chi(n)r^{sn},
\tag{10}
\end{equation}
where $n$ runs through the odd numbers of a complete residue system for the modulus $m$, then, since in (10) we may replace $n$ under the summation sign by $n+\mu$, we obtain
\[
f(r^s)=-r^{\mu s}\sum\chi(n)r^{sn},
\]
and since $r^\mu=-1$,
\[
f(r^s)=-(-1)^s f(r^s).
\]
Thus
\begin{equation}
f(r^s)=0,\qquad\text{if }s\text{ is even},
\tag{11}
\end{equation}
and hence in formulae (9) we may restrict the summation letter $s$ to the odd numbers which are smaller than $m$.

But from (10) one has
\[
f(r^s)=\chi(s)^{-1}\sum_n\chi(ns)r^{sn},
\]
and, replacing $sn$ by $n$, which is permitted for odd $s$,
\begin{equation}
f(r^s)=\chi(s)^{-1}f(r),\qquad\text{if }s\text{ is odd}.
\tag{12}
\end{equation}
The function $f(r)$, if $\alpha,\beta$ are the indices of $n$, has the expression
\begin{equation}
f(r)=\sum(\pm1)^\alpha\Theta^\beta r^n,
\tag{13}
\end{equation}
and is therefore equivalent to the cyclotomic resolvent considered in \S~17 of this volume,
\[
(\pm1,\Theta,r),
\]
for which in \S~17, (17), the relation was established:
\begin{equation}
(\pm1,\Theta,r)(\pm1,\Theta^{-1},r)=\pm m.
\tag{14}
\end{equation}
Here throughout the upper signs apply to the sums $X_1$, the lower to the sums $X_2$ [\S~195, (7)].

We now divide the series of numbers $s$ occurring in (9) into four parts. Let $t$ be a symbol which runs through the
series of positive odd numbers from $1$ to $\nu-1$; then the odd $s$ runs through the four number series
\[
t,
\quad t+\mu,
\quad m-t,
\quad m-t-\mu,
\]
and one has
\[
\begin{gathered}
\chi(t+\mu)=-\chi(t),\\
\chi_1(m-t)=\chi_1(t),\qquad \chi_1(m-\mu-t)=-\chi_1(t),\\
\chi_2(m-t)=-\chi_2(t),\qquad \chi_2(m-\mu-t)=\chi_2(t).
\end{gathered}
\]
From (12) it follows that
\[
\begin{gathered}
f(r^{t+\mu})=-f(r^t),\\
f_1(r^{m-t})=f_1(r^t),\qquad f_1(r^{m-\mu-t})=-f_1(r^t),\\
f_2(r^{m-t})=-f_2(r^t),\qquad f_2(r^{m-\mu-t})=f_2(r^t).
\end{gathered}
\]
Dividing the sums (9) accordingly into four parts each and also using the trigonometric equation
\[
\sin\frac{(t+\mu)\pi}{m}=\cos\frac{t\pi}{m},
\]
one obtains
\[
X_1=-\frac1m\sum_t f_1(r^t)\log\left(\tan\frac{t\pi}{m}\right)^2,
\qquad
X_2=\frac{i\pi}{m}\sum_t f_2(r^t).
\]
In this sum one has
\[
\frac{t\pi}{m}<\frac\pi4,
\qquad
\tan\frac{t\pi}{m}>0,
\]
and consequently, with regard to (12), (13), if we put $\Theta^{-1}$ for $\Theta$, whereby $\chi(t)^{-1}$ passes into $\chi(t)$, we obtain
\begin{equation}
X_1=-\frac2m(+1,\Theta^{-1},r)
\sum_{1\le t\le\nu-1}\chi_1(t)\log\tan\frac{t\pi}{m},
\tag{15}
\end{equation}
\begin{equation}
X_2=\frac{i\pi}{m}(-1,\Theta^{-1},r)
\sum_{1\le t\le\nu-1}\chi_2(t).
\tag{16}
\end{equation}

\section*{\S\,198. The class-number factor $A$.}

To calculate first $A$, we have to insert the expression (16) for $X_2$ into the first formula (8) of the preceding paragraph. This formula can also be represented as
\begin{equation}
\pi^{\nu'}A=2m^{1/4}2^{-\nu'}\prod X_2.
\tag{1}
\end{equation}

If we denote by $\alpha,\beta$ the indices of $t$, then $\chi_2(t)=(-1)^\alpha\Theta^\beta$;
\begin{equation}
\sum_t\chi_2(t)=\sum(-1)^\alpha\Theta^\beta=\varphi(\Theta)
\tag{2}
\end{equation}
is an integer of the field $\Omega_\nu$, and it becomes
\[
X_2=\frac{i\pi}{m}(-1,\Theta^{-1},r)\varphi(\Theta).
\]
Now for $\Theta$ all roots of the equation
\begin{equation}
\Theta^{\nu'}+1=0
\tag{3}
\end{equation}
are to be put, and for $A$ one obtains, with regard to \S~197, (14),
\begin{equation}
A=2^{1-\nu'}\prod_\Theta\varphi(\Theta).
\tag{4}
\end{equation}
The sum by which the number $\varphi(\Theta)$ is defined in (2) contains $\nu'$ terms of the form $(-1)^\alpha\Theta^\beta$, where $\alpha,\beta$ are to be determined so that
\begin{equation}
(-1)^\alpha5^\beta\equiv t\pmod m
\tag{5}
\end{equation}
and $t$ lies between $0$ and $\nu$. First it is clear that among the exponents $\beta$ no two can occur which are congruent modulo $\nu'$. For if $\beta\equiv\beta'\pmod{\nu'}$, then $5^\beta\equiv5^{\beta'}\pmod\mu$, and from
\[
(-1)^\alpha5^\beta\equiv t,
\qquad
(-1)^{\alpha'}5^{\beta'}\equiv t'\pmod m
\]
it follows that
\[
(-1)^\alpha t-(-1)^{\alpha'}t'\equiv0\pmod\mu.
\]
Since however $t$ and $t'$ are absolutely smaller than $\nu=\frac12\mu$, this is possible only when $t=t'$, $\alpha=\alpha'$. Thus $\beta$ in (2) runs through a complete residue system modulo $\nu'$.

But since $\beta$ is to be taken modulo $\nu$, the $\beta$ determined by (5) will come out partly smaller and partly greater than $\nu'$. If, then, we let $\beta_1$ run through the series of numbers $0,1,\ldots,\nu'-1$, then we obtain in (2) two kinds of terms:
\begin{equation}
\begin{array}{ll}
1. & (-1)^\alpha\Theta^\beta=(-1)^\alpha\Theta^{\beta_1},\quad \beta=\beta_1<\nu',\\[1mm]
2. & (-1)^\alpha\Theta^\beta= -(-1)^\alpha\Theta^{\beta_1},\quad \beta=\beta_1+\nu'\geqq\nu'.
\end{array}
\tag{6}
\end{equation}
Now by (5) the absolutely smallest residue of $(-1)^\alpha5^\beta$ modulo $m$ is positive; therefore if we determine a number $c_\beta=\pm1$ so that the absolutely smallest residue of $c_\beta5^\beta$ is positive, then in case (6), 1,
\[
c_{\beta_1}=(-1)^\alpha.
\]

Further, $5^{\nu'}\equiv1\pmod\mu$, and since $5^{\nu'}$ is not also congruent to $1$ modulo $m$, it follows that
\begin{equation}
5^{\nu'}\equiv1+\mu\pmod m.
\tag{7}
\end{equation}
Consequently, in case (6), 2,
\[
(-1)^\alpha5^\beta=(-1)^\alpha5^{\beta_1}5^{\nu'}\equiv(-1)^\alpha5^{\beta_1}+\mu\pmod m,
\]
and hence here the absolutely smallest residue of $(-1)^\alpha5^{\beta_1}$ modulo $m$ is negative. Thus in case (6), 2,
\[
c_{\beta_1}=-(-1)^\alpha
\]
is to be put.

From this, by (2), it follows, if we again write $\beta$ for $\beta_1$,
\begin{equation}
\varphi(\Theta)=\sum_{\beta=0}^{\nu'-1}c_\beta\Theta^\beta
=c_0+c_1\Theta+c_2\Theta^2+\cdots+c_{\nu'-1}\Theta^{\nu'-1},
\tag{8}
\end{equation}
and here
\begin{equation}
c_\beta=+1\quad\text{or}\quad=-1,
\tag{9}
\end{equation}
according as the absolutely smallest residue of $5^\beta$ is positive or negative. Thus $\varphi(\Theta)$ is an integer of the field $\Omega_\nu$.

It remains to determine the behaviour of this number toward the number $2$, or rather toward the prime factor $(1-\Theta)$ of $2$ contained in $\Omega_\nu$. For this purpose we form
\begin{equation}
(1-\Theta)\varphi(\Theta)=2\psi(\Theta),
\tag{10}
\end{equation}
and this gives
\begin{equation}
\psi(\Theta)=\frac{c_0+c_{\nu'-1}}{2}+\frac{c_1-c_0}{2}\Theta+\frac{c_2-c_1}{2}\Theta^2+\cdots+
\frac{c_{\nu'-1}-c_{\nu'-2}}{2}\Theta^{\nu'-1}.
\tag{11}
\end{equation}
The coefficients in this expression are all $0$ or $1$, and hence $\psi(\Theta)$ also is an integer of the field $\Omega_\nu$. But for this number one obtains, putting $\Theta\equiv1$,
\begin{equation}
\psi(\Theta)\equiv c_{\nu'-1}\equiv\pm1\pmod{(1-\Theta)},
\tag{12}
\end{equation}
and consequently $\psi(\Theta)$ is relatively prime to $2$.

The product occurring in (4) is nothing other than the norm, taken with respect to the field $\Omega_\nu$, of the number $\varphi(\Theta)$; and if we denote this norm by $N$, then by \S~169
\[
N(1-\Theta)=2.
\]

Thus by (10)
\[
N\varphi(\Theta)=2^{\nu'-1}N\psi(\Theta).
\]
It follows that
\begin{equation}
A=N\psi(\Theta),
\tag{13}
\end{equation}
and hence, with regard to (12), that $A$ is an odd integer.

If we denote by $A_\varkappa$ the number $A$ belonging to $m=2^\varkappa$, then from \S~197, (4), by complete induction, the expression for the first factor of the class number is
\begin{equation}
k=A_4A_5\cdots A_\varkappa,
\tag{14}
\end{equation}
and from this the theorem follows:
\begin{center}
\textbf{1. The first factor of the class number is an odd integer.}
\end{center}
The number $A_\varkappa$ is not difficult to calculate from formula (10) for the first values of $\varkappa$.

For $m=16$ one finds $\nu'=2$, $\beta=0,1$, $c_0=c_1=1$, consequently
\[
\psi(\Theta)=1,
\qquad
A_4=1.
\]
For $m=32$, $\nu'=4$, $\beta=0,1,2,3$,
\[
c_0=1,
\quad c_1=1,
\quad c_2=-1,
\quad c_3=-1,
\]
\[
\psi(\Theta)=-\Theta^2,
\]
so that also $A_5=1$.

For $m=64$, $\nu'=8$, $\beta=0,1,2,3,4,5,6,7$, the absolutely smallest residues of $5^\beta$ are
\[
1,\;5,\;25,\;-3,\;-15,\;-11,\;9,\;-19;
\]
thus the $c_\beta$ are
\[
+1,+1,+1,-1,-1,-1,+1,-1,
\]
and
\[
\psi(\Theta)=-\Theta^3+\Theta^6-\Theta^7,
\qquad
\psi(\Theta)\psi(-\Theta)=\Theta^6(\Theta^6-2i),
\]
from which it follows easily that
\[
A_6=17.
\]
A somewhat longer calculation gives still
\[
A_7=21121.
\]

\section*{\S\,199. The class-number factor $B$.}

The class-number factor $B$ must be calculated in an entirely different way. Here, if $\Theta$ is again a primitive root of the equation \S~198, (3), namely $\Theta^{\nu'}+1=0$, and
\begin{equation}
t\equiv(-1)^\alpha5^\beta\pmod m,
\tag{1}
\end{equation}
then one is to put
\begin{equation}
\chi_1(t)=\Theta^\beta,
\tag{2}
\end{equation}
and formula \S~197, (15) gives
\begin{equation}
X_1=-\frac2m(1,\Theta^{-1},r)\sum_t\Theta^\beta\log\tan\frac{t\pi}{m}.
\tag{3}
\end{equation}
Now for every odd $t$, if [with regard to (1)]
\[
r=e^{2\pi i/m},
\qquad
r^\nu=i,
\qquad
r^{t\nu}=(-1)^\alpha i,
\]
one has
\begin{equation}
\tan\frac{t\pi}{m}=(-1)^\alpha r^t\frac{1-r^t}{1+r^t},
\tag{4}
\end{equation}
and this, as we already saw in \S~178, is a unit of the real field $H_m$. The conjugate values of this unit with respect to this field are obtained by putting $t\equiv5^\beta\pmod m$ and by letting $\beta$ run through a complete residue system modulo $\mu$. We put accordingly
\begin{equation}
\tau_\beta=\tan\frac{5^\beta\pi}{m}.
\tag{5}
\end{equation}
By \S~198, (7), $5^{\beta+\nu'}\equiv5^\beta+\mu\pmod m$, and from this, by a known trigonometric formula, there follows
\begin{equation}
\tau_{\beta+\nu'}=-\frac1{\tau_\beta}.
\tag{6}
\end{equation}
If we therefore put
\begin{equation}
\lambda_\beta=\frac12\log\tau_\beta^2,
\tag{7}
\end{equation}
then by (6)
\begin{equation}
\lambda_{\beta+\nu'}=-\lambda_\beta,
\tag{8}
\end{equation}
and therefore
\begin{equation}
\Theta^{\beta+\nu'}\lambda_{\beta+\nu'}=\Theta^\beta\lambda_\beta.
\tag{9}
\end{equation}
Among the values of $\beta$ occurring in the sum (3), as we already saw in the preceding paragraph, no two are congruent modulo $\nu'$, and if therefore by means of
formula (9) we reduce all values of $\beta$ to the interval from $0$ to $\nu'-1$, then we obtain all numbers of this interval.

Accordingly formula (3) may be represented as
\begin{equation}
X_1=-\frac2m(1,\Theta^{-1},r)\sum_{\beta=0}^{\nu'-1}\Theta^\beta\lambda_\beta.
\tag{10}
\end{equation}
This must be inserted in the expression given in \S~197, (8),
\[
DB=2^{2^{\varkappa-4}(\varkappa-2)}\prod X_1=m^{\nu'}2^{-\nu'}\prod X_1,
\]
and then, if we put
\begin{equation}
U_\Theta=\sum_{\beta=0}^{\nu'-1}\Theta^\beta\lambda_\beta,
\tag{11}
\end{equation}
we obtain, with regard to \S~197, (14),
\begin{equation}
DB=\prod_\Theta U_\Theta.
\tag{12}
\end{equation}
This product may now be represented in the following way by a determinant which contains only the conjugate logarithms of the units $\tau$.

If we omit the index $\Theta$ in $U$, we may set up the system of equations
\[
\begin{aligned}
U&=\lambda_0+\lambda_1\Theta+\lambda_2\Theta^2+\cdots+\lambda_{\nu'-1}\Theta^{\nu'-1},\\
\Theta U&=-\lambda_{\nu'-1}+\lambda_0\Theta+\lambda_1\Theta^2+\cdots+\lambda_{\nu'-2}\Theta^{\nu'-1},\\
&\hspace{2.5cm}\cdots\\
\Theta^{\nu'-1}U&=-\lambda_1-\lambda_2\Theta-\lambda_3\Theta^2+\cdots+\lambda_0\Theta^{\nu'-1},
\end{aligned}
\]
from which, by elimination of the powers of $\Theta$, it follows that
\begin{equation}
\begin{vmatrix}
\lambda_0-U&\lambda_1&\lambda_2&\cdots&\lambda_{\nu'-1}\\
-\lambda_{\nu'-1}&\lambda_0-U&\lambda_1&\cdots&\lambda_{\nu'-2}\\
\cdots&\cdots&\cdots&\cdots&\cdots\\
-\lambda_1&-\lambda_2&-\lambda_3&\cdots&\lambda_0-U
\end{vmatrix}=0.
\tag{13}
\end{equation}
This is an equation of degree $\nu'$ for $U$, whose roots are the quantities $U_\Theta$. If one orders it according to powers of $U$, then, since $\nu'$ is even, the highest power $U^{\nu'}$ of $U$ receives the coefficient $1$, and one therefore obtains the product of the roots occurring in (12), when $U=0$ is put, from the determinant (13).

We shall arrange the determinant so obtained somewhat differently, leaving the first row in its place and writing the
remaining rows in the reverse order and with changed signs. In doing this $\frac12\nu'$ changes of sign are necessary, and therefore for the product (12) we obtain the determinant defined by the following equation, whose absolute value we denote by $T$:
\begin{equation}
(-1)^{\frac12\nu'}
\begin{vmatrix}
\lambda_0&\lambda_1&\cdots&\lambda_{\nu'-2}&\lambda_{\nu'-1}\\
\lambda_1&\lambda_2&\cdots&\lambda_{\nu'-1}&-\lambda_0\\
\lambda_2&\lambda_3&\cdots&-\lambda_0&-\lambda_1\\
\cdots&\cdots&\cdots&\cdots&\cdots\\
\lambda_{\nu'-1}&-\lambda_0&\cdots&-\lambda_{\nu'-3}&-\lambda_{\nu'-2}
\end{vmatrix}=\pm T,
\tag{14}
\end{equation}
where to the right of the diagonal row running from left below to right above the negative signs stand.

The determinant $T$ is here ordered in such a way that each row arises from the preceding one by the substitution $(r,r^{5^\beta})$. Each column therefore contains a system of conjugate logarithms.

By (12) one now has
\begin{equation}
B=\frac{T}{D}.
\tag{15}
\end{equation}
To obtain clarity concerning the nature of this result, it is necessary to consider somewhat more closely the number $D$ which was defined by \S~195, (12), \S~197, (6); and this requires entering into the theory of the units of the field $H_m$.

\clearpage
\section*{\S\,200. Normal units in $H_m$.}

We now introduce a simplifying notation, by putting
\begin{equation}
r^{5^j}=r_j,
\tag{1}
\end{equation}
and letting $\beta$ run through a complete system of residues modulo $\nu$. From this it follows that
\begin{equation}
r_{\beta+\nu'}=-r_\beta,
\tag{2}
\end{equation}
and in the form $(r,r_j)$ there are then indeed not all substitutions of the group of the field $\Omega_m$, but the substitutions $(r,r_j^{-1})$ must still be added; nevertheless $(r,r_j)$ supplies us with all substitutions of the group of $H_m$.

We shall now denote by $\mathfrak E(r)$ a unit of the field $H_m$. Among these there are also the units of the field $H_\mu$ (as
imprimitive numbers), and they are characterized by the equation $\mathfrak E(r)=\mathfrak E(-r)$.

By a \textbf{normal unit} of the field $H_m$ we shall now understand a unit which is not equal to $\pm1$, but which satisfies the equation
\begin{equation}
\mathfrak E(r)\mathfrak E(-r)=\pm1.
\tag{3}
\end{equation}

It is clear that a unit of the field $H_\mu$ can certainly not satisfy this condition. But not every unit in $H_m$ which does not belong to the field $H_\mu$ will be a normal unit.\footnote{In the above-mentioned work (Acta mathematica, vol. 8), the author designated the normal units as ``primitive units''. This name is rejected here, because it could give occasion to the error that every unit of the field $H_m$ which does not at the same time belong to the field $H_\mu$ is to be counted among them.}

The units $\tau_\beta$ [\S~199, (5)] do, however, belong to the normal units. Namely,
\begin{equation}
\tau_j=\tg\frac{5^j\pi}{m}=\frac{1-r_j}{1+r_j}r_j,
\tag{4}
\end{equation}
and if in this the substitution $(r,-r)=(r_\beta,r_{\beta+\nu'})$ is made, then $\tau_j$ goes over into
\begin{equation}
\tau_{\beta+\nu'}=-\frac1{\tau_\beta},
\tag{5}
\end{equation}
which corresponds to relation (3).

A system of $\nu'$ normal units
\[
\mathfrak E_0(r),\ \mathfrak E_1(r),\ldots,\mathfrak E_{\nu'-1}(r)
\]
shall be called an independent system if the determinant
\begin{equation}
\pm\sum\pm\log|\mathfrak E_0(r_0)|\log|\mathfrak E_1(r_1)|\cdots\log|\mathfrak E_{\nu'-1}(r_{\nu'-1})|
=L(\mathfrak E_0,\mathfrak E_1,\ldots,\mathfrak E_{\nu'-1})
\tag{6}
\end{equation}
has a value different from zero.

The absolute value of this determinant, $L$, we here also call the \textbf{regulator} of the system $\mathfrak E_0,\mathfrak E_1,\ldots,\mathfrak E_{\nu'-1}$.

The units $\tau_0,\tau_1,
\ldots,\tau_{\nu'-1}$ form an independent system of normal units; for their regulator
\[
L(\tau_0,\tau_1,\ldots,\tau_{\nu'-1})
\]
is just the determinant $T$ defined in the preceding paragraph [\S~199, (14), (15)]. And that this cannot be zero follows immediately from the fact that it occurs as a factor in the expression for the
class number. But the class number, since at least one class, namely the principal class, certainly always exists, is different from zero; hence $T$ too cannot be equal to zero. We therefore have the theorem:

\begin{center}\textbf{1. The numbers $\tau_0,\tau_1,\ldots,\tau_{\nu'-1}$ are an independent system of normal units of the field $H_m$.}\end{center}

If $\mathfrak E(r)$ is a normal unit, then, putting
\begin{equation}
l_\beta=\log|\mathfrak E(r_\beta)|,
\tag{7}
\end{equation}
one has, because of (3),
\begin{equation}
l_{\beta+\nu'}=-l_\beta.
\tag{8}
\end{equation}

If now $\mathfrak E_0,\mathfrak E_1,\ldots,\mathfrak E_{\nu'-1}$ is any independent system of normal units, we put
\begin{equation}
l_{i\beta}=\log|\mathfrak E_i(r_\beta)|,\qquad l_{i,\beta+\nu'}=-l_{i\beta},
\tag{9}
\end{equation}
so that the regulator of the system of the $\mathfrak E_i$ receives the expression
\begin{equation}
\pm\sum\pm l_{00}l_{11}\cdots l_{\nu'-1,\nu'-1}=L(\mathfrak E_0,\mathfrak E_1,\ldots,\mathfrak E_{\nu'-1})
\tag{10}
\end{equation}
and is a number different from zero, which we denote by $T_\mathfrak E$.

If $\mathfrak E(r)$ denotes any other normal unit, then we can determine the unknowns $\xi_0,\xi_1,\ldots,\xi_{\nu'-1}$ from the $\nu'$ linear equations which we obtain when, in
\begin{equation}
l_\beta=\xi_0l_{0\beta}+\xi_1l_{1\beta}+\cdots+\xi_{\nu'-1}l_{\nu'-1,\beta},
\tag{11}
\end{equation}
we let the index $\beta$ run from $0$ to $\nu'-1$; and because of (8) and (9), this equation is also correct for the remaining values of $\beta$, so that from it one may derive the formula, valid for all $r$:
\begin{equation}
\mathfrak E(r)=\pm\mathfrak E_0(r)^{\xi_0}\mathfrak E_1(r)^{\xi_1}\cdots\mathfrak E_{\nu'-1}(r)^{\xi_{\nu'-1}}.
\tag{12}
\end{equation}

Conversely, if $m_0,m_1,\ldots,m_{\nu'-1}$ denote arbitrary rational integers, then all numbers contained in the form
\[
\pm\mathfrak E_0(r)^{m_0}\mathfrak E_1(r)^{m_1}\cdots\mathfrak E_{\nu'-1}(r)^{m_{\nu'-1}}
\]
are normal units of the field $H_m$. Consequently the considerations which we carried out in detail in \S~186 may be applied to the system of equations (10), and this gives:

\begin{center}\textbf{2. The exponents $\xi_0,\xi_1,\ldots,\xi_{\nu'-1}$ determined from (10) or (11) are rational numbers.}\end{center}

From this one may infer, as in \S~187, that there are systems of $\nu'$ independent normal units whose regulator $T_0$ is as small as possible, and such systems are called \textbf{fundamental systems of normal units}.

If $\mathfrak E_0,\mathfrak E_1,\ldots,\mathfrak E_{\nu'-1}$ is such a fundamental system, then the exponents $\xi_0,\xi_1,\ldots,\xi_{\nu'-1}$ in (10) and (11) are obtained as \textbf{rational integers}. This too can be proved in the same way as in \S~187.

If, under this assumption concerning the $\mathfrak E_i$, we set up the equations (11) for any other system of independent units $\mathfrak E_0^{(1)},\mathfrak E_1^{(1)},\ldots,\mathfrak E_{\nu'-1}^{(1)}$, then we obtain expressions of the form
\[
L_{i,k}^{(1)}=\xi_{i,0}L_{0,k}+\xi_{i,1}L_{1,k}+\cdots+\xi_{i,\nu'-1}L_{\nu'-1,k}.
\]
Thus, if $T_1$ is the regulator of the system of the $\mathfrak E_i^{(1)}$, then, by the multiplication theorem for determinants, if $C$ denotes the absolute value of the determinant of the integral exponents $\xi_{i,k}$, hence a positive rational integer, there results
\begin{equation}
T_1=CT_0.
\tag{13}
\end{equation}

It cannot be proved, and for the higher values of $m$ probably also is not true, that $\tau_0,\tau_1,\ldots,\tau_{\nu'-1}$ form a fundamental system.\footnote{Here, as in all questions concerning units, it is very difficult to arrive at definite numerical results. According to the results contained in the rich work of Reuschle, \emph{Tafeln complexer Primzahlen, welche aus Wurzeln der Einheiten gebildet sind} (Berlin 1875), for $m=16$ the $\tau$ still form a fundamental system. For here the basic number of the field $H_m$ is equal to $2^{11}$, the degree of the field is equal to 4, and hence, by the note to \S~156, in every ideal class there is an ideal whose norm is less than 6. But according to Reuschle's tables every natural prime number below 100 can be decomposed in the field $H_m$ into actually existing prime factors. Consequently in all numbers below 6 only principal ideals certainly occur, and consequently there is here only the one class, the principal class. Since, as we have seen, the first class-number factor too is equal to 1, the cyclotomic field $\Omega_{16}$ is of one class. The same laws of primes hold in it as in the field of rational numbers. This is doubtful for the field $\Omega_{32}$ and certainly no longer applies in the field $\Omega_{64}$, in which the class number is a multiple of 17.}

For our purpose, however, it is important that these units behave, with respect to the denominator 2 in the exponents,
in a certain sense like a fundamental system; this is expressed more precisely by the following theorem:

\begin{center}\textbf{2. If a normal unit of the field $H_m$ is of the form
\[
\mathfrak E(r)=\tau_0^{e_0}\tau_1^{e_1}\cdots\tau_{\nu'-1}^{e_{\nu'-1}},
\]
where the exponents $e_0,e_1,
\ldots,e_{\nu'-1}$ are rational integers, and if it has one and the same sign with all its conjugate values, then all $e_0,e_1,\ldots,e_{\nu'-1}$ must be even numbers, and $\mathfrak E(r)$ is the square of a normal unit.}\end{center}

We can again conduct the proof by complete induction; but first we shall transform the expression for $\tau_j$ somewhat.

The conjugate values of the unit $\mathfrak E(r)$ are obtained in the form
\[
\mathfrak E(r_\beta)=\tau_\beta^{e_0}\tau_{\beta+1}^{e_1}\cdots\tau_{\beta+\nu'-1}^{e_{\nu'-1}},
\]
and these expressions are therefore supposed to have the same sign for all values of $\beta$.

By \S~199, (5),
\[
\tau_j=\tg\frac{5^j\pi}{m}=-\frac12\frac{\sin\left(\frac{2\pi}{m}5^j\right)}{\left[\cos\left(\frac{5^j\pi}{m}\right)\right]^2}.
\]
Here, for abbreviation, we put
\begin{equation}
\sigma_j=\sin\left(\frac{2\pi}{m}5^j\right),
\tag{14}
\end{equation}
so that $\sigma_j$ always has the same sign as $\tau_j$.

From this we form the product
\begin{equation}
S_\beta=\sigma_\beta^{e_0}\sigma_{\beta+1}^{e_1}\cdots\sigma_{\beta+\nu'-1}^{e_{\nu'-1}},
\tag{15}
\end{equation}
and our assumption is therefore that $S_\beta$ has the same sign for all values of $\beta$.

It is to be shown that the rational integers $e_0,e_1,\ldots,e_{\nu'-1}$ must all be even.

The theorem is correct for $m=8$; for in this case
\[
\sigma_0=\sin\frac{2\pi}{8},\qquad \sigma_1=\sin\frac{10\pi}{8}=-\sin\frac{6\pi}{8},
\]
therefore $\sigma_0$ is positive and $\sigma_1$ negative. Here $S_\beta=\sigma_\beta^{e_0}$, and $S_0$ and $S_1$ can have the same sign only when $e_0$ is even.

After this we apply complete induction. For simplification we put $\nu'=2\nu''$ and recall the notation
\begin{equation}
m=2^\kappa,\qquad \mu=\frac12m,\qquad \nu=\frac12\mu,
\qquad \nu'=\frac12\nu,
\qquad \nu''=\frac12\nu',
\tag{16}
\end{equation}
and assume that our theorem has already been proved when $\mu$ takes the place of $m$.

Then
\[
5^\nu\equiv1\pmod m,
\qquad 5^{\nu'}\equiv1+\mu\pmod m,
\qquad 5^{\nu''}\equiv1+\nu\pmod\mu,
\]
\[
5^{\nu''}\equiv1+\nu+\mu\pmod m,
\]
and from this it follows that
\[
\sigma_{\beta+\nu}=-\sigma_\beta,
\]
\[
\sigma_{\beta+\nu''}=-\cos\left(\frac{2\pi}{m}5^\beta\right),
\]
\begin{equation}
\sigma_\beta\sigma_{\beta+\nu''}=-\frac12\sin\left(\frac{2\pi}{\mu}5^\beta\right)=-\frac12\sigma'_\beta,
\tag{17}
\end{equation}
\begin{equation}
\sigma'_{\beta+\nu''}=-\sigma'_\beta,
\tag{18}
\end{equation}
where the $\sigma'_\beta$ arise from the $\sigma_\beta$ by the transition from $\kappa$ to $\kappa-1$. We now form, according to (17), (18), the product $S_\beta S_{\beta+\nu''}$, and obtain, if for abbreviation we also put
\[
e'=e_0+e_1+\cdots+e_{\nu'-1},
\qquad e''=e_0+e_1+\cdots+e_{\nu''-1},
\]
\[
S_\beta S_{\beta+\nu''}=(-1)^{\nu''}\left(\frac12\right)^{e'}
(\sigma'_\beta)^{e_0+e_{\nu''}}(\sigma'_{\beta+1})^{e_1+e_{\nu''+1}}\cdots
(\sigma'_{\beta+\nu''-1})^{e_{\nu''-1}+e_{\nu'-1}}.
\]
This, however, is a product of the same form as $S_\beta$, with $\mu$ in place of $m$, and it likewise has the same sign with all its conjugates. Therefore by the assumption made,
\begin{equation}
e_0\equiv e_{\nu''},\ e_1\equiv e_{\nu''+1},\ldots,e_{\nu''-1}\equiv e_{\nu'-1}\pmod2.
\tag{19}
\end{equation}
Putting, then,
\[
S'_\beta=(\sigma'_\beta)^{e_0}(\sigma'_{\beta+1})^{e_1}\cdots(\sigma'_{\beta+\nu''-1})^{e_{\nu''-1}},
\]
\[
R_\beta=\sigma_{\beta+\nu''}^{\frac12(e_0-e_{\nu''})}
\sigma_{\beta+\nu''+1}^{\frac12(e_1-e_{\nu''+1})}\cdots
\sigma_{\beta+\nu'-1}^{\frac12(e_{\nu''-1}-e_{\nu'-1})},
\]
it follows from (15) and (17) that
\[
S_\beta R_\beta=\left(-\frac12\right)^{e''}S'_\beta,
\]
and hence $S'_\beta$, just as $S_\beta$, has the same sign with all its conjugate values. But $S'_\beta$ arises from $S_\beta$ by putting $\mu$ in place of $m$, and therefore by our assumption
\[
e_0\equiv0,\ e_1\equiv0,\ldots,e_{\nu''-1}\equiv0\pmod2,
\]
and, with the aid of (19),
\[
e_{\nu''}\equiv0,\ e_{\nu''+1}\equiv0,\ldots,e_{\nu'-1}\equiv0\pmod2,
\]
which completes the proof of Theorem 2.

If now in formula (12) for $\mathfrak E_0,\mathfrak E_1,\ldots,\mathfrak E_{\nu'-1}$ we choose the system $\tau_0,\tau_1,\ldots,\tau_{\nu'-1}$, and put the exponents $\xi$ into the form
\[
\xi_0=\frac{e_0}{e},\quad \xi_1=\frac{e_1}{e},\ldots,\quad \xi_{\nu'-1}=\frac{e_{\nu'-1}}{e},
\]
where $e,e_0,e_1,\ldots,e_{\nu'-1}$ are rational integers without a common divisor, then we obtain
\[
\mathfrak E(r)^e=\pm\tau_0^{e_0}\tau_1^{e_1}\cdots\tau_{\nu'-1}^{e_{\nu'-1}}.
\]
It follows from our theorem that $e$ must be odd; for if it were even, then $\mathfrak E(r)^e$ would be the square of a real quantity, and consequently $\pm\mathfrak E(r)^e$ would have the same sign with its conjugates. But then, according to our theorem, all $e_0,e_1,\ldots,e_{\nu'-1}$ would have to be divisible by 2, contrary to the assumption that the exponents $e,e_0,
\ldots,e_{\nu'-1}$ should have no common divisor. We therefore have the theorem:

\begin{center}\textbf{3. If $\mathfrak E(r)$ is any normal unit of the field $H_m$, then the odd number $e$ and the rational integers $e_0,e_1,\ldots,e_{\nu'-1}$ can be so determined that}
\[
\mathfrak E(r_\beta)^e=\pm\tau_\beta^{e_0}\tau_{\beta+1}^{e_1}\cdots\tau_{\beta+\nu'-1}^{e_{\nu'-1}}
\tag{20}
\]
\textbf{is obtained.}\end{center}

This result leads to some important consequences:

\begin{center}\textbf{4. If $T$ is the regulator of the system of the $\tau_j$, and $T_0$ the regulator of a fundamental system of normal units in $H_m$, i.e. the minimal value which the regulator of any system of independent normal units can assume, then}
\[
T=CT_0,
\qquad C\equiv1\pmod2,
\tag{21}
\]
\textbf{where $C$ denotes an odd natural number.}\end{center}

For the fact that $T:T_0=C$ is an integer follows from formula (13).

If, however, we apply Theorem 3 to all elements of the fundamental system $\mathfrak E_i(r)$, then a system of integers $e_{i,k}$ and an odd number $e$ result such that
\[
\mathfrak E_i(r)^e=\pm\tau_0^{e_{i,0}}\tau_1^{e_{i,1}}\cdots\tau_{\nu'-1}^{e_{i,\nu'-1}}.
\]
Taking logarithms of these equations, and then forming the regulator $T_0$, we obtain
\[
e^{\nu'}T_0=\pm\sum\pm e_{0,0}e_{1,1}\cdots e_{\nu'-1,\nu'-1}T,
\]
hence
\[
e^{\nu'}=\pm C\sum\pm e_{0,0}e_{1,1}\cdots e_{\nu'-1,\nu'-1},
\]
and since in this equation the left-hand side is odd, $C$ cannot be even, as was to be proved.

If all conjugate values of the normal unit $\mathfrak E(r_\beta)$ have one and the same sign, then the same is also true of $\mathfrak E(r_\beta)^e$, and consequently in this case the exponents $e_0,e_1,\ldots,e_{\nu'-1}$ in formula (20) must be even numbers by 2. If one then writes formula (20) in the form
\[
\mathfrak E(r)=\pm\tau_0^{e_0}\tau_1^{e_1}\cdots\tau_{\nu'-1}^{e_{\nu'-1}}\mathfrak E(r)^{1-e},
\]
one obtains from it the theorem:

\begin{center}\textbf{5. A normal unit of the field $H_m$ which has the same sign with its conjugates is, apart from its sign, the square of a normal unit.}\end{center}

\section*{\S\,201. Fundamental system of units of the field $H_m$.}

It remains to determine the denominator $D$ in the formula for $B$ [\S~197, (6)], which was defined by the equation
\[
E=E'D,
\]
where $E$ and $E'$ denoted the regulators of the fields $H_m,H_\mu$.

Here, therefore, the matter is no longer normal units, but units of the field $H_m$ in general.

We assume a fundamental system of units in the field $H_m$:
\begin{equation}
\varepsilon_1(r),\ \varepsilon_2(r),\ldots,\varepsilon_{\nu-1}(r),
\tag{1}
\end{equation}
and let the conjugate logarithms of this system be
\begin{equation}
l_{1,k},\ l_{2,k},\ldots,l_{\nu-1,k},\qquad k=0,1,2,\ldots,\nu-1.
\tag{2}
\end{equation}
Alongside this we consider a fundamental system in the field $H_\mu$,
\begin{equation}
\varepsilon'_1,\ \varepsilon'_2,\ldots,\varepsilon'_{\nu'-1},
\tag{3}
\end{equation}
with the conjugate logarithms
\begin{equation}
l'_{1,k},\ l'_{2,k},\ldots,l'_{\nu'-1,k},\qquad k=0,1,2,\ldots,\nu'-1.
\tag{4}
\end{equation}

In forming the regulators, one of the conjugate values, say $k=0$, is omitted (cf. \S~186), and therefore we obtain
\begin{equation}
E=\pm\sum\pm l_{1,1}l_{2,2}\cdots l_{\nu-1,\nu-1},
\qquad
E'=\pm\sum\pm l'_{1,1}l'_{2,2}\cdots l'_{\nu'-1,\nu'-1}.
\tag{5}
\end{equation}

Now we add to the system (3) a fundamental system of normal units of the field $H_m$,
\begin{equation}
\mathfrak E_0,\mathfrak E_1,\ldots,\mathfrak E_{\nu'-1},
\tag{6}
\end{equation}
and assert that the system
\begin{equation}
\varepsilon'_1,\varepsilon'_2,
\ldots,\varepsilon'_{\nu'-1},\mathfrak E_0,\mathfrak E_1,
\ldots,\mathfrak E_{\nu'-1}
\tag{7}
\end{equation}
is an independent system of units of the field $H_m$ (even if not a fundamental system).

If we denote by
\begin{equation}
L_{0,k},L_{1,k},\ldots,L_{\nu'-1,k},\qquad k=0,1,\ldots,\nu'-1
\tag{8}
\end{equation}
the conjugate logarithms of the system (6), then by the preceding paragraph
\begin{equation}
\pm\sum\pm L_{0,0}L_{1,1}\cdots L_{\nu'-1,\nu'-1}=T_0
\tag{9}
\end{equation}
is the regulator of this system; and in view of the definition of normal units [\S~200, (9)] we have
\begin{equation}
L_{i,k+\nu'}=-L_{i,k},
\tag{10}
\end{equation}
and since the $\varepsilon'_i$ as numbers of the field $H_\mu$ remain unchanged by the substitution $(r,-r)$,
\begin{equation}
l'_{i,k+\nu'}=l'_{i,k}.
\tag{11}
\end{equation}

Accordingly the regulator $R$ of the system (7) is obtained from the determinant
\[
\small
\left|\begin{array}{cccc|cccc}
l'_{1,1}&l'_{2,1}&\cdots&l'_{\nu'-1,1}&L_{0,1}&L_{1,1}&\cdots&L_{\nu'-1,1}\\
\cdots&\cdots&&\cdots&\cdots&\cdots&&\cdots\\
l'_{1,\nu'-1}&l'_{2,\nu'-1}&\cdots&l'_{\nu'-1,\nu'-1}&L_{0,\nu'-1}&L_{1,\nu'-1}&\cdots&L_{\nu'-1,\nu'-1}\\
l'_{1,0}&l'_{2,0}&\cdots&l'_{\nu'-1,0}&-L_{0,0}&-L_{1,0}&\cdots&-L_{\nu'-1,0}\\
l'_{1,1}&l'_{2,1}&\cdots&l'_{\nu'-1,1}&-L_{0,1}&-L_{1,1}&\cdots&-L_{\nu'-1,1}\\
\cdots&\cdots&&\cdots&\cdots&\cdots&&\cdots\\
l'_{1,\nu'-1}&l'_{2,\nu'-1}&\cdots&l'_{\nu'-1,\nu'-1}&-L_{0,\nu'-1}&-L_{1,\nu'-1}&\cdots&-L_{\nu'-1,\nu'-1}
\end{array}\right|.
\]
This determinant has $\nu-1$ rows. We add the first to the $(\nu'+1)$-st, the second to the $(\nu'+2)$-nd, and so on, the $(\nu'-1)$-st to the last. In the last $(\nu'-1)$ rows the $L$ thereby cancel, and the $l'$ acquire the factor 2 [by (10), (11)]; and by a determinant theorem (vol. I, \S~23, V.) this gives
\begin{equation}
R=2^{\nu'-1}E'T_0.
\tag{12}
\end{equation}

Since, accordingly, $R$ does not vanish, and hence in (7) there is a system of $\nu-1$ independent units of the field $H_m$, we can [by \S~186, (5)] determine the rational (integral or fractional) numbers $m_{i,k},M_{i,k}$ from the equations
\begin{equation}
2l_{i,k}=m_{1,i}l'_{1,k}+\cdots+m_{\nu'-1,i}l'_{\nu'-1,k}
+M_{0,i}L_{0,k}+M_{1,i}L_{1,k}+\cdots+M_{\nu'-1,i}L_{\nu'-1,k}.
\tag{13}
\end{equation}
But it is easily shown, by means of the relations (10), (11), that the numbers $m_{i,k},M_{i,k}$ here must be \textbf{integral}. For
\[
l_{i,k}+l_{i,k+\nu'}=\log|\varepsilon_i(r_k)\varepsilon_i(-r_k)|
=m_{1,i}l'_{1,k}+\cdots+m_{\nu'-1,i}l'_{\nu'-1,k},
\]
\[
l_{i,k}-l_{i,k+\nu'}=\log|\varepsilon_i(r_k)\varepsilon_i(-r_k)^{-1}|
=M_{0,i}L_{0,k}+M_{1,i}L_{1,k}+\cdots+M_{\nu'-1,i}L_{\nu'-1,k}.
\]
But the units $\varepsilon_i(r)\varepsilon_i(-r)$ belong to the field $H_\mu$, and since by assumption the $\varepsilon'_i$ form a fundamental system of this field, the $m_{i,k}$ must be rational integers by \S~187.

The units $\varepsilon_i(r)\varepsilon_i(-r)^{-1}$ are normal units of the field $H_m$, and since the $\mathfrak E_i$ are supposed to be a fundamental system of normal units, the $M_{i,k}$ are integers by \S~200.

For better survey we put the equations (13), once again without the index $k$ which distinguishes the conjugate values from one another, here:
\begin{equation}
2l_i=m_{1,i}l'_1+m_{2,i}l'_2+\cdots+m_{\nu'-1,i}l'_{\nu'-1}
+M_{0,i}L_0+M_{1,i}L_1+\cdots+M_{\nu'-1,i}L_{\nu'-1}.
\tag{14}
\end{equation}
If we now form the regulator $E$ according to (13), and denote by $M$ the absolute value of the determinant of the $\nu-1$ rows of the numbers $m_{k,i},M_{k,i}$,
\[
M=\pm\sum\pm m_{1,1}\cdots m_{\nu'-1,\nu'-1}M_{0,0}M_{1,1}\cdots M_{\nu'-1,\nu'-1},
\]
then there follows
\begin{equation}
2^{\nu-1}E=MR,
\tag{15}
\end{equation}
and from this, in view of (12),
\begin{equation}
E=2^{-\nu'}ME'T_0.
\tag{16}
\end{equation}

It can now further be shown that $M$ is a power of 2 and is divisible by $2^{\nu'}$.

For since the system (1) was assumed to be a fundamental system in $H_m$, there are rational integers $n_{k,i},N_{k,i}$ satisfying the equations
\begin{equation}
\begin{aligned}
l'_i&=n_{1,i}l_1+n_{2,i}l_2+\cdots+n_{\nu-1,i}l_{\nu-1},&& i=1,2,\ldots,\nu'-1,\\
L_i&=N_{1,i}l_1+N_{2,i}l_2+\cdots+N_{\nu-1,i}l_{\nu-1},&& i=0,1,2,\ldots,\nu'-1.
\end{aligned}
\tag{17}
\end{equation}
If from these we form the determinant $R$, and denote by $N$ the absolute value of the determinant of the $n_{k,i},N_{k,i}$, then
\[
R=NE,
\]
and hence by (15)
\[
MN=2^{\nu-1};
\]
from this it follows that each of the two integers $M,N$ must be a power of 2.

To decide by which power of 2 the determinant $M$ is divisible, we first determine a system of rational integers $a_i$ without common divisor which satisfies the linear equations
\begin{equation}
\sum_{i=1}^{\nu-1}a_im_{k,i}=0\qquad(k=1,2,\ldots,\nu'-1).
\tag{18}
\end{equation}
If we then put
\[
\sum_{i=1}^{\nu-1}a_iM_{k,i}=\xi_k,
\]
we obtain from the equations (13)
\[
2\sum_{i=1}^{\nu-1}a_il_{i,k}=\xi_0L_{0,k}+\xi_1L_{1,k}+\cdots+\xi_{\nu'-1}L_{\nu'-1,k}.
\]
From this one obtains by (10)
\[
\sum_{i=1}^{\nu-1}a_il_{i,k+\nu'}=-\sum_{i=1}^{\nu-1}a_il_{i,k},
\]
and this means that
\[
\varepsilon_1^{a_1}\varepsilon_2^{a_2}\cdots\varepsilon_{\nu-1}^{a_{\nu-1}}
\]
is a normal unit of the field $H_m$. It then follows further, since $\mathfrak E_0,\mathfrak E_1,\ldots,\mathfrak E_{\nu'-1}$ is a fundamental system of normal units, that the numbers $\xi_0,\xi_1,\ldots,\xi_{\nu'-1}$ must be even. Hence:

\begin{center}\textbf{1. The system of equations (18) implies the congruences
\[
\sum_{i=1}^{\nu-1}a_iM_{k,i}\equiv0\pmod2\qquad(k=0,1,\ldots,\nu'-1).
\tag{19}
\]}\end{center}

But in the equations (18), whose number is only $\nu'-1$, there are $\nu-1$ unknowns $a_i$. Therefore there are several independent systems of solutions, which we may state more precisely as follows:

\begin{center}\textbf{2. There are $\nu'$ integral solutions of the system (18):}
\[
\begin{array}{cccc}
a_{1,1}&a_{2,1}&\cdots&a_{\nu-1,1}\\
a_{1,2}&a_{2,2}&\cdots&a_{\nu-1,2}\\
\cdots&\cdots&&\cdots\\
a_{1,\nu'}&a_{2,\nu'}&\cdots&a_{\nu-1,\nu'}
\end{array}
\tag{20}
\]
\textbf{such that from the matrix (20) a determinant of $\nu'$ rows can be formed which is an odd number.}\end{center}

For first there is a solution of (18),
\[
a_{1,1},a_{2,1},\ldots,a_{\nu-1,1},
\]
in which an odd number occurs. Suppose, as is permitted, since the ordering of the indices is immaterial here,
that this odd number is $a_{1,1}$; then we determine a second solution
\[
0,a_{2,2},a_{3,2},\ldots,a_{\nu-1,2},
\]
in which again an odd number, say $a_{2,2}$, occurs. Then we determine a third solution
\[
0,0,a_{3,3},\ldots,
\]
and continue setting equal to zero as long as the number of remaining unknowns is still greater than the number of equations.

In the last system, therefore, $\nu'-1$ zeros will occur, leaving $\nu'$ unknowns.

We can take these rows of numbers $a_{i,k}$ for the system (20), for then the determinant
\begin{equation}
\pm\sum\pm a_{1,1}a_{2,2}\cdots a_{\nu',\nu'}
\tag{21}
\end{equation}
reduces to the diagonal term $a_{1,1}a_{2,2}\cdots a_{\nu',\nu'}$, all whose factors are odd.

If now the requirement of Theorem 2 is fulfilled, the matrix (20) can be completed, by adding $\nu'-1$ rows, to a determinant of $\nu-1$ rows which likewise has an odd integral value:
\begin{equation}
\alpha=\pm\sum\pm a_{1,1}a_{2,2}\cdots a_{\nu',\nu'}a_{\nu'+1,\nu'+1}\cdots a_{\nu-1,\nu-1}.
\tag{22}
\end{equation}
It is only necessary, when the determinant (21) is assumed to be odd, to put zeros throughout the rows that have been added, except for the elements in the diagonal row, $a_{\nu'+1,\nu'+1},\ldots,a_{\nu-1,\nu-1}$, which may be taken equal to 1. Then $\alpha$ receives the same value as the determinant (21).

But if now, by the multiplication theorem for determinants, we form the product $\alpha M$, then there results a determinant in which the $\nu-1$ sums
\[
\sum_{i=1}^{\nu-1}a_i m_{k,i},
\qquad
\sum_{i=1}^{\nu-1}a_i M_{k,i},
\]
\[
k=1,2,\ldots,\nu'-1,
\qquad
k=0,1,2,\ldots,\nu'-1
\]
form a row for each value $s=1,2,\ldots,\nu'$; and in it, therefore, $\nu'$ rows occur whose elements are all divisible by 2.

Consequently $\alpha M$, and therefore also $M$, is divisible by $2^{\nu'}$, and if we put
\begin{equation}
M=2^{\nu'+\theta},\qquad N=2^{\nu'-\theta-1},
\tag{23}
\end{equation}
then $\theta$ is a non-negative rational integer.

Putting this into (16) gives
\[
E=2^\theta E'T_0.
\]
But in \S~197, (6), one had put
\[
E=E'D,
\]
and hence it follows that
\[
D=2^\theta T_0.
\]
From this, by \S~199, (15) and \S~200, (21), one obtains for the class-number factor $B$
\begin{equation}
B=2^{-\theta}\frac{T}{T_0}=2^{-\theta}C,
\tag{24}
\end{equation}
where $C$ is an odd integer.

Accordingly, by \S~197, (4), for the second factor of the class number $h_1$, which, as we have seen, is an integer, one obtains
\begin{equation}
h_1=2^{-\theta}h'_1C.
\tag{25}
\end{equation}
If we now assume as proved that $h'_1$ is odd, then, since $C$ also is odd, it follows that $\theta=0$ and $h_1$ is odd; and both are thus proved by complete induction.

The derived results are thereby simplified as follows:

\begin{center}\textbf{3. One has}
\[
M=2^{\nu'},\qquad D=T_0,
\qquad B=\frac{T}{T_0};
\tag{26}
\]
\textbf{the class-number factor $B$ is equal to the ratio of the regulator $T$ of the units $\tau$ to the regulator $T_0$ of a fundamental system of normal units, and is an odd integer.}\end{center}

If we denote by $B_\kappa$ the number $B$ belonging to $m=2^\kappa$, then the class number $h_1$, as follows by complete induction, is
\begin{equation}
h_1=B_4B_5\cdots B_\kappa,
\tag{27}
\end{equation}
and hence is also an odd number. Thus, with regard to \S~198, 1., the main theorem is proved:

\begin{center}\textbf{A. The class number in the field $\Omega_m$, when $m$ is a power of 2, is odd.}\end{center}

\section*{\S\,202. Positive units.}

The results found last show that the system of independent units \S~201, (7),
\[
\varepsilon'_1,\varepsilon'_2,\ldots,\varepsilon'_{\nu'-1},\mathfrak E_0,\mathfrak E_1,\ldots,\mathfrak E_{\nu'-1},
\]
in the field $H_m$ is not a fundamental system. For otherwise in the formulae \S~201, (13), all the numbers $m_{k,i},M_{k,i}$ would have to be divisible by 2, and the determinant $M$, whose value we found to be $2^{\nu'}$, would be divisible by $2^{\nu-1}$.

In place of the system of equations (18), \S~201, we now consider the following congruences:
\begin{equation}
\begin{aligned}
\sum_{i=1}^{\nu-1}b_i m_{1,i}&\equiv1\pmod2,\\
\sum_{i=1}^{\nu-1}b_i m_{2,i}&\equiv0\pmod2,\\
&\ \vdots\\
\sum_{i=1}^{\nu-1}b_i m_{\nu'-1,i}&\equiv0\pmod2,
\end{aligned}
\tag{1}
\end{equation}
and prove that it is possible to satisfy them by integral values of the $b_i$.

If this were not possible, then from the $\nu'-2$ congruences
\[
\sum_{i=1}^{\nu-1}b_i m_{k,i}\equiv0\pmod2\qquad k=2,3,\ldots,\nu'-1
\]
the last congruence
\[
\sum_{i=1}^{\nu-1}b_i m_{1,i}\equiv0\pmod2
\]
would have to follow; and then, as in the preceding paragraph, it would further follow that
\[
\sum_{i=1}^{\nu-1}b_i M_{k,i}\equiv0\pmod2\qquad k=0,1,\ldots,\nu'-1.
\]

Then in place of the matrix \S~201, (20), one could put a matrix of the $b_{k,i}$ with $\nu'+1$ rows, and on the path taken in the preceding paragraph it would follow that $M$ would have to be divisible by $2^{\nu'+1}$, which is not the case.

But if we apply the system of multipliers $b_i$ which satisfies the conditions (1) to the equations (14) of the preceding paragraph, multiplying by $b_i$ and then summing, there results, when we omit multiples of 2 and express this also here, where irrational numbers, namely logarithms, occur, by the congruence sign:
\begin{equation}
l'_1\equiv L_0\sum M_{0,i}b_i+L_1\sum M_{1,i}b_i+\cdots+L_{\nu'-1}\sum M_{\nu'-1,i}b_i\pmod2.
\tag{2}
\end{equation}
If we wish to transform the congruence into an equation, then a sum of quantities
\[
l_1,l_2,\ldots,l_{\nu-1},\quad l'_1,l'_2,\ldots,l'_{\nu'-1}
\]
is to be added, with even integers as coefficients.

The same consideration may also be applied to $l'_2,l'_3,\ldots,l'_{\nu'-1}$; and from this, when one passes back from the logarithms to the numbers, formula (2) may be expressed in words as follows:

\begin{center}\textbf{1. Every unit of the field $H_\mu$ is representable as the product of a normal unit of the field $H_m$ and the square of a unit in $H_m$.}\end{center}

If we now apply Theorem 2, \S~200 to this, then it follows that every unit $\mathfrak E(r)$ of the field $H_\mu$ which has the same sign with all its conjugates must, apart from sign, be the square of a unit in $H_m$. But if
\[
\pm\mathfrak E(r)=\Delta(r)^2,
\]
where $\mathfrak E(r)$ is a unit in $H_\mu$ and $\Delta(r)$ a unit in $H_m$, then $\Delta(r)$ too must be contained in $H_\mu$.

For if we put
\[
\Delta(r)=\Delta_1(r^2)+r\Delta_2(r^2),
\]
then its square can be contained in $H_\mu$, and so remain unchanged under the change of sign of $r$, only if either $\Delta_1$ or $\Delta_2=0$. But it cannot be that $\Delta(r)=r\Delta_2(r^2)$, for otherwise $\Delta_2(r^2)$ would be a unit in $\Omega_\mu$ and would have to be,
by \S~178, 1., the product of a real unit and a root of unity $r^{2\lambda}$; and since $\Delta(r)$ too is real, $r^{2\lambda+1}$ would have to be real, which is impossible. Hence $\Delta(r)$ is contained in $H_\mu$.

Since now $\mu$, as well as $m$, can be any power of 2, the second main theorem follows from this:

\begin{center}\textbf{B. A unit of the field $H_m$ which is positive with all its conjugates is the square of a unit of the same field.}\end{center}

On the two theorems A. and B., however, in \S~184 we made the proof of the main theorem on Abelian number fields (\S~179, I.) depend.

\clearpage
\begin{center}
{\large Twenty-fifth Section.}\\[0.8em]
{\Large Transcendent Numbers.}\\[1.2em]
{\large §. 203.}\\[0.4em]
{\large Countable Sets.}
\end{center}

In the introduction to our work the general concept of numbers was defined, and with this general concept of number we have at first operated, for example in the proof of the existence of roots. In the further course of our investigations we have then occupied ourselves only with algebraic numbers, without giving account to ourselves whether the extent of the number-domain was thereby exhausted, or whether there are also non-algebraic numbers. The existence of non-algebraic numbers, which are also called transcendent numbers, was first proved by Liouville. Other proofs of this theorem are due to G. Cantor\footnote{G. Cantor, Crelle's Journal, vol. 77 (1873). Ueber eine Eigenschaft des Inbegriffes aller reellen algebraischen Zahlen.}.

We attach ourselves here to the concept, already set forth in the introduction, of a set or manifold; by this is to be understood a system of elements of any kind whatever, which is so precisely delimited that of every arbitrary object it is fully decided whether it belongs to the system or not.

We distinguish finite and infinite sets, and as the first and most important example of an infinite set we cite the totality of the natural numbers $1,2,3,\ldots$. Then the following definition holds:

\textbf{1. Definition.} A set is called countable if its elements can be put, with the natural number sequence
or with a part of the same, into a reciprocal one-to-one relation\footnote{Cantor, loc. cit. The concept of countable sets coincides with Dedekind's concept, defined without presupposing the number-system, of simply infinite systems. (Dedekind, Was sind und was sollen die Zahlen? §. 6. Braunschweig 1887. Second unchanged edition 1893.)}.

Every finite set is accordingly countable, and in the counting the natural number sequence is used only up to a certain greatest number. In what follows we consider preferably infinite sets.

We can also give the definition for an infinite countable set the expression that it is such a set in which to each element a definite number of the natural number sequence can be assigned as name, in such a way that every number of the number sequence is also used, hence a set in which there is a first, second, third $\ldots$ hundredth $\ldots$ element.

Finally we can also say that an infinite countable set is one which can be ordered into a sequence in such a way that a first element is present, and that after each element a definite other element of the set follows, and before each element, with the exception of the first, a definite other element precedes. Such a sequence we can call a countable ordering.

It is clear that a countable set can be counted not only in one way, but in infinitely many different ways.

Besides the natural number sequence itself, which is certainly countable, we can cite as a second example the system of the rational positive proper fractions, which among other ways can be counted in the following way:
\[
\frac12,\ \frac13,\ \frac23,\ \frac14,\ \frac34,\ \frac15,\ \frac25,\ \frac35,\ \frac45,\ \frac16,\ \frac56,\ldots
\]
that is, in such a way that every larger denominator follows the smaller denominator, and that the fractions with the same denominator are ordered according to the size of the numerator. But if one wanted to order the fractions according to the size of their numerical value, one would not obtain a countable ordering.

\textbf{2. The totality of all algebraic numbers is a countable set.}

To prove this important theorem, we recall that every algebraic number $\Theta$ is the root of one and only one irreducible equation
\begin{equation}\tag{1}
f(\Theta)=a_0\Theta^n+a_1\Theta^{n-1}+\cdots+a_n=0
\end{equation}
in which $a_0,a_1,\ldots,a_n$ are rational integers without a common divisor, and $a_0$ is different from zero and positive. The degree $n$ of the equation (1) is a positive integer, therefore at least $=1$. The numbers $a_1,\ldots,a_{n-1}$ can in part be zero.

We shall now determine the signs $\pm$ so that $\pm a_1,\pm a_2,\pm a_3,\ldots,\pm a_n$ are not negative, and call the sum
\begin{equation}\tag{2}
N=(n-1)+a_0\pm a_1\pm a_2\cdots\pm a_n
\end{equation}
the height of the algebraic number $\Theta$. The height is then always a positive integer.

Now it is easy to see that for a given value of the height $N$ there can always be only a finite number of algebraic numbers. For first, by (2), $n$ can never be greater than $N$, and for given $N$ and $n$ the numbers $a_0,a_1,\ldots,a_n$ can be determined only in a finite number of ways. Among the functions $f(\Theta)$ so determined one then has to retain only the irreducible ones. If one now orders the algebraic numbers in such a way that one puts the numbers of smaller height before those of greater height, that among numbers of equal height one puts first those whose real part is smaller, and among numbers of equal height and equal real part lets those with smaller imaginary part precede, then we have a countable ordering of the algebraic numbers, and it is proved that the totality of all algebraic numbers is a countable set.

For example, we obtain:
\[
\begin{array}{llll}
N=1, & n=1, & a_0=1, & a_1=0,\\
N=2, & n=1, & a_0=1, & a_1=\pm1,\\
N=3, & n=1, & a_0=1, & a_1=\pm2,\\
     &      & a_0=2, & a_1=\pm1,\\
     & n=2, & a_0=1, & a_1=0,\quad a_2=1,
\end{array}
\]
and the beginning of the ordered sequence of algebraic numbers is therefore:
\[
0,\ -1,\ +1,\ -2,\ -\frac12,\ -\sqrt{-1},\ \sqrt{-1},\ \frac12,\ 2,\ldots .
\]

Every part of the natural number sequence is a countable set; for one need only order the numbers of such a part according to their size in order to obtain a countable ordering.

From this, however, it follows that every part of a countable set is itself a countable set. It follows from this, among other things, that the set of real algebraic numbers is also countable.

\begin{center}
{\large §. 204.}\\[0.4em]
{\large Uncountable Sets.}
\end{center}

We now come secondly to the proof of the theorem that among sets of numbers there are also those which are not countable, and we shall prove specifically:

\begin{center}
\textbf{That the totality of all real numbers, even if we restrict ourselves to a finite interval, is not countable.}
\end{center}

For this purpose we consider any countable set of real numbers distinct from one another, which, put into a countable ordering $\Omega$, we shall denote as
\[
(\Omega)\qquad \omega_1,\omega_2,\omega_3,\omega_4,\ldots
\]
(where it is to be remembered that this is not meant to be a succession according to size).

For brevity we shall call, of two elements of the sequence $\Omega$, the one with the smaller index the earlier, and the one with the greater index the later.

We now take any two real numbers $\alpha,\beta$, so that $\alpha<\beta$, and show that in the interval $\delta=(\alpha,\beta)$ there is at least one number which does not occur in the sequence $\Omega$. If we have proved such a number for every interval, then there are also infinitely many of them, since the same inference can be applied to every part of the interval $(\alpha,\beta)$.

First it is clear that our assertion is correct if in some finite part of the interval $\delta$ only a finite number of numbers of the sequence $\Omega$ is contained, and we can therefore pass at once to the case in which in every, however small, part of the interval $\delta$ there lies an infinite set of numbers of the sequence $\Omega$.

We denote by $\alpha_1,\beta_1$ the two earliest numbers of the sequence $\Omega$ which lie in the interval $\delta$, assume $\alpha_1<\beta_1$, and set $\delta_1=\beta_1-\alpha_1$, so that $\delta_1=(\alpha_1,\beta_1)$ is a part of the interval $\delta$.

We denote in the same way by $\alpha_2,\beta_2$ the two earliest numbers of $\Omega$ which lie in the interior of the interval $\delta_1$ (excluding the endpoints), assume $\alpha_2<\beta_2$, and set $\delta_2=\beta_2-\alpha_2$.

In this way we can continue and obtain an unbounded sequence of intervals
\[
\delta,\ \delta_1,\ \delta_2,\ \delta_3,\ldots,
\]
each of which includes all the following ones, and two sequences of numbers
\[
\alpha,\alpha_1,\alpha_2,\ldots,\qquad
\beta,\beta_1,\beta_2,\ldots
\]
which, with the possible exception of the first, $\alpha$ and $\beta$, all belong to the sequence $\Omega$. The $\alpha,\alpha_1,\alpha_2,\ldots$ form an increasing, the $\beta,\beta_1,\beta_2,\ldots$ a decreasing number-sequence, and at the same time every $\alpha$ is smaller than every $\beta$.

\textbf{1. From this it follows that the numbers $\alpha_\nu$ have an upper bound $a$, the numbers $\beta_\nu$ a lower bound $b$, and that $a$ in any case is not greater than $b$.} (Compare vol. I, §. 38, 1.)

But possibly $a=b$.

From the way in which the intervals $\delta,\delta_1,\delta_2,\ldots$ are formed, the following also emerges:

\textbf{2. If any number $\omega$ of the sequence $\Omega$ lies in the interval $\delta_\nu$, excluding its endpoints $\alpha_\nu,\beta_\nu$, then $\omega$ is later in the sequence $\Omega$ than the number-pair $\alpha_\nu,\beta_\nu$ belonging to that same sequence.}

For $\alpha_\nu,\beta_\nu$ were precisely the two earliest numbers of $\Omega$ lying in the interval $\delta_{\nu-1}$. From this it follows:

\textbf{3. The number-pairs $\alpha_\nu,\beta_\nu$ belonging to the sequence $\Omega$ are later members of the sequence $\Omega$ the larger the index $\nu$ is, and since we have assumed that the sequence of intervals $\delta_\nu$ does not break off, we can for sufficiently large $\nu$ push $\alpha_\nu,\beta_\nu$ arbitrarily far out in the sequence $\Omega$.}

It now follows very simply that no number $g$ which coincides with one of the numbers $a,b$, or also, if $a$ and $b$ are different, lies between them, can belong to the sequence $\Omega$.

For the number $g$ lies in the interior of each of the intervals $\delta_\nu$. Suppose that $g$ occurs in $\Omega$, and by 3. go so far with $\nu$ that $\alpha_\nu,\beta_\nu$ occur in $\Omega$ later than $g$; then, by 2., $g$ cannot lie in the interval $\delta_\nu$, and thereby the impossibility of our assumption is shown.

Thus it follows that the totality of the numbers of an interval $\alpha,\beta$ does not form a countable set.

\medskip\hrule\medskip

This fact can also be proved in another way, which in a certain respect is still simpler and which we shall present in few words. We do not restrict the generality if we take as basis the interval from 0 to 1. We think of all numbers of this interval as represented by infinite decimal fractions. Among them the finite decimal fractions are also contained, if from a certain place on we put all digits equal to zero. In order to make the representation by decimal fractions unique, let it still be fixed that for a finite decimal fraction this representation is always chosen, hence for example that $0{,}4999\ldots$ is not to be put for $0{,}5000\ldots$.

We now want to assume that these decimal fractions form a countable set; they can therefore be ordered in a countable sequence, which we write as follows:
\[
(\Omega)\qquad
\begin{array}{rcl}
\omega_1&=&0,\alpha_1^{(1)}\alpha_2^{(1)}\alpha_3^{(1)}\ldots\\
\omega_2&=&0,\alpha_1^{(2)}\alpha_2^{(2)}\alpha_3^{(2)}\ldots\\
\omega_3&=&0,\alpha_1^{(3)}\alpha_2^{(3)}\alpha_3^{(3)}\ldots\\
&&\ldots
\end{array}
\]
where the $\alpha_\mu^{(\nu)}$ denote digits of the decimal system.

It is now very easy to exhibit a decimal fraction (or also arbitrarily many) which are not contained in the sequence $\Omega$. We need only form
\[
\eta=0,\beta_1\beta_2\beta_3\ldots,
\]
where the $\beta_\nu$ are digits of the decimal system satisfying the single condition that, for every $\nu$, $\beta_\nu$ is different from $\alpha_\nu^{(\nu)}$.
This number $\eta$, which after all also belongs to the interval $(0,1)$, cannot agree with any number of the sequence $\Omega$.

One can generalize the formation of $\eta$ still further by taking the first $\beta$'s, up to an arbitrarily far removed place, arbitrarily, and only from there on requiring the law: $\beta_\nu$ different from $\alpha_\nu^{(\nu)}$, to hold.

Since it has now been proved that the real algebraic numbers form a countable set, and that in every interval there are numbers which do not belong to a given countable set, it follows from this quite immediately:

\begin{center}
\textbf{In every real interval there are transcendent numbers.}
\end{center}

\begin{center}
{\large §. 205.}\\[0.4em]
{\large Transcendence of the Number $e$.}
\end{center}

A much more difficult task is now, for a definitely presented number, to decide whether it is algebraic or transcendent. Here interest has concentrated chiefly on the two numbers, so frequently occurring in analysis, $e$, that is, the base of the natural logarithm system, and Ludolph's number $\pi$, the ratio of the circumference of the circle to the diameter.

For the number $e$ the question was decided by Hermite in a famous paper, which became the foundation for the later investigations on the number $\pi$\footnote{Hermite, Sur la fonction exponentielle. Comptes rendus, vol. LXXVII, 1873.}. The decision for the number $\pi$, which because of its relation to the ancient and famous problem of the quadrature of the circle was of quite special interest, still offered insurmountable difficulties for a long time. Finally Lindemann gave the proof that $\pi$ too belongs to the transcendent numbers. The proof given by Lindemann, however, at first still offered great difficulties for understanding, which were removed by later investigations of Weierstrass, Hilbert, Hurwitz, and Gordan\footnote{Lindemann, Ueber die Zahl $\pi$. Mathem. Annalen, vol. 20, 1882. Weierstrass, Zu Lindemann's Abhandlung ``Ueber die Ludolph'sche Zahl''. Sitzungsbericht der Berliner Akademie, 3 December 1885. The works of Hilbert, Hurwitz, and Gordan are all three found in volume 43 of the Mathematische Annalen (1893), the first two also in the Göttinger Nachrichten of 1893.}
The proof of Weierstrass has simplified the form of the proof quite considerably, and this proof is to be set forth in what follows.

We denote the product of the natural numbers from 1 to $n$ by
\[
\Pi(n)=1.2.3\cdots n.
\]
If $x$ is an arbitrary real or complex number, then it can first be shown that, for an arbitrarily given value of $x$, the terms of the sequence
\begin{equation}\tag{1}
\frac{x^n}{\Pi(n)}
\end{equation}
approach the limit zero as $n$ grows without bound, and indeed more strongly than the terms of a decreasing geometric progression. For if $k$ is an integer greater than the absolute value of $x$, then
\[
\left|\frac{x^n}{\Pi(n)}\right|
=
\left|\frac{x^k}{\Pi(k)}\frac{x}{k+1}\frac{x}{k+2}\cdots\frac{x}{n}\right|
<
\left|\frac{x^k}{\Pi(k)}\right|\left|\frac{x}{k}\right|^{\,n-k}.
\]
From this it follows that the infinite series
\begin{equation}\tag{2}
e^x=1+x+\frac{x^2}{\Pi(2)}+\frac{x^3}{\Pi(3)}+\cdots
\end{equation}
converges for all values of $x$, and by it we define the exponential function $e^x$, whose simplest properties we here presuppose from analysis. In particular, for two arbitrary values $x,y$ the relation
\[
e^{x+y}=e^x e^y
\]
holds, from which it then follows that $e^n$, for a positive integer $n$, is the $n$th power of the number
\begin{equation}\tag{3}
e=2+\frac1{\Pi(2)}+\frac1{\Pi(3)}+\frac1{\Pi(4)}+\cdots
=2{,}718281828459\ldots
\end{equation}
If now $r$ denotes any positive integer, we can write formula (2) in the form
\begin{equation}\tag{4}
\Pi(r)e^x=\Pi(r)+\frac{\Pi(r)}{\Pi(1)}x+\frac{\Pi(r)}{\Pi(2)}x^2+\cdots+x^r+x^rU_r,
\end{equation}
where
\begin{equation}\tag{5}
U_r=\frac{x}{r+1}+\frac{x^2}{(r+1)(r+2)}+\frac{x^3}{(r+1)(r+2)(r+3)}+\cdots,
\end{equation}
and this formula also holds for $r=0$, if one assumes $\Pi(0)=1$.

If we denote the absolute value of $x$ by $\xi$, then, by the theorem proved in the introduction that the absolute value of a sum is never greater than the sum of the absolute values of its parts, the absolute value of $U_r$ is certainly not greater than
\[
\frac{\xi}{r+1}
+\frac{\xi^2}{(r+1)(r+2)}
+\frac{\xi^3}{(r+1)(r+2)(r+3)}
+\cdots,
\]
hence smaller than
\[
1+\frac{\xi}{1}+\frac{\xi^2}{1.2}+\frac{\xi^3}{1.2.3}+\cdots=e^\xi.
\]
If therefore we put
\[
U_r=q_re^\xi,
\]
then $q_r$ is a function of $x$ about which we establish only so much, that its absolute value is smaller than 1.

Accordingly we write formula (4) for $r=0,1,\ldots,n$, and write the terms, ordered in reverse order, in the following survey:
\begin{align}
\Pi(n)e^x&=q_nx^n e^\xi+x^n+nx^{n-1}+n(n-1)x^{n-2}+\cdots+\Pi(n),\notag\\
\Pi(n-1)e^x&=q_{n-1}x^{n-1}e^\xi+x^{n-1}+(n-1)x^{n-2}
 +(n-1)(n-2)x^{n-3}+\cdots,\tag{6}\\
&\hspace{5em}\vdots\notag\\
e^x&=q_0e^\xi+1.\notag
\end{align}

Now take an arbitrary integral function of degree $n$ in $x$:
\begin{equation}\tag{7}
f(x)=c_nx^n+c_{n-1}x^{n-1}+c_{n-2}x^{n-2}+\cdots+c_0,
\end{equation}
and its derivatives
\begin{align}
f'(x)&=nc_nx^{n-1}+(n-1)c_{n-1}x^{n-2}+\cdots,\\
f''(x)&=n(n-1)c_nx^{n-2}+(n-1)(n-2)c_{n-1}x^{n-3}+\cdots,\tag{8}\\
&\hspace{5em}\vdots\notag\\
f^{(n)}(x)&=\Pi(n)c_n,\notag
\end{align}
and put, for abbreviation,
\begin{equation}\tag{9}
F(x)=f(x)+f'(x)+f''(x)+\cdots+f^{(n)}(x).
\end{equation}

Then $F(x)$ is an integral function of $x$ of degree $n$. Finally set
\begin{align}
Q(x)&=c_nq_nx^n+c_{n-1}q_{n-1}x^{n-1}+\cdots+c_0q_0,\tag{10}\\
P&=c_n\Pi(n)+c_{n-1}\Pi(n-1)+\cdots+c_0,\tag{11}
\end{align}
so that $Q(x)$ depends on $x$, even though it is not expressible rationally by $x$; but $P$ is independent of $x$.

If we now multiply the equations (6) in order by $c_n,c_{n-1},\ldots,c_0$ and add, it follows that
\begin{equation}\tag{12}
e^xP=F(x)+e^\xi Q(x).
\end{equation}
This formula is now the starting point of the further conclusions.

Suppose that $e$ were an algebraic number. Then there must exist an equation, of degree $m$:
\begin{equation}\tag{13}
C_0+C_1e+C_2e^2+\cdots+C_me^m=0,
\end{equation}
whose coefficients $C_0,C_1,\ldots,C_m$ are rational integers, of which $C_0$ and $C_m$ are different from zero. The point is to show the impossibility of this assumption.

For this purpose we put in equation (12), for $x$ successively, the integers $0,1,2,\ldots,m$, so that $\xi$ becomes identical with $x$, multiply by $C_0,C_1,\ldots,C_m$, and add. Then, by (13), since $P$ is independent of $x$, there results
\begin{align}
0={}&C_0F(0)+C_1F(1)+\cdots+C_mF(m)\notag\\
&+C_0Q(0)+C_1eQ(1)+\cdots+C_me^mQ(m),\tag{14}
\end{align}
and now it is to be proved, from a suitable assumption about the still arbitrary function $f(x)$, that equation (14) is impossible.

We choose a prime number $p$ greater than $m$\footnote{That there is always a prime number $p$ greater than an arbitrary number $\mu$ was already proved by Euclid (Elements, Book IX, no. XX, vol. 2 of Heiberg's edition). The proof is simply that the integer $\Pi(\mu)+1$, which is plainly greater than $\mu$, is divisible by no prime number not greater than $\mu$, since on division by each of the numbers $2,3,\ldots,\mu$ this number leaves the remainder 1. It is therefore impossible that no prime number lies above $\mu$.}, and put
\begin{equation}\tag{15}
f(x)=\frac{x^{p-1}(1-x)^p(2-x)^p\cdots(m-x)^p}{\Pi(p-1)},
\end{equation}
so that the degree $n$ of $f(x)$ is equal to $(m+1)p-1$, and we now prove two things:
\begin{enumerate}[label=\arabic*)]
\item $C_0F(0)+C_1F(1)+\cdots+C_mF(m)$ is an integer different from zero, hence, apart from sign, at least equal to 1;
\item $C_0Q(0)+C_1eQ(1)+\cdots+C_me^mQ(m)$ is smaller than 1;
\end{enumerate}
both under the presupposition that $p$ is suitably disposed of. If both of these are proved, then one recognizes the impossibility of equation (14), and hence also that of equation (13), from which (14) had been inferred.

If we order the numerator of $f(x)$ by powers of $x$, an expression of the form
\begin{equation}\tag{16}
f(x)=\frac{A_{p-1}x^{p-1}+A_px^p+A_{p+1}x^{p+1}+\cdots}{\Pi(p-1)}
\end{equation}
is obtained, where $A_{p-1},A_p,A_{p+1},\ldots$ are integers, and $A_{p-1}=\Pi(m)^p$, and hence is certainly not divisible by $p$. Therefore, if one compares (16) with the Taylor expansion
\[
f(x)=f(0)+xf'(0)+\frac{x^2}{\Pi(2)}f''(0)+\cdots,
\]
one has
\[
f(0)=f'(0)=f''(0)=\cdots=f^{(p-2)}(0)=0,
\]
\[
f^{(p-1)}(0)=A_{p-1},\quad f^{(p)}(0)=pA_p,\quad
f^{(p+1)}(0)=p(p+1)A_{p+1},\ldots,
\]
therefore
\[
F(0)=A_{p-1}+pA_p+p(p+1)A_{p+1}+\cdots,
\]
an integer not divisible by $p$. But if we order $f(x)$ by powers of $x-1$, it follows that
\[
f(x)=\frac{B_p(x-1)^p+B_{p+1}(x-1)^{p+1}+\cdots}{\Pi(p-1)},
\]
where $B_p,B_{p+1},\ldots$ are again integers. It follows from this, as above, by comparison with
\[
f(x)=f(1)+(x-1)f'(1)+\frac{(x-1)^2}{\Pi(2)}f''(1)+\cdots,
\]
that
\[
F(1)=pB_p+p(p+1)B_{p+1}+\cdots,
\]
and hence $F(1)$ is an integer divisible by $p$, and in exactly the same way one obtains that $F(2),F(3),\ldots,F(m)$ are integers divisible by $p$. Since one can now also choose $p$ so large that $C_0$ is not divisible by $p$, it follows

that $C_0F(0)+C_1F(1)+\cdots+C_mF(m)$ is an integer not divisible by $p$, hence also not vanishing, and with this 1) is proved.

If we can now still show that $Q(x)$, for every positive $x$, can be made arbitrarily small by increasing $p$, it follows that, for sufficiently large $p$, the expression 2) is certainly smaller than 1, and our proof is complete.

But if we go back to the expression (10) for $Q(x)$ and denote by $\gamma_n,\gamma_{n-1},\ldots,\gamma_0$ the absolute values of the coefficients $c_n,c_{n-1},\ldots,c_0$ of $f(x)$, then we see, since the $q_n,q_{n-1},\ldots$ are in absolute value smaller than 1, that for every positive $x$ the absolute value of $Q(x)$ is smaller than
\begin{equation}\tag{17}
\psi(x)=\gamma_nx^n+\gamma_{n-1}x^{n-1}+\cdots+\gamma_0.
\end{equation}
The coefficients $c_n,c_{n-1},\ldots,c_0$ of the function $f(x)$ in (15), however, differ from the $\gamma_n,\gamma_{n-1},\ldots,\gamma_0$ only by sign. Therefore, if we replace $x$ by $-x$, thus form the function
\[
f(-x)=\frac{x^{p-1}(1+x)^p(2+x)^p\cdots(m+x)^p}{\Pi(p-1)},
\]
this function has the same coefficients as $f(x)$, but all with positive signs. Thus it is no other function than $\psi(x)$.

If we therefore also put
\[
X=x(1+x)(2+x)\cdots(m+x),
\]
then
\begin{equation}\tag{18}
\psi(x)=\frac{X}{x}\cdot\frac{X^{p-1}}{\Pi(p-1)},
\end{equation}
and this approaches, as was proved at the beginning of this paragraph, the limit zero as $p$ grows without bound.

Thus $e$ is a transcendent number.

\begin{center}
{\large §. 206.}\\[0.4em]
{\large Transcendence of the Number $\pi$.}
\end{center}

With the same aids one can now also prove the transcendence of the number $\pi$. As definition of this number we use the fact that it is the smallest positive number satisfying the equation
\begin{equation}\tag{1}
e^{i\pi}=-1,
\end{equation}
when $e^x$ is defined by the formula §. 205, (2).

Suppose then that $\pi$, and consequently also $i\pi$, is an algebraic number. Then $i\pi$ is one of the roots of an irreducible rational equation $\chi(x)=0$ whose coefficients are rational integers.

Let all the roots of this algebraic equation be denoted by $\beta_1,\beta_2,\ldots,\beta_\nu$, and let the coefficient of $x^\nu$ in $\chi$ be denoted by $a$; then
\begin{equation}\tag{2}
\chi(x)=a(x-\beta_1)(x-\beta_2)\cdots(x-\beta_\nu),
\end{equation}
and the products $a\beta_1,a\beta_2,\ldots,a\beta_\nu$ are algebraic integers, hence their integral symmetric functions are rational integers (§. 133).

Since now the number $i\pi$ is to occur among the $\beta$, it follows from (1) that
\[
(1+e^{\beta_1})(1+e^{\beta_2})\cdots(1+e^{\beta_\nu})=0,
\]
and if we carry out the multiplication there results an equation
\[
1+\sum e^{\beta_i}+\sum e^{\beta_i+\beta_h}
+\sum e^{\beta_i+\beta_h+\beta_k}+\cdots=0.
\]
Among the exponents in this sum the number zero can occur several times; we shall assume $(C-1)$ times, so that $C$ is a positive integer, at least $=1$. The remaining exponents $\beta_i,\beta_i+\beta_h,\beta_i+\beta_h+\beta_k,\ldots$, which in part may also be equal among themselves, we shall denote by $\alpha_1,\alpha_2,\ldots,\alpha_\mu$, so that the equation
\begin{equation}\tag{3}
C+e^{\alpha_1}+e^{\alpha_2}+\cdots+e^{\alpha_\mu}=0
\end{equation}
holds. Here the $\alpha_1,\alpha_2,\ldots,\alpha_\mu$ are algebraic numbers which, when multiplied by the rational integer $a$, pass into algebraic integers.

The symmetric functions of all the sums $\beta_i,\beta_i+\beta_h,\beta_i+\beta_h+\beta_k,\ldots$ are at the same time symmetric functions of $\beta_1,\ldots,\beta_\nu$, and hence rational numbers. These sums are consequently also the roots of a rational equation; and since one can separate off the root 0, as often as it is present, the $\alpha_1,\alpha_2,\ldots,\alpha_\mu$ also are the roots of a rational equation; the elementary symmetric functions of $a\alpha_1,a\alpha_2,\ldots,a\alpha_\mu$ are rational integers.

The absolute values of the numbers
\[
\alpha_1,\alpha_2,\ldots,\alpha_\mu
\]
shall now be denoted by
\[
a_1,a_2,\ldots,a_\mu .
\]

If we now put in equation (12) of the preceding paragraph $x=0,\alpha_1,\alpha_2,\ldots,\alpha_\mu$ and apply equation (3), we get
\begin{align}
0={}&CF(0)+F(\alpha_1)+F(\alpha_2)+\cdots+F(\alpha_\mu)\notag\\
&+CQ(0)+e^{\alpha_1}Q(\alpha_1)+e^{\alpha_2}Q(\alpha_2)+\cdots+e^{\alpha_\mu}Q(\alpha_\mu),
\tag{4}
\end{align}
and we have to show the impossibility of such an equation by a suitable choice of $f(x)$.

Let $p$ again be a prime number destined to increase without bound, and put
\begin{equation}\tag{5}
f(x)=
\frac{a^{\mu p+p-1}x^{p-1}(x-\alpha_1)^p(x-\alpha_2)^p\cdots(x-\alpha_\mu)^p}{\Pi(p-1)},
\end{equation}
so that $f(x)$ is a function with rational coefficients.

We now form, as in the preceding paragraph, by ordering according to powers of $x$,
\[
f(x)=\frac{A_{p-1}x^{p-1}+A_px^p+A_{p+1}x^{p+1}\cdots}{\Pi(p-1)},
\]
where the $A_{p-1},A_p,\ldots$ are rational integers, and
\[
A_{p-1}=(-1)^\mu a^{\mu p+p-1}\alpha_1^p\alpha_2^p\cdots\alpha_\mu^p.
\]
If then we take $p$ greater than each of the rational integers
\[
a,\quad a^\mu\alpha_1\alpha_2\cdots\alpha_\mu,
\]
then $A_{p-1}$ is not divisible by $p$, and
\[
F(0)=A_{p-1}+pA_p+p(p+1)A_{p+1}+\cdots
\]
is obtained, that is, an integer not divisible by $p$. Likewise, if $p$ is taken large enough, $C$ is not divisible by $p$.

On the other hand, we order the numerator of $f(x)$ according to powers of $a(x-\alpha_1)$ and obtain
\[
f(x)=\frac{B_pa^p(x-\alpha_1)^p+B_{p+1}a^{p+1}(x-\alpha_1)^{p+1}+\cdots}{\Pi(p-1)},
\]
where $B_p,B_{p+1},\ldots$ are no longer rational, but certainly algebraic integers, since the numerator of $f(x)$ is an integral function of $ax,a\alpha_1,\ldots,a\alpha_\mu$.

From this it follows, as in the preceding paragraph, that
\[
F(\alpha_1)=pB_pa^p+p(p+1)B_{p+1}a^{p+1}+\cdots.
\]
If one forms in the same way $F(\alpha_2),\ldots,F(\alpha_\mu)$, and observes that the sum $B_p(\alpha_1)+B_p(\alpha_2)+\cdots+B_p(\alpha_\mu)$ is a rational integer, it follows that
\[
F(\alpha_1)+F(\alpha_2)+\cdots+F(\alpha_\mu)
\]
is a rational integer divisible by $p$. From this it finally follows that
\[
CF(0)+F(\alpha_1)+F(\alpha_2)+\cdots+F(\alpha_\mu)
\]
is a rational integer not divisible by $p$, and hence also not vanishing, which therefore, apart from sign, is at least $=1$.

If we can now finally still prove that, under the present assumption about $f$, the function $Q(x)$ can be made arbitrarily small for every finite $x$ by sufficiently increasing $p$, then, exactly as above, the impossibility of equation (4) results.

This, however, can be seen as follows: instead of the function
\begin{equation}\tag{6}
f(x)=c_nx^n+c_{n-1}x^{n-1}+\cdots+c_0
\end{equation}
we consider the function
\[
\psi(x)=
\frac{a^{\mu p+p-1}x^{p-1}(x+a_1)^p(x+a_2)^p\cdots(x+a_\mu)^p}{\Pi(p-1)}
=\gamma_nx^n+\gamma_{n-1}x^{n-1}+\cdots+\gamma_0,
\]
which now has only positive (but not necessarily rational) coefficients. The coefficients $c_n,c_{n-1}\ldots$ are formed by multiplication and addition from the numbers $a,-\alpha_1,-\alpha_2,\ldots,-\alpha_\mu$, and the corresponding coefficients $\gamma_n,\gamma_{n-1}\ldots$ are obtained from them by replacing $-\alpha_1,-\alpha_2,\ldots,-\alpha_\mu$ by their absolute values $a_1,a_2,\ldots,a_\mu$. From this it follows, by the theorem of the introduction already mentioned above, that the coefficients $\gamma_n,\gamma_{n-1},\ldots$ are certainly not smaller than the absolute values of the corresponding $c_n,c_{n-1}\ldots$.

Now for every finite $x$, whose absolute value is $\xi$, the absolute value of $Q(x)$ is, by §. 205, (10), smaller, or at least not greater than
\[
\gamma_n\xi^n+\gamma_{n-1}\xi^{n-1}+\cdots+\gamma_0=\psi(\xi),
\]
and that $\psi(\xi)$ sinks below every bound if $p$ is large enough follows from the representation
\[
\psi(x)=\frac{X}{ax}\frac{X^{p-1}}{\Pi(p-1)},
\]
when
\[
X=a^{\mu+1}x(x+a_1)(x+a_2)\cdots(x+a_\mu)
\]
is put.

Thus it is proved:

\begin{center}
\textbf{The number $\pi$ is a transcendent number. The ``quadrature of the circle'' cannot be solved by a geometric construction in which only algebraic curves and surfaces are used.}
\end{center}

\begin{center}
{\large §. 207.}\\[0.4em]
{\large The General Theorem of Lindemann on the Exponential Function.}
\end{center}

The transcendence of the numbers $e$ and $\pi$, which has hereby been proved, is contained as a special case in a very general theorem on the exponential function, which Lindemann announced in the paper mentioned above, and a fully carried out proof of which is contained in the cited paper of Weierstrass. We shall prove this theorem here at the close, using the same aids that we have used for the two special cases. The theorem reads as follows:

\textbf{I. There exists no equation of the form}
\begin{equation}\tag{1}
C_1e^{z_1}+C_2e^{z_2}+\cdots+C_me^{z_m}=0,
\end{equation}
\textbf{in which the coefficients $C_1,C_2,\ldots,C_m$ are algebraic numbers and the exponents $z_1,z_2,\ldots,z_m$ are mutually different algebraic numbers, unless all coefficients $C_1,C_2,\ldots,C_m$ are equal to zero.}

In order to prove it, we first derive an auxiliary theorem.

Let
\[
\alpha_1,\alpha_2,\ldots,\alpha_r
\]
be arbitrary real or imaginary, but mutually different, quantities, and let
\[
A_1,A_2,\ldots,A_r
\]
likewise be arbitrary quantities, not all vanishing. The same is to hold for the two sequences of quantities
\[
\beta_1,\beta_2,\ldots,\beta_s,\qquad
B_1,B_2,\ldots,B_s.
\]
We denote by
\[
\gamma_1,\gamma_2,\ldots,\gamma_t
\]
the mutually different ones among the $rs$ sums $\alpha_i+\beta_k$, and set
\begin{align}
A&=A_1e^{\alpha_1}+A_2e^{\alpha_2}+\cdots+A_re^{\alpha_r},\tag{2}\\
B&=B_1e^{\beta_1}+B_2e^{\beta_2}+\cdots+B_se^{\beta_s},\notag\\
AB&=C_1e^{\gamma_1}+C_2e^{\gamma_2}+\cdots+C_te^{\gamma_t}.\tag{3}
\end{align}
The auxiliary theorem to be proved is then:

\textbf{1. The coefficients $C_1,C_2,\ldots,C_t$ cannot all vanish.}

In the proof of this theorem we can plainly assume that none of the coefficients $A_1,A_2,\ldots,A_r,B_1,B_2,\ldots,B_s$ vanishes (since any vanishing ones can simply be left out).

For the sake of shorter expression, we shall for the moment say, if $a,b$ are two different complex numbers, that $a$ is smaller than $b$, $(a<b)$, if the real part of $a$ is smaller than the real part of $b$, or if the real parts are equal and the imaginary part of $a$ is smaller than the imaginary part of $b$.

If, then, in this sense $a<b$, $b<c$, then also $a<c$, and if $a<b$, $c<d$, then $a+c<b+d$.

Among every finite sequence of mutually different complex numbers there is then a definite smallest one; and if therefore $\alpha_1$ is the smallest among the numbers $\alpha$, and $\beta_1$ the smallest among the numbers $\beta$, then $\alpha_1+\beta_1$ is the smallest among the numbers $\gamma$, and this sum is equal to none of the other sums $\alpha_i+\beta_k$. Thus $C_1=A_1B_1$, and $C_1$ is different from zero, as was to be proved.

This theorem can now be generalized immediately by complete induction.

\textbf{2. If}
\[
\begin{aligned}
A'&=A'_1e^{\alpha'_1}+A'_2e^{\alpha'_2}+\cdots\\
A''&=A''_1e^{\alpha''_1}+A''_2e^{\alpha''_2}+\cdots\\
A'''&=A'''_1e^{\alpha'''_1}+A'''_2e^{\alpha'''_2}+\cdots\\
&\hspace{3em}\cdots
\end{aligned}
\]
\textbf{are sums of the form (2) in arbitrary number, and $\gamma_1,\gamma_2,\ldots$ are the mutually different ones among the sums $\alpha'_h+\alpha''_i+\alpha'''_k+\cdots$, then in the product}
\begin{equation}\tag{4}
A'A''A'''\cdots=C_1e^{\gamma_1}+C_2e^{\gamma_2}+\cdots
\end{equation}
\textbf{not all coefficients $C_1,C_2,\ldots$ are equal to zero.}

This auxiliary theorem first permits us, in the proof of theorem (1), to make the simplifying assumption that the coefficients $C_1,C_2,\ldots,C_m$ are rational integers.

Namely, if an equation
\begin{equation}\tag{5}
X_1e^{x_1}+X_2e^{x_2}+\cdots=0
\end{equation}
with algebraic coefficients $X_1,X_2,\ldots$ and mutually different exponents $x_1,x_2,\ldots$ exists, then we can always assume that the $X_1,X_2,\ldots$ belong to one and the same algebraic field $\Omega$.

We denote by $u_1,u_2,\ldots$ variables and form the norm of the linear form $X_1u_1+X_2u_2+\cdots$:
\begin{equation}\tag{6}
N(X_1u_1+X_2u_2+\cdots)=\Phi(u_1,u_2,\ldots).
\end{equation}
This is an integral homogeneous function of the variables $u$ with rational coefficients, whose degree is equal to the degree of the field $\Omega$. If in (6) we put
\[
u_1=e^{x_1},\quad u_2=e^{x_2},\ldots,
\]
then every product of variables $u$ passes into a quantity of the form $e^z$, where $z$ is a sum of numbers $x$, and $\Phi(u_1,u_2,\ldots)$ becomes an expression of the form $C_1e^{z_1}+C_2e^{z_2}+\cdots$ whose coefficients are rational numbers; if the $X_1,X_2,\ldots$ do not all vanish, then by 2. these coefficients also cannot all be zero after the terms which are equal among the terms of the function $\Phi(e^{x_1},e^{x_2},\ldots)$ have been collected into a single term $Ce^z$. If, however, equation (5) holds, then also $\Phi(e^{x_1},e^{x_2},\ldots)=0$, and consequently
\[
C_1e^{z_1}+C_2e^{z_2}+\cdots=0.
\]
If among the $C_1,C_2,\ldots$ there are fractional numbers, then by multiplying the whole equation by the principal denominator we can transform them into integers.

\textbf{3. Thus it remains only to prove that equation (1) cannot exist for any system of rational integers $C_1,C_2,\ldots,C_m$, not all of which vanish.}

Accordingly we now assume that there exists an equation
\begin{equation}\tag{7}
A_1e^{x_1}+A_2e^{x_2}+\cdots=0,
\end{equation}
whose coefficients $A_1,A_2,\ldots$ are rational integers,
not all $=0$, and in which the exponents $x_1,x_2,\ldots$ are mutually different algebraic numbers. We denote by $\Omega$ a normal field to which all these numbers $x_1,x_2,\ldots$ belong, and by $\Theta$ a primitive number of this field, further by
\[
\sigma'=(\Theta,\Theta'),\quad \sigma''=(\Theta,\Theta''),\ldots
\]
the substitutions of the field $\Omega$. If by one of these substitutions, $\sigma'$, the numbers $x_1,x_2,\ldots$ pass into $x'_1,x'_2,\ldots$, then these also must be mutually different.

Let $u$ now be a variable and
\begin{equation}\tag{8}
U=A_1e^{ux_1}+A_2e^{ux_2}+\cdots .
\end{equation}
Here we carry out all substitutions $\sigma',\sigma'',\ldots$ and denote the functions thereby arising by $U',U'',\ldots$. The product of all these functions, although it is not merely a question of algebraic numbers, we can call the norm of $U$. This product has the form
\begin{equation}\tag{9}
N(U)=C_1e^{uz_1}+C_2e^{uz_2}+\cdots,
\end{equation}
where the $C_1,C_2,\ldots$ are likewise rational integers. The $z_1,z_2,\ldots$ are numbers of the field $\Omega$, which, if we collect terms with equal exponents into a single term, we may assume to be mutually different. At the same time, because of (7),
\begin{equation}\tag{10}
C_1e^{z_1}+C_2e^{z_2}+\cdots=0.
\end{equation}
But the expression on the right-hand side of (9) still has, as its origin as a norm shows, the property of remaining unchanged when any one of the substitutions $\sigma',\sigma'',\ldots$ is made. If, however, one expands (9) according to powers of $u$, this invariance must hold for each term of this series expansion; hence, if $h$ is any positive integer exponent, for
\[
C_1z_1^h+C_2z_2^h+\cdots.
\]
But this sum is a number of the field $\Omega$, and therefore a rational number. We can collect this as follows:

If $g(z)$ denotes any integral function of $z$ with rational coefficients, then
\[
C_1g(z_1)+C_2g(z_2)+\cdots
\]
is a rational number.

\textbf{4. The proof of theorem I. is hereby reduced to the proof of the following special case: if there exists an equation}
\begin{equation}\tag{11}
C_1e^{z_1}+C_2e^{z_2}+\cdots+C_me^{z_m}=0,
\end{equation}
\textbf{in which the $z_1,z_2,\ldots,z_m$ are mutually different algebraic numbers, in which $C_1,C_2,\ldots,C_m$ and the combinations}
\begin{equation}\tag{12}
C_1g(z_1)+C_2g(z_2)+\cdots+C_mg(z_m)
\end{equation}
\textbf{are rational numbers whenever $g(z)$ is any integral function with rational coefficients, then the $C_1,C_2,\ldots,C_m$ must all vanish.}

To prove this theorem, first determine a positive rational integer $c$ so that
\begin{equation}\tag{13}
x_1=cz_1,\quad x_2=cz_2,\ldots,\quad x_m=cz_m
\end{equation}
become algebraic integers (§. 133, 5.). We assume the coefficients $C_1,C_2,\ldots,C_m$ to be rational integers, and form an integral function
\begin{equation}\tag{14}
\varphi(x)=ax^r+a_1x^{r-1}+\cdots+a_r
\end{equation}
with rational integer coefficients, which has the following properties:
\begin{enumerate}[label=\arabic*)]
\item The numbers $x_1,x_2,\ldots,x_m$ occur among the roots of $\varphi(x)=0$.
\item The sum
\[
C_1\varphi'(x_1)+C_2\varphi'(x_2)+\cdots+C_m\varphi'(x_m)=k,
\]
which by the assumptions (12), (13) is a rational integer, is different from zero.
\end{enumerate}

To see that such a function $\varphi(x)$ always exists, one first takes a function $\chi(x)$ which satisfies only condition 1), and the second condition that none of the numbers $x_1,x_2,\ldots,x_m$ is a double root of $\chi(x)$. Such a function plainly exists [for example, in order to form a function $\chi$, one can free the norm of the product $(x-x_1)(x-x_2)\cdots(x-x_m)$ from common divisors with its derivatives]. If one then sets
\[
\varphi(x)=x^h\chi(x),\quad
\varphi'(x_1)=x_1^h\chi'(x_1),\quad
\varphi'(x_2)=x_2^h\chi'(x_2),\ldots,
\]
then the sum
\[
C_1\varphi'(x_1)+\cdots+C_m\varphi'(x_m)
=
C_1\chi'(x_1)x_1^h+\cdots+C_m\chi'(x_m)x_m^h
\]
certainly cannot vanish for every exponent $h$, since otherwise, contrary to the assumption, $C_1\chi'(x_1)\ldots C_m\chi'(x_m)$ would all have to be equal to zero.

After the existence of a function $\varphi(x)$ as required has thus been established, we apply formula §. 205, (12):
\[
e^xP(x)=F(x)+e^\xi Q(x).
\]
In it we put $x=z_1,z_2,\ldots,z_m$, and denote the absolute values of $z_1,z_2,\ldots,z_m$ by $\xi_1,\xi_2,\ldots,\xi_m$.

Then, by (11), there results
\begin{align}
0={}&C_1F(z_1)+C_2F(z_2)+\cdots+C_mF(z_m)\notag\\
&+C_1e^{\xi_1}Q(z_1)+C_2e^{\xi_2}Q(z_2)+\cdots+C_me^{\xi_m}Q(z_m),
\tag{15}
\end{align}
and, if $f(x)$ is an arbitrary integral function of degree $n$,
\begin{equation}\tag{16}
F(z)=f(z)+f'(z)+\cdots+f^{(n)}(z).
\end{equation}
Putting, with $p$ denoting a natural prime number which can be taken arbitrarily large, $x=cz$ [by (13)] and
\begin{equation}\tag{17}
f(z)=\frac{\varphi(x)^{p-1}\varphi'(x)}{\Pi(p-1)},
\end{equation}
where $\varphi(x)$ is the function satisfying conditions 1), 2), we get the required choice of $f$.

By ordering according to powers of $(x-x_1)$ let there result
\[
\varphi(x)^{p-1}\varphi'(x)=
\varphi'(x_1)^p(x-x_1)^{p-1}
+A_p(x_1)(x-x_1)^p
+A_{p+1}(x_1)(x-x_1)^{p+1}+\cdots,
\]
and in this the $A_p(x_1),A_{p+1}(x_1),\ldots$ are algebraic integers, indeed rational functions of $x_1$. On the other hand, by Taylor's theorem,
\[
f(z)=f(z_1)+(z-z_1)f'(z_1)+\frac{(z-z_1)^2}{1.2}f''(z_1)+\cdots
\]
and $c(z-z_1)=x-x_1$. The comparison then gives
\[
f(z_1)=0,\quad f'(z_1)=0,\ldots,\quad f^{(p-2)}(z_1)=0,
\]
\[
f^{(p-1)}(z_1)=c^{p-1}\varphi'(x_1)^p,\quad
f^{(p)}(z_1)=pc^pA_p(x_1),\quad
f^{(p+1)}(z_1)=p(p+1)c^{p+1}A_{p+1}(x_1),\ldots,
\]
and consequently
\[
F(z_1)=c^{p-1}\varphi'(x_1)^p+pc^pA_p(x_1)+p(p+1)c^{p+1}A_{p+1}(x_1)+\cdots;
\]
here $x_1,z_1$ can be replaced by $x_2,z_2$ or by $x_3,z_3$, and so on.

Accordingly the rational integer
\[
C_1F(z_1)+C_2F(z_2)+\cdots+C_mF(z_m)
\]
is congruent modulo $p$ to
\[
c^{p-1}\left[C_1\varphi'(x_1)^p+C_2\varphi'(x_2)^p+\cdots+C_m\varphi'(x_m)^p\right].
\]
Since now $C_1,C_2,\ldots,C_m$ are rational integers, one has $C_1^p\equiv C_1$, $C_2^p\equiv C_2,\ldots\pmod p$, and by application of the polynomial theorem there results
\[
\begin{aligned}
C_1\varphi'(x_1)^p+C_2\varphi'(x_2)^p+\cdots+C_m\varphi'(x_m)^p
&\equiv
\left[C_1\varphi'(x_1)+C_2\varphi'(x_2)+\cdots+C_m\varphi'(x_m)\right]^p\\
&\equiv k^p \pmod p.
\end{aligned}
\]
Thus we obtain the congruence
\begin{equation}\tag{18}
C_1F(z_1)+C_2F(z_2)+\cdots+C_mF(z_m)\equiv c^{p-1}k^p\pmod p.
\end{equation}
Now the numbers $c,k$ are independent of $p$, and we can therefore assume $p$ so large that it divides neither $c$ nor $k$.

\textbf{5. Then the sum $C_1F(z_1)+C_2F(z_2)+\cdots+C_mF(z_m)$ is a rational integer different from zero, hence in absolute value at least equal to 1.}

We now denote by $\varphi_1(x)$ the integral function which is obtained from $\varphi(x)$ by replacing the negative ones among the coefficients $a,a_1,a_2,\ldots$ by their positive values. The function
\[
f_1(z)=\frac{\varphi_1(x)^{p-1}\varphi_1'(x)}{\Pi(p-1)},
\]
which is obtained according to (17), then has only positive coefficients, and these coefficients are in absolute value certainly not smaller than the corresponding coefficients of $f(z)$, because the coefficients of $f_1(z)$ arise from the same numbers by addition as those which, in forming the coefficients of $f(z)$, are partly added and partly subtracted.

From this, however, by formula §. 205, (10), if we denote by $\xi$ the absolute value of $x$, for the absolute value of $Q(z)$ we get
\[
\abs{Q(z)}\leq \frac{\varphi_1(\xi)^{p-1}\varphi_1'(\xi)}{\Pi(p-1)}.
\]

\textbf{6. and therefore $Q(z)$ can be made arbitrarily small for every finite $z$ by sufficiently increasing $p$.}

By 5. and 6. it is now shown that, if $p$ is large enough, equation (15) cannot hold. Thus 4., and with it the whole theorem I., is proved.

This theorem now allows many applications. It gives us first the transcendence of $e$, if we take $C_1,C_2,\ldots,z_1,z_2,\ldots$ as rational integers. It further gives the transcendence of $\pi$. For from the equation $1+e^{i\pi}=0$ it follows by I. that $i\pi$, and consequently also $\pi$, cannot be algebraic.

It follows further from it:

\begin{center}
\textbf{For every algebraic number $x$, with the exception of $x=0$, $X=e^x$ is a transcendent number.}

\textbf{For every algebraic $X$, with the exception of $X=1$, every natural logarithm $x=\log X$ is a transcendent number.}

\textbf{For every arc which stands to the radius in an algebraically expressible ratio $x$, with the exception of $x=0$, $X=\sin x$ is a transcendent number.}
\end{center}

This follows by I. from $2iX=e^{ix}-e^{-ix}$.

The same holds for the other trigonometric lines, $\cos x$, $\tg x$, and for the chord $\frac12\sin\frac{x}{2}$. To mention one more: the transcendent equation $\tg x=ax$ has, for algebraic $a$ different from 0, only transcendent numbers as roots.

\bigskip
\begin{center}\rule{0.28\linewidth}{0.4pt}\end{center}

\clearpage
\begin{center}
{\Large Addenda.}\\[1em]
{\large I.}\\[0.5em]
{\large Irreducibility of the cyclotomic equation.}
\end{center}

The last part of §. 134 of the first volume, which treats the irreducibility of the cyclotomic equation in the general case in which the degree of the root of unity is a number composed of several different prime numbers, requires a correction, which is to be supplied here; and we ask the reader to complete the cited passage according to what is here carried out.

If one puts $n=ab$, with $a$ and $b$ relatively prime, and if $\rho,\alpha,\beta$ are primitive $n$th, $a$th, $b$th roots of unity, then $\rho=\alpha\beta$. At the cited place it was proved that, if $\rho$ is a root of a rational equation $\Phi(x)=0$, then among the roots of this equation $\rho=\alpha\beta$ all primitive $a$th roots of unity $\alpha$ and all primitive $b$th roots of unity $\beta$ must occur. But it does not follow from this, as was erroneously assumed there, that among them \emph{all} primitive $n$th roots of unity occur, since it was not shown that every $\alpha$ occurs combined with every $\beta$.

To fill this gap, there may first be given here a wholly elementary proof\footnote{Arndt, Crelle's Journal, vol. 56, p. 178 (1859). Lebesgue, Liouville's Journal, 2nd series, vol. 4, p. 105 (1859).}.

If $p$ is a prime number, then all binomial coefficients for the $p$th power
\[
B_h^{(p)}=\frac{p(p-1)\cdots(p-h+1)}{1\cdot2\cdots h},
\]
with the exception of $B_0^{(p)}$ and $B_p^{(p)}$, are integers divisible by $p$,
and consequently, if $u,v$ are indeterminate quantities and the congruence is referred only to the coefficients,
\[
(u+v)^p\equiv u^p+v^p \pmod p.
\]
If in this formula we replace $v$ again by a sum $v+w+\cdots$, we obtain, by complete induction (or also from the polynomial theorem),
\[
(u+v+w+\cdots)^p\equiv u^p+v^p+w^p+\cdots \pmod p,
\]
and, if we apply this formula repeatedly and understand by $k$ an arbitrary positive exponent,
\begin{equation}\tag{1}
(u+v+w+\cdots)^{p^k}\equiv u^{p^k}+v^{p^k}+w^{p^k}+\cdots \pmod p.
\end{equation}

Let $x$ now be a variable, and let $a_1,a_2,\ldots,a_\nu$ be rational integers,
\[
f(x)=x^\nu+a_1x^{\nu-1}+a_2x^{\nu-2}+\cdots+a_\nu
\]
a whole function of $x$ whose roots may be $\alpha,\beta,\gamma,\ldots$, so that
\[
-a_1=\sum\alpha,\qquad a_2=\sum\alpha\beta,\qquad -a_3=\sum\alpha\beta\gamma\cdots
\]
are the elementary symmetric functions of $\alpha,\beta,\gamma,\ldots$.

We form the function
\[
F(x)=x^\nu+A_1x^{\nu-1}+A_2x^{\nu-2}+\cdots+A_\nu,
\]
whose roots are the $p^k$th powers of the roots of $f$, thus
\[
\alpha^{p^k},\ \beta^{p^k},\ \gamma^{p^k},\ldots .
\]
The coefficients of $F$,
\[
A_1=-\sum\alpha^{p^k},\quad
A_2=\sum\alpha^{p^k}\beta^{p^k},\quad
A_3=-\sum\alpha^{p^k}\beta^{p^k}\gamma^{p^k}\cdots,
\]
are, as integral symmetric functions of $\alpha,\beta,\gamma,\ldots$, rationally and integrally expressible by $a_1,a_2,\ldots$, and hence are integers. The quotients
\[
\frac{A_1-a_1^{p^k}}{p},\qquad
\frac{A_2-a_2^{p^k}}{p},\qquad
\frac{A_3-a_3^{p^k}}{p},\ldots
\]
are likewise integers by formula (1), and hence by Fermat's theorem
\begin{equation}\tag{2}
F(x)\equiv f(x) \pmod p.
\end{equation}

These propositions are now to be applied to the function $f_n(x)$ whose roots are all primitive $n$th roots of unity $\alpha,\beta,\gamma,\ldots$,
 and which, as was proved in §. 133 of the first volume, has integral coefficients $a_1,a_2,\ldots$.

If $p$ is a prime number dividing $n$, then
\begin{equation}\tag{3}
\frac{x^n-1}{x^{\frac np}-1}
=x^{\frac np(p-1)}+x^{\frac np(p-2)}+\cdots+x^{\frac np}+1,
\end{equation}
and this function vanishes for $x=\alpha,\beta,\gamma,\ldots$, and is therefore divisible by $f_n(x)$. We shall put it equal to $f_n(x)g(x)$, where $g(x)$ is then (by vol. I, §. 2) a whole function of $x$ with integral coefficients.

If now
\[
n=p^k n',
\]
where $n'$ is no longer divisible by $p$, and if $\alpha'$ is a primitive $n'$th root of unity, then
\[
\alpha'^{\frac np}=1,
\]
and consequently, if one puts $x=\alpha'$ [by (3)],
\begin{equation}\tag{4}
f_n(\alpha')g(\alpha')=p.
\end{equation}

From this we first derive a proof of the irreducibility of $f_n(x)$ under the assumption that $n$ is a power of a prime, although for this case the proof given in vol. I, §. 134 is unobjectionable.

Assume that $f_n(x)$ decomposes into two rational factors:
\begin{equation}\tag{5}
f_n(x)=\varphi(x)\psi(x).
\end{equation}
Then all roots, both of $\varphi(x)$ and of $\psi(x)$, are $n$th roots of unity, and their $n$th powers are therefore equal to 1. If now $n=p^k$, then from $\varphi(x)$ and $\psi(x)$ we can derive two whole functions $\Phi(x)$ and $\Psi(x)$ whose roots are the $n$th powers of the roots of $\varphi(x)$ and $\psi(x)$, and therefore are all equal to 1; and from (2) one obtains
\[
\varphi(x)\equiv\Phi(x),\qquad \psi(x)\equiv\Psi(x) \pmod p,
\]
whence for $x=1$ it follows that
\[
\varphi(1)\equiv0,
\qquad
\psi(1)\equiv0 \pmod p.
\]
From this (5) gives
\begin{equation}\tag{6}
f_n(1)\equiv0 \pmod {p^2}.
\end{equation}
But in the present case $f_n(1)=p$ (vol. I, §. 133, V.), which contradicts formula (6). Our assumption was therefore inadmissible, and $f_n(x)$ is irreducible.

To prove the irreducibility for an arbitrary composite $n$, we apply complete induction. We put
\[
n=p^k n',
\]
and understand by $p$ a prime number which divides $n$ but not $n'$, and assume the irreducibility of $f_{n'}(x)$ as already proved. If $f_n(x)$ is then decomposable into two rational factors $\varphi(x)$, $\psi(x)$, neither of which is constant,
\begin{equation}\tag{7}
f_n(x)=\varphi(x)\psi(x),
\end{equation}
then from $\varphi(x)$ we derive the function $\Phi(x)$ whose roots are the $p^k$th powers of the roots of $\varphi(x)$, and which satisfies the condition
\[
\varphi(x)\equiv\Phi(x) \pmod p,
\]
or, if we denote by $\chi(x)$ a second integral function of $x$,
\[
\varphi(x)=\Phi(x)+p\chi(x).
\]
The roots of $\Phi(x)$, however, are primitive roots of unity of degree $n'$, and therefore among them there must be at least one, $\alpha'$, for which $\Phi(\alpha')$ vanishes, so that
\[
\varphi(\alpha')=p\chi(\alpha').
\]
Because of the assumed irreducibility of $f_{n'}(x)$, however, this equation must hold for \emph{all} roots of unity $\alpha'$; and for the same reason, if $\vartheta(x)$ is an integral function, for all $\alpha'$ one has
\[
\psi(\alpha')=p\vartheta(\alpha').
\]
Accordingly (7) gives
\[
f_n(\alpha')=p^2\chi(\alpha')\vartheta(\alpha'),
\]
and by (4),
\[
1=p\chi(\alpha')\vartheta(\alpha')g(\alpha').
\]
But the product $\chi(x)\vartheta(x)g(x)$ is a whole integral function of $x$, whose remainder on division by $f_{n'}(x)$ likewise has integral coefficients and may be denoted by $\Theta(x)$. This remainder is of lower degree than $f_{n'}(x)$, and the function
\[
p\Theta(x)-1
\]
vanishes for all roots of $f_{n'}(x)$. Consequently the equation
\[
p\Theta(x)=1
\]
must be identical, which is plainly impossible, since $\Theta(x)$, for an integral $x$, passes into an integer which, multiplied by $p$, cannot be equal to 1.

Thus the irreducibility of $f_n(x)$ is proved in general.

\medskip\hrule\medskip

After the irreducibility of the cyclotomic equation $f_n(x)=0$ has thus been proved in general, it follows, as in §. 134 of the first volume, that it remains irreducible when an arbitrary root of unity is adjoined whose degree is relatively prime to $n$.

Kronecker goes, in a certain sense, the opposite way\footnote{Mém. sur les facteurs irréduct. de l'expression $(x^n-1)$. Liouville's Journal, vol. 19 (1854). A simple proof of the irreducibility of the general cyclotomic equation based on the theory of higher congruences was given by Dedekind [Crelle's Journal für Mathematik, vol. 58 (1859)].}. He first proves, under the assumption that $n$ is a power of a prime $p$, the irreducibility of $f_n(x)$ after adjunction of a root of unity whose degree is not divisible by $p$, and derives from this the general proof of irreducibility, as we shall now briefly show.

\medskip\hrule\medskip

Let now
\[
n=p^k
\]
be a power of a prime number $p$, and let
\begin{equation}\tag{8}
f_n(x)=x^{p^{k-1}(p-1)}+x^{p^{k-1}(p-2)}+\cdots+x^{p^{k-1}}+1
\end{equation}
be the function whose roots are the primitive $n$th roots of unity. Further let $\alpha$ be any root of unity whose degree $m$ is not divisible by $p$, and let $\nu$ be the degree of the field $R(\alpha)$. It is to be proved that $f_n(x)$ is irreducible in the field $R(\alpha)$.

Every number of the field $R(\alpha)$ can be represented in one and only one way in the form
\begin{equation}\tag{9}
\varphi(\alpha)=m_0+m_1\alpha+\cdots+m_{\nu-1}\alpha^{\nu-1},
\end{equation}
where $m_0,m_1,\ldots,m_{\nu-1}$ are rational numerical coefficients.

If $\chi(x)=0$ is the rational equation of least degree whose root is $\alpha$, and
\begin{equation}\tag{10}
\chi(x)=x^\nu+c_1x^{\nu-1}+\cdots+c_{\nu-1}x+c_\nu,
\end{equation}
then $c_1,c_2,\ldots$ are integers, since $\chi(x)$ must be a divisor of $x^m-1$ (by Gauss's theorem, vol. I, §. 2).

From this it follows that, if $\varphi(x)$ denotes a whole function of $x$ with integral rational coefficients, then, even if the degree of $\varphi(x)$ was originally higher than $\nu$, the coefficients $m_0,m_1,\ldots,m_{\nu-1}$ in (9) become integers; for these coefficients are obtained when one seeks the remainder of the division of $\varphi(x)$ by $\chi(x)$.

We now assume that $f_n(x)$ is decomposable in the field $R(\alpha)$, and thus that
\begin{equation}\tag{11}
f_n(x)=\varphi(x,\alpha)\psi(x,\alpha),
\end{equation}
where $\varphi(x,\alpha)$ and $\psi(x,\alpha)$ are two whole functions of $x$, neither of which is constant, and whose coefficients are contained in $R(\alpha)$. From (8) and (11), if we put $x=1$, we obtain
\begin{equation}\tag{12}
p=\varphi(1,\alpha)\psi(1,\alpha),
\end{equation}
where $\varphi(1,\alpha)$ and $\psi(1,\alpha)$ are representable, as numbers of the field $R(\alpha)$, in the form (9). If we denote the least common denominators of these two representations by $a$ and $b$, we get
\begin{equation}\tag{13}
\begin{aligned}
a\varphi(1,\alpha)&=a_0+a_1\alpha+a_2\alpha^2+\cdots+a_{\nu-1}\alpha^{\nu-1}=A(\alpha),\\
b\psi(1,\alpha)&=b_0+b_1\alpha+b_2\alpha^2+\cdots+b_{\nu-1}\alpha^{\nu-1}=B(\alpha),
\end{aligned}
\end{equation}
where $a,a_0,a_1,\ldots,a_{\nu-1}$ are integers without a common divisor, and $A$ is a function-symbol. The same holds for $b,b_0,b_1,\ldots,b_{\nu-1}$ and $B$.

The roots of the equation $\varphi(x,\alpha)=0$ are found among the primitive $n$th roots of unity, and these are all powers of one among them, $r$. Accordingly there are certain powers $r^h,r^{h'},r^{h''},\ldots$ such that
\[
\varphi(x,\alpha)=(x-r^h)(x-r^{h'})(x-r^{h''})\cdots,
\]
and hence, by (13),
\[
A(\alpha)=a(1-r^h)(1-r^{h'})(1-r^{h''})\cdots .
\]
If we take the $n$th power of this, then, by (1), taking into account the equation $r^n=1$, we obtain
\[
[A(\alpha)]^n=p\chi(r),
\]
where $\chi(x)$ denotes a whole function of the variable $x$ with integral coefficients.

Now, understanding by $t$ a variable, we set
\[
[t-\chi(1)][t-\chi(r)]\cdots[t-\chi(r^{n-1})]
=t^n+A_1t^{n-1}+A_2t^{n-2}+\cdots+A_n,
\]
where the coefficients $A_1,A_2,\ldots,A_n$, as whole symmetric functions of the $n$ quantities $\chi(1),\chi(r),\ldots,\chi(r^{n-1})$, are rational integers. But if in this equation we put
\[
t=\chi(r)=\frac{[A(\alpha)]^n}{p},
\]
then the left hand side vanishes, and multiplication by $p^n$ gives
\begin{equation}\tag{14}
[A(\alpha)]^{n^2}=pC(\alpha),
\end{equation}
where $C(\alpha)$ denotes a number of the form (9) with integral coefficients.

If formula (14) is repeatedly raised to the $p$th power and proposition (1) is applied, then for every power not smaller than $n^2$ one obtains
\begin{equation}\tag{15}
A(\alpha^{p^\lambda})=p\xi(\alpha),
\end{equation}
where the whole function $\xi(\alpha)$ likewise has integral coefficients.

Here we may assume $\lambda$ divisible by $\varphi(m)$; then, since $m$ is not divisible by $p$, by the generalized theorem of Fermat (vol. II, §. 14) one has
\[
p^\lambda\equiv1 \pmod m,
\qquad
\alpha^{p^\lambda}=\alpha,
\]
so that from (14) it follows that
\begin{equation}\tag{16}
A(\alpha)=p\xi(\alpha).
\end{equation}
Thus in (13) the coefficients $a_0,a_1,\ldots,a_{\nu-1}$ are all divisible by $p$, and consequently $a$ is not divisible by $p$. In the same way it is shown that $b_0,b_1,\ldots,b_{\nu-1}$ are divisible by $p$, while $b$ is indivisible by $p$.

Consequently from (12), if one puts $A(\alpha)B(\alpha)=p^2\Theta(\alpha)$, one obtains
\begin{equation}\tag{17}
ab=p\Theta(\alpha),
\end{equation}
where $\Theta(\alpha)$ denotes a function of $\alpha$ of the form (9) with integral coefficients. But the equation (17), since it is of at most the $(\nu-1)$st degree, would have to be identical with respect to $\alpha$, and $ab$ would have to be divisible by $p$, which is impossible.

It follows from this that our assumption that $f_n(x)$ is decomposable in the field $R(\alpha)$ into two factors according to formula (11) is inadmissible.

\medskip\hrule\medskip

This now again gives a general proof that the equation whose roots are all primitive $m$th roots of unity, and whose degree is $\varphi(m)$, is irreducible even when $m$ is an arbitrary composite number. This is equivalent to the following proposition:

\begin{quote}
If a whole function $\varphi(x)$ with rational coefficients vanishes when a primitive $m$th root of unity is substituted for $x$, then $\varphi(x)$ also vanishes when any other primitive $m$th root of unity is substituted for $x$.
\end{quote}

If $m$ is a power of a prime, then the proposition has already been proved above. We therefore apply complete induction, assume it proved for $m$th roots of unity, and derive it for $mn$th roots of unity, where $n=p^k$ and $p$ is a prime not dividing $m$.

Every $mn$th root of unity $\beta$ can be represented in the form
\[
\beta=\alpha r,
\]
where $\alpha$ is an $m$th and $r$ an $n$th root of unity. What is to be proved can be expressed as follows: if $\beta$ is the root of a rational equation $\Phi(x)=1$, then all other primitive roots of unity of degree $mn$ must also be roots of this equation.

If, as above, $f_n(x)=0$ is the equation whose roots are the quantities $r$, then, as we have proved, $f_n(x)$ is not decomposable into rational factors, nor into such factors as depend rationally on $\alpha$. The function $f_n(\alpha^{-1}x)$, which vanishes for $x=\beta$, likewise cannot be decomposed into rational factors; for from a decomposition of the form
\[
f_n(\alpha^{-1}x)=\varphi(x,\alpha)\psi(x,\alpha)
\]
one would obtain
\[
f_n(x)=\varphi(\alpha x,\alpha)\psi(\alpha x,\alpha),
\]
which cannot hold. Since accordingly $f_n(\alpha^{-1}x)$ has a root
in common with $\Phi(x)$, $\Phi(x)$ must be divisible by $f_n(\alpha^{-1}x)$, and we can therefore, for an indeterminate $x$, put
\begin{equation}\tag{18}
\Phi(x)=f_n(\alpha^{-1}x)\Psi(x,\alpha),
\end{equation}
where $\Psi(x,\alpha)$ is a whole rational function of $x$.

Let $\alpha,\alpha_1,\alpha_2,\ldots$ be all the primitive $m$th roots of unity, which, as we have assumed, are roots of an irreducible rational equation
\[
F_m(x)=0.
\]
Then, by this assumption, in the identity (18) $\alpha$ can be replaced by every other $\alpha_1,\alpha_2,\ldots$, and it follows that $\Phi(x)$ is divisible by each of the factors
\[
f_n(\alpha^{-1}x),\quad f_n(\alpha_1^{-1}x),\quad f_n(\alpha_2^{-1}x)\cdots .
\]
But no two of these factors can have a common divisor; for if, for instance,
\[
f_n(\alpha^{-1}x)=0,
\qquad
f_n(\alpha_1^{-1}x)=0
\]
had a common root $x$, then there would have to exist two primitive $n$th roots of unity $r,r_1$ such that
\[
x=\alpha r=\alpha_1r_1.
\]
Thus one would have $\alpha^n=\alpha_1^n$, and consequently, since $n$ is relatively prime to $m$, $\alpha=\alpha_1$.

Accordingly $\Phi(x)$ is divisible by the product
\[
F_m^{(n)}(x)=f_n(\alpha^{-1}x)f_n(\alpha_1^{-1}x)f_n(\alpha_2^{-1}x)\cdots,
\]
whose roots are all the $\beta$; and therefore this function is also indecomposable, as was to be proved.

\begin{center}
{\large II.}\\[0.5em]
{\large The irreducibility of the cyclotomic equation\\
and the theorem on the prime numbers contained in a linear form.}
\end{center}

Whereas in the first addendum, which is to be regarded as a supplement to the first volume, no use was deliberately made of the theory of whole algebraic numbers, we give in conclusion a proof of the irreducibility of the cyclotomic equation that rests on other foundations, namely on the theorems on the class number (§. 192), and that is noteworthy because of the generalizations which it permits.

Let $m$ be any natural number, and let $n$ be the system of positive numbers prime to $m$; among these there are
\begin{equation}\tag{1}
\varphi(m)=\mu
\end{equation}
which are smaller than $m$; we denote them by $a$.

If we unite all numbers $n$ congruent modulo $m$ to one such $a$ into a class $A$, we obtain $\mu$ number-classes
\begin{equation}\tag{2}
A_1,A_2,\ldots,A_\mu,
\end{equation}
which, under composition by multiplication, form an Abelian group $\mathfrak R$ of degree $\mu$. This group was already considered in §. 16 of this volume.

The $\mu$ characters of this group we denote by
\begin{equation}\tag{3}
\chi_1,\chi_2,\ldots,\chi_\mu,
\end{equation}
and put, if $\chi$ is one of these characters and $n$ denotes a number from the class $A$,
\begin{equation}\tag{4}
\chi(A)=\chi(n).
\end{equation}
These characters, which were determined more closely in §. 16, are all $\mu$th roots of unity.

Here, however, we do not need to make use of the special expressions of these characters. Rather, we rely on the following theorem, which holds for all Abelian groups:

\begin{quote}
\textbf{1.} If $A$ is an element of $f$th degree of an Abelian group of $\mu$th degree, and if $\mu=ef$, then all $\chi_i(A)$ are $f$th roots of unity, and among them each $f$th root of unity occurs exactly $e$ times.
\end{quote}

The theorem is contained implicitly in theorem 7, §. 12 of this volume. For if we apply this theorem to the group
\[
T=1,A,A^2,\ldots,A^{f-1},
\]
whose index with respect to $\mathfrak R$ is equal to $e$, then it follows that there is a group $\mathfrak F$ of $e$ characters $\xi$ which satisfy the conditions
\[
\xi(A)=1.
\]
If, therefore, the group $X$ of the characters $\chi$ is decomposed into the cosets
\[
\mathfrak F_1,\mathfrak F_2,\ldots,\mathfrak F_f,
\]
then $\chi_1(A)$ and $\chi_2(A)$ are equal to one another or different from one another according as $\chi_1$ and $\chi_2$ occur in the same or in different ones of these cosets. Moreover, since, on account of $\chi(A)^f=\chi(A^f)=1$, all $\chi(A)$ are roots of unity of degree $f$, and there are only $f$ different such roots, our theorem is thereby proved.

If $n$ is contained in the class $A$, then the degree $f$ of $A$ is the exponent to which $n$ belongs modulo $m$, that is, the least positive exponent for which
\begin{equation}\tag{5}
n^f\equiv1 \pmod m
\end{equation}
holds. If therefore $n$ belongs to the exponent $f$, then all $\chi_i(n)$ are $f$th roots of unity, and each $f$th root of unity occurs among them $e$ times.

Let $a$ denote any one of the $\varphi(m)$ residues of the numbers $n$. Then, in the form
\begin{equation}\tag{6}
\mu_x=mx+a-m
\end{equation}
every number of the number-class (2) represented by $a$ is contained; if by $a$ we understand the least positive residue, then we obtain the totality of the numbers of this class by letting $x$ run through all positive integral values.

Now
\[
\frac{\mu_x}{m}-x=\frac{a-m}{m},
\]
and therefore we can apply the theorem §. 191, 6. to the quantities $\mu_x$, where we have to put $n=x$, $\gamma=1:m$, $c_n=(a-m):m$. Thus
\begin{equation}\tag{7}
A(s)=\sum_{x=0}^{\infty}\frac{1}{(mx+a)^s}
     =\sum_{x=1}^{\infty}\frac{1}{(mx+a-m)^s},
\end{equation}
as long as $s>1$, is a continuous function of $s$, and
\begin{equation}\tag{8}
A(s)=\frac{1}{m(s-1)}+C(s),
\end{equation}
where $C(s)$ remains finite for $s=1$\footnote{By using known theorems on $\Gamma$-functions, one finds from §. 191, (20):
\[
C(1)=\frac1m\,\frac{\Gamma'\!\left(\frac am\right)}{\Gamma\!\left(\frac am\right)}.
\]
}.

Now we let $A$ run through all classes (2), denote by $\chi$ any one of the characters of the group $\mathfrak R$, and put
\begin{equation}\tag{9}
Q(s)=\sum_A \chi(A)A(s),
\end{equation}
or, what is the same,
\begin{equation}\tag{10}
Q(s)=\sum_n \frac{\chi(n)}{n^s},
\end{equation}
where the sum extends over all numbers $n$. The number of these sums $Q$ is as large as the number of characters, that is, $=\varphi(m)$. If a distinction is necessary, we denote them by
\[
Q_1,Q_2,\ldots,Q_\mu.
\]
One of them, $Q_1$, corresponds to the principal character and has the expression
\[
Q_1(s)=\sum_n\frac1{n^s}.
\]
Now [§. 11, (21)]
\[
\sum_A \chi(A)=\mu\quad \hbox{or}=0,
\]
according as $\chi$ is the principal character or not. Furthermore, from (8) one obtains

\[
Q(s)=\frac{1}{m(s-1)}\sum_A \chi(A)+\sum_A\chi(A)C(s),
\]
and hence the theorem:

\begin{quote}
\textbf{2.} The sums $Q_2,\ldots,Q_\mu$ receive finite values for $s=1$; $Q_1$ becomes infinite, and indeed
\[
\lim_{s=1}(s-1)Q_1(s)=\frac{\mu}{m}.
\]
\end{quote}

To transform the sums $Q(s)$, we denote by $p$ each prime number not dividing $m$ and multiply together all infinite series of the form
\[
1+\frac{\chi(p)}{p^s}+\frac{\chi(p^2)}{p^{2s}}+\cdots
=\frac{1}{1-\chi(p)p^{-s}}.
\]
In this way we obtain [cf. §. 192, (8)]
\begin{equation}\tag{11}
Q(s)=\prod_p \frac{1}{1-\chi(p)p^{-s}}.
\end{equation}
If now $p$ belongs to the exponent $f$, so that $f$ is the least positive number satisfying the congruence
\begin{equation}\tag{12}
p^f\equiv1 \pmod m,
\end{equation}
and if $\mu=ef$, then among the $\chi(p)$ each $f$th root of unity occurs exactly $e$ times; consequently, for a variable $x$,
\[
\prod_{\chi}[1-\chi(p)x]=(1-x^f)^e,
\]
and therefore from (11) one obtains
\begin{equation}\tag{13}
Q_1Q_2\cdots Q_\mu=\prod_p \frac{1}{(1-p^{-fs})^e}.
\end{equation}

We now introduce the cyclotomic field $\Omega_m$, without presupposing the irreducibility of the cyclotomic equation, and denote its still unknown degree by $\nu$. Since, however, the primitive $m$th root of unity $r$ satisfies a rational equation of degree $\mu$, one has
\begin{equation}\tag{14}
\nu\leq\mu.
\end{equation}
The numbers $\omega$ of the field $\Omega_m$ are all representable in the form
\begin{equation}\tag{15}
a\omega=a_0+a_1r+\cdots+a_{\nu-1}r^{\nu-1}
\end{equation}
with rational integer coefficients $a,a_0,a_1,\ldots,a_{\nu-1}$,
and if $\omega$ is a whole number, then for $a$ we can put a fixed integer, whose closer knowledge is not required. Its value is equal to the square root of the quotient of the discriminants
\[
\Delta(1,r,r^2,\ldots,r^{\nu-1}):\Delta(\omega_1,\omega_2,\ldots,\omega_\nu),
\]
where $\omega_1,\omega_2,\ldots,\omega_\nu$ is a minimal basis (§. 144, 145).

We exclude from the system of the $p$ all prime numbers occurring in the discriminant of the cyclotomic equation and all occurring in $m$ and in $a$; their number is in any case finite. Let $\mathfrak p$ denote a prime ideal occurring in $p$. Then the congruence
\[
r^p\equiv r \pmod{\mathfrak p}
\]
is fulfilled precisely when
\begin{equation}\tag{16}
p\equiv1 \pmod m,
\end{equation}
that is, when $p$ belongs to the exponent 1, hence when $r^p=r$; under the same condition, by (15), one also has, for every whole number $\omega$ in $\Omega_m$,
\[
\omega^p\equiv\omega \pmod{\mathfrak p}.
\]
But by §. 162 this is the necessary and sufficient condition that the prime number $p$, not occurring in the field-discriminant of $\Omega_m$, decomposes into prime ideals of the first degree only.

Thus, if we denote by
\[
P_f(s)=\prod (1-p^{-sf})
\]
the product extending over all prime numbers $p$ belonging to the exponent $f$, it follows from §. 192 that, if $f>1$, $P_f(1)$ is finite and different from zero, and that
\begin{equation}\tag{17}
\frac{s-1}{P_1(s)^\nu}
\end{equation}
for $s=1$ is finite and different from zero.

Now, by (13), if $S$ denotes the finite product $\prod(1-p^{-sf})^{-e}$ extending over the excluded prime numbers $p$,
\[
Q_1Q_2\cdots Q_\mu=S\prod_f\frac{1}{P_f(s)^e},
\]
and, since $e=\mu$ for $f=1$, it follows from (17) that
\begin{equation}\tag{18}
(s-1)^{\frac{\mu}{\nu}}Q_1Q_2\cdots Q_\mu
\end{equation}
for $s=1$ is finite and different from zero.

On the other hand, by 2.,
\begin{equation}\tag{19}
(s-1)Q_1Q_2\cdots Q_\mu
\end{equation}
is finite for $s=1$, and if we therefore divide (19) by (18), it follows that
\[
(s-1)^{1-\frac{\mu}{\nu}}\quad\hbox{is finite for }s=1,
\]
that is,
\begin{equation}\tag{20}
\nu\geq\mu.
\end{equation}
Together with (14), this gives $\nu=\mu$, whereby the following theorem is proved:

\begin{quote}
\textbf{3.} The degree of the cyclotomic field $\Omega_m$ is equal to $\varphi(m)$; hence the cyclotomic equation of degree $\varphi(m)$, whose roots are the primitive $m$th roots of unity, is irreducible.
\end{quote}

This consideration, however, also renders us another important service, by furnishing a proof of the famous theorem that \emph{in every arithmetic progression whose difference and initial term are natural numbers without common divisor, infinitely many prime numbers are contained}\footnote{The first complete and general proof of this theorem was given by Dirichlet (Abhandlungen der Berliner Akademie for the year 1837; collected works no. XXI). Dedekind drew attention to the essential simplification of this proof by use of Kummer's formulas for the class number in cyclotomic fields (Vorlesungen über Zahlentheorie, 3rd edition, p. 596).}.

For it now follows from (18) that the product
\[
(s-1)Q_1Q_2\cdots Q_\mu
\]
has, for $s=1$, a finite and non-zero value; and since by 2. none of the factors is infinite, it follows that
\begin{quote}
\textbf{4.} The sums
\[
(s-1)Q_1(s),\quad Q_2(s),\ldots,Q_\mu(s)
\]
have finite and non-zero values for $s=1$.
\end{quote}

If we apply to all factors of (11) the formula
\[
\log[1-\chi(p)p^{-s}]
=\frac{\chi(p)}{p^s}+\frac12\frac{\chi(p^2)}{p^{2s}}+\frac13\frac{\chi(p^3)}{p^{3s}}+\cdots
\]
(where the imaginary part of the logarithm is to be taken between $\pm\pi i$), then, if the imaginary part of $\log Q(s)$ is suitably determined, it follows that
\begin{equation}\tag{21}
\log Q(s)=\sum\frac{\chi(p)}{p^s}+\frac12\sum\frac{\chi(p^2)}{p^{2s}}+\frac13\sum\frac{\chi(p^3)}{p^{3s}}+\cdots .
\end{equation}
If now $A$ denotes any one of the classes (2), and $a$ a number from $A$, then by §. 11, (6) of this volume
\begin{equation}\tag{22}
\sum_\chi \chi(A)\chi(n)=\sum_\chi \chi(an)=\mu\quad\hbox{or}=0,
\end{equation}
according as $an\equiv1$ or not $\equiv1$ modulo $m$, that is, according as $n$ is contained in the class $A^{-1}$ or not. If, therefore, we multiply formula (21) by $\chi(A)$ and sum with respect to $\chi$, then it follows that
\begin{equation}\tag{23}
\frac1\mu\sum_\chi \chi(A)\log Q(s)=
\sum\frac1{p^s}+\frac12\sum\frac1{p^{2s}}+\frac13\sum\frac1{p^{3s}}+\cdots,
\end{equation}
where on the right the first, second, third, $\ldots$ sum refers to all prime numbers $p$ whose first, second, third, $\ldots$ power is contained in $A^{-1}$.

The second part of this sum,
\[
R(s)=\frac12\sum\frac1{p^{2s}}+\frac13\sum\frac1{p^{3s}}+\cdots,
\]
is enlarged if each of its partial sums is extended over \emph{all} prime numbers $p$, and consequently
\[
R(s)<\frac12\sum\left(\frac1{p^{2s}}+\frac1{p^{3s}}+\cdots\right)
=\frac12\sum\frac1{p^s(p^s-1)},
\]
and therefore still more
\[
R(s)<\frac12\sum_{n=2}^{\infty}\frac1{n^{2s}}\frac1{1-n^{-s}},
\]
or, since $1-n^{-s}>\frac12$,
\[
R(s)<\sum_{n=2}^{\infty}\frac1{n^{2s}}.
\]
From this one concludes that $R$ remains finite for $s=1$ (§. 191, 2.).

If now in formula (23),
\[
\frac1\mu\sum_\chi \chi(A)\log Q(s)=\sum\frac1{p^s}+R(s),
\]
we let $s$ pass to 1, then $\log Q_1(s)$ becomes infinite by 4., while the remaining terms of the left hand side, $\log Q_2(s),\ldots,\log Q_\mu(s)$, remain finite. Since $R(s)$ likewise remains finite, the sum $\sum p^{-s}$ must become infinite; and this is certainly possible only if the sum consists of infinitely many terms. Hence it is proved that:

\begin{quote}
\textbf{5.} In each of the number-classes $A$ (modulo $m$) infinitely many prime numbers are contained.
\end{quote}

This can also be expressed by saying that the linear form $mx+a$, in which $m$ and $a$ are integers without a common divisor, represents a prime number for infinitely many integral values of $x$.

\clearpage
\section*{Register and Corrections}
\begin{center}\small The Roman numerals denote the volume, the Arabic numerals the page.\end{center}
\begin{multicols}{2}\small\raggedright
\idxitem{Abelian equations I, 533.}
\idxitem{— cubic II, 107.}
\idxitem{— biquadratic II, 110.}
\idxitem{— groups I, 476, 536; II, 6, 32.}
\idxitem{— fields II, 588, 648.}
\idxitem{Derived functions I, 52.}
\idxitem{Absolute norm II, 502, 536.}
\idxitem{Absolute value of an imaginary number I, 18.}
\idxitem{— of a functional II, 501.}
\idxitem{Countable sets II, 745.}
\idxitem{Addition I, 14.}
\idxitem{Adjunction I, 451.}
\idxitem{Similar substitutions II, 154.}
\idxitem{Equivalent functionals II, 552.}
\idxitem{— numbers I, 371.}
\idxitem{— numbers and functionals in the quadratic field II, 610.}
\idxitem{Affect I, 480.}
\idxitem{Algebraic functions II, 526.}
\idxitem{— fields I, 455; II, 527.}
\idxitem{— numbers II, 489.}
\idxitem{Algorithm of the greatest common divisor for numbers I, 2.}
\idxitem{— for integral functions I, 34.}
\idxitem{Alternating group I, 496, 600.}
\idxitem{— on five letters II, 138, 280.}
\idxitem{— function I, 496.}
\idxitem{Arithmetic progressions I, 47.}
\idxitem{Aronhold systems of bitangents II, 367.}
\idxitem{Associated functionals II, 510.}
\idxitem{Associated numbers I, 586.}
\idxitem{Solution of equations by radicals I, 595.}
\idxitem{Axes of a ternary substitution group II, 438.}
\idxitem{Bases of the functionals II, 532, 581.}
\idxitem{— of the ideals II, 550.}
\idxitem{Basis of an Abelian group II, 33.}
\idxitem{— of an algebraic field II, 528.}
\idxitem{— of the natural logarithm system II, 752.}
\idxitem{Basis form of a functional II, 535.}
\idxitem{— of zero II, 535.}
\idxitem{— in the quadratic field II, 604, 608.}
\idxitem{Removal of the second term of an equation I, 114.}
\idxitem{Bernoulli's approximation method for solving equations I, 341.}
\idxitem{Bertrand's theorem on bounds for the index of permutation groups II, 143.}
\idxitem{Bezout's theorem I, 161.}
\idxitem{Bezoutiant I, 225, 265.}
\idxitem{— and reality of roots I, 252.}
\idxitem{Binomial theorem I, 42.}
\idxitem{Binomial coefficients I, 42.}
\idxitem{Binary forms I, 189.}
\idxitem{— linear substitution II, 191.}
\idxitem{Biquadratic Abelian equations II, 110.}
\idxitem{Biquadratic forms I, 199.}
\idxitem{— equations I, 119, 201, 528; II, 319.}
\idxitem{— cyclotomic fields II, 101.}
\idxitem{Bring-Jerrard form of the equation of the fifth degree I, 176, 233.}
\idxitem{Brioschi normal form of the equation of the fifth degree I, 236.}
\idxitem{Fractions I, 10.}
\idxitem{Literal calculation I, 20.}
\idxitem{Budan-Fourier theorem I, 299.}
\idxitem{Canonical form of the ternary cubic form II, 332.}
\idxitem{Cardan formula I, 116.}
\idxitem{Descartes' (Harriot's) theorem on the number of real roots I, 308.}
\idxitem{Casus irreducibilis of cubic equations I, 608.}
\idxitem{Cayley's expression of Cardan's formula I, 118.}
\idxitem{Characters of an Abelian group II, 44.}
\idxitem{Characteristic of a system of functions I, 285.}
\idxitem{Classes of functionals II, 553.}
\idxitem{— ideals II, 553.}
\idxitem{— quadratic forms I, 400.}
\idxitem{— quadratic irrational numbers I, 383.}
\idxitem{Class number of an algebraic field II, 554.}
\idxitem{— in the field of eighth roots of unity II, 717.}
\idxitem{Class number formula II, 702.}
\idxitem{— in the cyclotomic field II, 705.}
\idxitem{— in the real cyclotomic field II, 711.}
\idxitem{— in the field of the \(2^n\)-th roots of unity II, 719.}
\idxitem{Class-number factors II, 716.}
\idxitem{Class-number factor \(A\) II, 722.}
\idxitem{— \(B\) II, 726.}
\idxitem{Commutative groups I, 476; II, 6.}
\idxitem{Commutative normal divisors of a metacyclic group II, 27.}
\idxitem{Complexes of bitangents II, 357.}
\idxitem{Complex roots (their enclosure) I, 292.}
\idxitem{— numbers I, 18.}
\idxitem{— Gaussian I, 585.}
\idxitem{Complex pairs, complex triples II, 363.}
\idxitem{Composition in a group II, 3.}
\idxitem{— of parts II, 12.}
\idxitem{— of quadratic irrational numbers II, 613.}
\idxitem{— of permutations I, 474.}
\idxitem{— of substitutions I, 470.}
\idxitem{— of linear substitutions II, 153.}
\idxitem{Composition series of a group II, 17.}
\idxitem{Configuration of the ternary substitution group of degree 168 II, 451.}
\idxitem{Congruence of numbers I, 358, 368; II, 539.}
\idxitem{Congruences of the first degree I, 368.}
\idxitem{— in algebraic fields II, 542.}
\idxitem{Congruence group II, 250.}
\idxitem{Congruence field II, 245.}
\idxitem{— of the second degree II, 259.}
\idxitem{Congruence roots I, 426.}
\idxitem{Conjugate functions I, 505.}
\idxitem{— groups I, 506; II, 10.}
\idxitem{— fields I, 460; II, 493.}
\idxitem{— logarithms II, 672.}
\idxitem{— poles of a linear substitution group II, 201.}
\idxitem{Constancy of the index series II, 18.}
\idxitem{Contragredient groups II, 480.}
\idxitem{— transformations II, 157, 478.}
\idxitem{Contravariants II, 480.}
\idxitem{Covariants I, 187.}
\idxitem{— of ternary forms II, 327.}
\idxitem{— of ternary cubic forms II, 333.}
\idxitem{Cubic Abelian equations II, 107.}
\idxitem{— binary forms I, 192.}
\idxitem{— ternary forms II, 331.}
\idxitem{— equations I, 116, 522.}
\idxitem{— trigonometric solution I, 349.}
\idxitem{— cyclotomic fields II, 93.}
\idxitem{Algebraic curves II, 324.}
\idxitem{Cyclic functions I, 531, 538.}
\idxitem{— equations I, 538.}
\idxitem{— groups I, 476.}
\idxitem{— of linear substitutions II, 209.}
\idxitem{Cyclic groups in the congruence field II, 269, 275.}
\idxitem{— permutations I, 476, 493.}
\idxitem{Dedekind ideals II, 547.}
\idxitem{Dedekind's theorem on the field discriminant II, 578.}
\idxitem{Derived functions I, 51.}
\idxitem{Derivative of a product I, 54.}
\idxitem{— of functions of several variables I, 59.}
\idxitem{Determinants I, 68.}
\idxitem{— bordered I, 81.}
\idxitem{— from subdeterminants I, 97.}
\idxitem{Determinant of a system of linear equations I, 84.}
\idxitem{— of a quadratic form I, 180.}
\idxitem{Dense sets I, 4.}
\idxitem{Dihedral group II, 210.}
\idxitem{— in the congruence field II, 276.}
\idxitem{Differential quotients I, 54, 60.}
\idxitem{Differences I, 46.}
\idxitem{Difference product of the \(n\)-th roots of unity I, 423.}
\idxitem{Dirichlet's theorem on units II, 670.}
\idxitem{Discrete sets I, 4.}
\idxitem{Discriminant I, 150.}
\idxitem{— of the cubic form I, 37, 153, 193.}
\idxitem{— of the biquadratic form I, 155, 200.}
\idxitem{— of quadratic irrational numbers I, 378; II, 602.}
\idxitem{— of the cyclotomic equation for a prime degree I, 421.}
\idxitem{— of an algebraic curve II, 325.}
\idxitem{— of an algebraic field II, 529.}
\idxitem{— of the cyclotomic field II, 619, 643.}
\idxitem{Discriminants in an algebraic field II, 528.}
\idxitem{Discriminant surface I, 249.}
\idxitem{Division of numbers I, 1.}
\idxitem{— of integral functions I, 28.}
\idxitem{Divisors of an Abelian group II, 48.}
\idxitem{— of a group I, 476, 501; II, 7.}
\idxitem{Rotations II, 186.}
\idxitem{Rotation group II, 189.}
\idxitem{Triangular coordinates II, 322.}
\idxitem{Double points of a curve II, 326.}
\idxitem{Bitangents II, 327.}
\idxitem{— of a curve of fourth order II, 351.}
\idxitem{Intersection of groups I, 510; II, 9.}
\idxitem{Proper and improper equivalence I, 371.}
\idxitem{Proper and improper orthogonal substitutions II, 195.}
\idxitem{Simple groups I, 511; II, 136.}
\idxitem{— of degree 60 II, 138.}
\idxitem{— Abelian fields II, 649.}
\idxitem{Simplicity of the alternating group I, 600.}
\idxitem{— congruence groups II, 254.}
\idxitem{Identity of a group II, 5.}
\idxitem{Units I, 399.}
\idxitem{— in the field \(R(i)\) I, 586.}
\idxitem{— in algebraic fields II, 510.}
\idxitem{— in the real cyclotomic field II, 645.}
\idxitem{— in the field of eighth roots of unity II, 718.}
\idxitem{— independent II, 677.}
\idxitem{— fundamental system II, 682.}
\idxitem{Roots of unity I, 114, 408.}
\idxitem{— third I, 118.}
\idxitem{— primitive I, 411.}
\idxitem{— in cyclotomic fields II, 638.}
\idxitem{— in the congruence field II, 260.}
\idxitem{Elimination from linear equations I, 87.}
\idxitem{— from higher equations I, 157, 159.}
\idxitem{— from several equations I, 161.}
\idxitem{Final equation I, 157.}
\idxitem{Finite groups II, 4.}
\idxitem{— linear substitution groups II, 195.}
\idxitem{Finiteness of the invariant system of a linear group II, 169.}
\idxitem{Opposite elements of a group II, 5.}
\idxitem{Euler's theorem on homogeneous functions I, 62.}
\idxitem{Exponent system of units II, 679.}
\idxitem{Exponent system of numbers II, 685.}
\idxitem{Exponential functions II, 752.}
\idxitem{Factors of integral functions I, 39, 103.}
\idxitem{— of natural primes in algebraic fields II, 576.}
\idxitem{Fermat theorem I, 427; II, 55.}
\idxitem{— in the congruence field II, 247.}
\idxitem{— in algebraic fields II, 544.}
\idxitem{Forms I, 58.}
\idxitem{Form problem of the linear substitution group II, 175.}
\idxitem{Form system of the cubic form I, 196.}
\idxitem{— of the biquadratic form I, 203, 206.}
\idxitem{Integral functions I, 23.}
\idxitem{— decomposition into linear factors I, 103.}
\idxitem{— decomposition into prime functions I, 164, 452; II, 495.}
\idxitem{— homogeneous I, 56.}
\idxitem{— in a field I, 452; II, 520.}
\idxitem{Function congruences II, 243.}
\idxitem{Functional determinants I, 185.}
\idxitem{Functionals II, 499.}
\idxitem{Functional classes II, 552.}
\idxitem{Frobenius theorems on groups II, 129, 134.}
\idxitem{Fundamental theorem of algebra I, 121, 127, 296.}
\idxitem{Fundamental systems of units II, 681.}
\idxitem{— of normal units II, 731.}
\idxitem{Galois group I, 477, 511; II, 588.}
\idxitem{— resolvents I, 467.}
\idxitem{— field I, 465; II, 588.}
\idxitem{Integral functions of one variable I, 23.}
\idxitem{— of several variables I, 24.}
\idxitem{— in a field II, 520.}
\idxitem{— algebraic numbers II, 491.}
\idxitem{— functionals II, 504.}
\idxitem{Gauss sums I, 575.}
\idxitem{Fractional functions I, 32.}
\idxitem{Ordered sets I, 4.}
\idxitem{Weight of an invariant I, 187.}
\idxitem{Lattice II, 561.}
\idxitem{Equations I, 101.}
\idxitem{— linear homogeneous I, 81.}
\idxitem{— linear non-homogeneous I, 89.}
\idxitem{— of the third degree I, 116, 118, 522.}
\idxitem{— of the fourth degree I, 119, 201, 528.}
\idxitem{— of the fifth degree I, 233, 621; II, 404.}
\idxitem{— of the seventh degree II, 476, 481.}
\idxitem{Degree of a group I, 472, 475; II, 4.}
\idxitem{— of an element of a group II, 10.}
\idxitem{— of a permutation I, 500.}
\idxitem{— of a prime ideal and prime functional II, 516.}
\idxitem{Graffe's method of approximate solution of an equation I, 344.}
\idxitem{Bounds I, 121.}
\idxitem{Greatest common divisor of numbers I, 2.}
\idxitem{— of integral functions I, 34.}
\idxitem{— of groups I, 510.}
\idxitem{— of functionals II, 512, 519.}
\idxitem{Greatest normal divisor of a group II, 17.}
\idxitem{Fundamental curve of the group \(G_{168}\) II, 456.}
\idxitem{Fundamental forms of the polyhedron groups II, 205.}
\idxitem{— cyclic groups II, 210.}
\idxitem{— dihedral groups II, 210.}
\idxitem{— tetrahedron group II, 214.}
\idxitem{— octahedron group II, 216, 218.}
\idxitem{— icosahedron group II, 221, 230.}
\idxitem{Fundamental ideal II, 579, 593.}
\idxitem{Fundamental number of a field II, 529.}
\idxitem{Group (general) II, 3.}
\idxitem{— of the bitangent equation II, 380.}
\idxitem{— of the triple equations II, 344.}
\idxitem{— of the tetrahedron II, 212.}
\idxitem{— of the octahedron II, 216.}
\idxitem{— of the icosahedron II, 220.}
\idxitem{— of an equation I, 477.}
\idxitem{— of a normal field II, 588.}
\idxitem{— of the cyclotomic fields II, 68.}
\idxitem{— of an ideal in the normal field II, 589.}
\idxitem{— of the ideal classes II, 557.}
\idxitem{— of the 168th degree II, 282, 433.}
\idxitem{Groups, Abelian I, 476, 536; II, 6.}
\idxitem{— finite II, 4.}
\idxitem{— of permutations I, 475.}
\idxitem{— of substitutions I, 472.}
\idxitem{— of degree \(p^n\) II, 127.}
\idxitem{— of degree \(pq\) II, 141.}
\idxitem{— of linear substitutions II, 152.}
\idxitem{— of orthogonal substitutions II, 184.}
\idxitem{— of linear fractional substitutions II, 189, 195.}
\idxitem{Group characters II, 44.}
\idxitem{Group table II, 115.}
\idxitem{Half-metacyclic groups I, 616.}
\idxitem{Principal axes of a ternary linear substitution II, 441.}
\idxitem{— of the group \(G_{168}\) II, 444.}
\idxitem{Principal equation of the fifth degree I, 175, 230.}
\idxitem{Principal series II, 24.}
\idxitem{Principal subdeterminants I, 256.}
\idxitem{Hermite's solution of Sturm's problem I, 280.}
\idxitem{Hessian curve II, 330.}
\idxitem{— determinant I, 186.}
\idxitem{Hesse-Cayley designation of the bitangents II, 369.}
\idxitem{Hilbert theorem II, 165.}
\idxitem{Homogeneous functions I, 56.}
\idxitem{Ideals II, 547.}
\idxitem{Ideal numbers according to Kummer II, 551.}
\idxitem{Identical group I, 496.}
\idxitem{— substitution II, 152.}
\idxitem{Icosahedral equation II, 232, 418.}
\idxitem{Icosahedral group II, 220, 279.}
\idxitem{Icosahedral invariants II, 232.}
\idxitem{Icosahedral resolvents II, 422, 426.}
\idxitem{Imaginary numbers I, 17.}
\idxitem{Imaginary form of the linear congruence group II, 266.}
\idxitem{Imaginary congruence field of the second degree II, 259.}
\idxitem{Imaginary quadratic irrational numbers I, 379.}
\idxitem{Imprimitive groups I, 483, 516.}
\idxitem{— fields I, 460, 483.}
\idxitem{Index of an invariant II, 163.}
\idxitem{— of a divisor of a group I, 502; II, 8.}
\idxitem{Index moduli II, 61.}
\idxitem{Index series of a group II, 17.}
\idxitem{Indices I, 431.}
\idxitem{— of the elements of a group II, 44.}
\idxitem{— modulo an odd prime power II, 54.}
\idxitem{— modulo a power of 2 II, 58.}
\idxitem{Invariants I, 186.}
\idxitem{— of a finite linear substitution group II, 161.}
\idxitem{— absolute and relative II, 162, 177.}
\idxitem{— in the wider sense II, 179.}
\idxitem{— of ternary forms II, 328.}
\idxitem{— of curves of third order II, 337.}
\idxitem{— of the group \(G_{168}\) II, 453, 461.}
\idxitem{Invariant curves II, 454.}
\idxitem{Inverse substitution II, 155.}
\idxitem{Irrational ratios I, 10.}
\idxitem{— numbers I, 12, 370.}
\idxitem{Irreducibility I, 453.}
\idxitem{— of pure equations I, 607.}
\idxitem{— of the cyclotomic equation I, 417; II, 768.}
\idxitem{Isomorphic groups I, 477; II, 6.}
\idxitem{Isomorphism (multistage) II, 15.}
\idxitem{Jacobi's estimate of the number of roots between two bounds I, 311.}
\idxitem{Jordan's theorem on composite groups II, 18.}
\idxitem{Continued fractions I, 358.}
\idxitem{— for equivalent numbers I, 374.}
\idxitem{Chains of principal subdeterminants I, 257.}
\idxitem{— Sturm chains I, 272.}
\idxitem{Least common multiple of numbers I, 3.}
\idxitem{— of functionals II, 519.}
\idxitem{— of groups II, 29.}
\idxitem{Field I, 449; II, 493.}
\idxitem{Spherical functions I, 273.}
\idxitem{Kummer theorem on resolvents II, 632.}
\idxitem{Cyclotomic equation I, 422, 554.}
\idxitem{Cyclotomic field II, 67.}
\idxitem{— with given group II, 79.}
\idxitem{— of prime-power degree II, 617.}
\idxitem{Cyclotomy I, 408, 554.}
\idxitem{Crystallographic groups II, 241.}
\idxitem{Lagrange theorem on functions admitting the permutations of a group I, 508.}
\idxitem{Laguerre theorems for equations with only real roots I, 322.}
\idxitem{Legendre symbol I, 441.}
\idxitem{Linear congruence groups I, 611; II, 250, 302.}
\idxitem{— homogeneous II, 303.}
\idxitem{— real II, 261.}
\idxitem{— imaginary form II, 266.}
\idxitem{— of degree 168 II, 282.}
\idxitem{— binary for modulus 3 II, 305.}
\idxitem{— ternary for modulus 2 II, 306.}
\idxitem{Linear fractional substitutions II, 189.}
\idxitem{Linear functions I, 30.}
\idxitem{Linear equations I, 81.}
\idxitem{Linear transformation I, 178.}
\idxitem{Linear substitutions I, 178, 371; II, 151.}
\idxitem{Lüroth theorem II, 404.}
\idxitem{Manifold I, 4.}
\idxitem{Matrix I, 82.}
\idxitem{Set I, 4.}
\idxitem{Measurable set I, 8.}
\idxitem{Metacyclic equations I, 597; II, 292.}
\idxitem{— of prime degree I, 609.}
\idxitem{— of the fifth degree I, 621, 648.}
\idxitem{— of the sixth degree II, 296.}
\idxitem{— of the eighth degree II, 314.}
\idxitem{— of the ninth degree II, 305.}
\idxitem{Metacyclic groups I, 598; II, 27.}
\idxitem{— functions I, 617, 635.}
\idxitem{Minimal basis of a field II, 529.}
\idxitem{— of the cyclotomic field II, 620.}
\idxitem{— of a quadratic field II, 603.}
\idxitem{Minimum I, 122.}
\idxitem{— of a quadratic form II, 559.}
\idxitem{— of the fundamental number of a field II, 569.}
\idxitem{Modulus I, 358.}
\idxitem{Multiplicative substitutions II, 152.}
\idxitem{Multiplication of numbers I, 14.}
\idxitem{— of determinants I, 92.}
\idxitem{— of trigonometric functions I, 433.}
\idxitem{Convergents of a continued fraction I, 362.}
\idxitem{Approximation method for computing roots by regula falsi I, 331.}
\idxitem{— for computing roots by Daniel Bernoulli I, 341.}
\idxitem{— Graffe I, 344.}
\idxitem{— Newton I, 335.}
\idxitem{— by continued fractions I, 401.}
\idxitem{Approximation method for solving trinomial equations I, 352.}
\idxitem{Natural irrationalities I, 516.}
\idxitem{Cosets I, 502; II, 8.}
\idxitem{Negative numbers I, 16.}
\idxitem{Ninth roots of unity I, 582.}
\idxitem{Newton formulas for the power sums I, 142.}
\idxitem{Newton rule for estimating roots I, 305.}
\idxitem{Non-metacyclic equations I, 603.}
\idxitem{Norm I, 461; II, 494, 498.}
\idxitem{— absolute II, 502, 538.}
\idxitem{— of Gaussian complex numbers I, 586.}
\idxitem{— of an ideal II, 551.}
\idxitem{— of a field I, 465.}
\idxitem{— of a quadratic irrational number I, 377.}
\idxitem{Normal units in the cyclotomic field II, 729.}
\idxitem{Normal forms of the linear substitution groups II, 181.}
\idxitem{Normal equation I, 465.}
\idxitem{Normal field I, 464; II, 588.}
\idxitem{Normal divisor of a group I, 511; II, 11.}
\idxitem{Zero I, 16.}
\idxitem{Upper bound for the roots I, 316.}
\idxitem{Octahedron group II, 216.}
\idxitem{— in the congruence field II, 277.}
\idxitem{Orthogonal substitution II, 184.}
\idxitem{Partial discriminant II, 597.}
\idxitem{Partial fundamental ideal II, 585, 597.}
\idxitem{Partial norm II, 584.}
\idxitem{Partial trace II, 585.}
\idxitem{Partial resolvents I, 512.}
\idxitem{Pell equation I, 395.}
\idxitem{Period of a permutation I, 500.}
\idxitem{Periods of cyclotomy I, 557, 579; II, 74.}
\idxitem{— of the roots of a cyclic equation I, 545.}
\idxitem{— of the reduced numbers I, 388.}
\idxitem{Periodic continued fractions I, 387.}
\idxitem{Permutations I, 64, 473, 500.}
\idxitem{— their analytic representation I, 613; II, 299.}
\idxitem{— of first and second kind I, 65, 496.}
\idxitem{Permutation groups I, 475, 489; II, 117.}
\idxitem{— of degree 168 on seven letters II, 473.}
\idxitem{Polars I, 188.}
\idxitem{Poles of linear fractional substitutions II, 196.}
\idxitem{— of a group II, 198.}
\idxitem{— in the congruence field II, 274.}
\idxitem{— of a ternary group II, 438.}
\idxitem{Polyhedron groups II, 202.}
\idxitem{— of the second kind II, 234.}
\idxitem{— in the congruence field II, 273.}
\idxitem{Polynomial coefficients I, 51.}
\idxitem{Polynomial theorem I, 50.}
\idxitem{Positive units in the cyclotomic field II, 742.}
\idxitem{Power residues I, 430.}
\idxitem{Power sums I, 140.}
\idxitem{Primary divisors of a cyclotomic field II, 71.}
\idxitem{— of an Abelian group II, 72.}
\idxitem{Prime factors of the numbers of an algebraic field II, 523.}
\idxitem{— of natural prime numbers II, 576.}
\idxitem{— of cyclotomic resolvents II, 635.}
\idxitem{Prime functionals II, 515, 576.}
\idxitem{— relative II, 514.}
\idxitem{Prime ideals II, 550.}
\idxitem{— in the quadratic field II, 604.}
\idxitem{— in the relatively normal field II, 588.}
\idxitem{— in the divisors of a normal field II, 599.}
\idxitem{— in the cyclotomic fields II, 619, 622.}
\idxitem{— in the real cyclotomic field II, 643.}
\idxitem{Primitive congruence roots I, 428; II, 546.}
\idxitem{— roots of unity I, 410.}
\idxitem{— functions I, 25; II, 497.}
\idxitem{— roots of prime numbers I, 429.}
\idxitem{— of squares of primes II, 57.}
\idxitem{— of prime functionals II, 546.}
\idxitem{— of prime ideals II, 591.}
\idxitem{— of a congruence field II, 248.}
\idxitem{Primitive and imprimitive characters in cyclotomic fields II, 720.}
\idxitem{Primitive groups I, 484.}
\idxitem{Primitive fields I, 463.}
\idxitem{Primitive form problems II, 180.}
\idxitem{Prime numbers I, 1.}
\idxitem{— relative I, 359.}
\idxitem{— in the field \(R(i)\) I, 587.}
\idxitem{— in the field of third roots of unity I, 593.}
\idxitem{Squares in the congruence field II, 248.}
\idxitem{Quadratic forms I, 179.}
\idxitem{— equations I, 116.}
\idxitem{— irrational numbers I, 377.}
\idxitem{— residues I, 444.}
\idxitem{— fields II, 601.}
\idxitem{Squaring of the circle II, 760.}
\idxitem{Rational roots of an equation I, 403.}
\idxitem{Rational numbers I, 12.}
\idxitem{Rational ratio I, 10.}
\idxitem{Domain of rationality I, 452.}
\idxitem{Space of \(n\) dimensions II, 560.}
\idxitem{Reality of roots I, 241.}
\idxitem{— for quadratic and cubic equations I, 243.}
\idxitem{— for biquadratic equations I, 246.}
\idxitem{— for cyclic equations I, 551.}
\idxitem{— for metacyclic equations I, 647.}
\idxitem{— of the bitangents of a curve of fourth order II, 391.}
\idxitem{Reality conditions of the orthogonal group II, 193.}
\idxitem{Reality relations of the inflection points of a curve of third order II, 340.}
\idxitem{— for triple equations II, 348.}
\idxitem{Operations of calculation I, 1.}
\idxitem{Calculating with integral functions I, 24.}
\idxitem{— with numbers I, 14.}
\idxitem{Reciprocity law of quadratic residues I, 443.}
\idxitem{Reducible and irreducible functions I, 453.}
\idxitem{— modulo a prime modulus II, 243.}
\idxitem{Reducible and irreducible equations I, 404, 482.}
\idxitem{Reduced numbers I, 380, 384.}
\idxitem{— units II, 680.}
\idxitem{Real cyclotomic fields II, 641.}
\idxitem{Real radicals I, 606.}
\idxitem{Regula falsi I, 334.}
\idxitem{Regular fields II, 188.}
\idxitem{Regulator of a system of units II, 678, 729.}
\idxitem{— of the field II, 682.}
\idxitem{Pure equations I, 109.}
\idxitem{Series II, 694.}
\idxitem{Relatively prime numbers I, 2.}
\idxitem{— prime functionals II, 514.}
\idxitem{Relative norm II, 584.}
\idxitem{Relative trace II, 585.}
\idxitem{Resolvents I, 511.}
\idxitem{— of the biquadratic equation I, 120.}
\idxitem{Resolvents of the composition series II, 289.}
\idxitem{— of equations of seventh degree II, 476.}
\idxitem{— of equations of eighth degree II, 308.}
\idxitem{— of the group \(G_{168}\) II, 466, 473.}
\idxitem{— of the icosahedral equation II, 419.}
\idxitem{— in cyclotomy I, 564; II, 63, 651.}
\idxitem{— of Lagrange I, 542.}
\idxitem{— for metacyclic equations I, 631.}
\idxitem{— with a parameter II, 404.}
\idxitem{Remainder of division by integral functions I, 28.}
\idxitem{Residue system I, 359; II, 541.}
\idxitem{Resultants I, 156.}
\idxitem{Resultant of two quadratic functions I, 36, 158.}
\idxitem{Rolle theorem on the number of real roots I, 319.}
\idxitem{Secular equation I, 276.}
\idxitem{Complete quotients of a continued fraction I, 361.}
\idxitem{Cut I, 5.}
\idxitem{Seven-systems of bitangents II, 367.}
\idxitem{Seventeen-gon I, 568.}
\idxitem{Singular points of a curve II, 325.}
\idxitem{Trace I, 461; II, 494.}
\idxitem{Steiner complexes of bitangents II, 357.}
\idxitem{Continuity I, 5.}
\idxitem{— of integral functions I, 105.}
\idxitem{— of roots I, 132.}
\idxitem{Sturm functions I, 279.}
\idxitem{— chains I, 271.}
\idxitem{Sturm problem I, 270.}
\idxitem{Linear substitution I, 92, 372; II, 151.}
\idxitem{— of ratios II, 158.}
\idxitem{Substitutions of a normal field I, 468; II, 588.}
\idxitem{Substitution determinant I, 92, 178.}
\idxitem{Sylow theorems II, 121, 125.}
\idxitem{Symmetric determinants I, 69.}
\idxitem{— functions I, 138, 144, 147.}
\idxitem{— fundamental functions I, 139.}
\idxitem{— group I, 492.}
\idxitem{Symmetric group, its normal divisors I, 602.}
\idxitem{Syzygetic and azygetic systems of bitangents II, 362.}
\idxitem{— complexes of bitangents II, 366.}
\idxitem{Tangents of a curve II, 326.}
\idxitem{Taylor expansion I, 59.}
\idxitem{Ternary forms II, 322.}
\idxitem{— linear groups of degree 168 II, 433.}
\idxitem{— congruence group of degree 168 II, 475.}
\idxitem{Tetrahedron group II, 212.}
\idxitem{— in the congruence field II, 276.}
\idxitem{Divisibility of integral functions I, 32.}
\idxitem{— of integral functionals II, 509.}
\idxitem{— of whole numbers II, 510.}
\idxitem{Divisors of whole numbers I, 1.}
\idxitem{— of whole numbers, greatest common I, 359.}
\idxitem{— of integral functions I, 27; II, 497.}
\idxitem{— of integral functions, greatest common I, 34.}
\idxitem{— of a group I, 476, 501; II, 7.}
\idxitem{— of the icosahedral group II, 227.}
\idxitem{— of a field I, 450.}
\idxitem{— of the linear congruence group II, 271.}
\idxitem{Relatively prime numbers I, 2, 359.}
\idxitem{— functions I, 35.}
\idxitem{— functionals II, 514.}
\idxitem{Partial denominator of a continued fraction I, 361.}
\idxitem{Division of the angle I, 435.}
\idxitem{Total resolvents I, 512.}
\idxitem{Law of inertia of quadratic forms I, 183, 255.}
\idxitem{Transcendental functions for solving the icosahedral equation II, 429.}
\idxitem{Transcendental numbers II, 745.}
\idxitem{Transcendence of the number \(e\) II, 751.}
\idxitem{— of the number \(\pi\) II, 756.}
\idxitem{Transformation of a group I, 506.}
\idxitem{— of the quadratic form into a sum of squares I, 181.}
\idxitem{— of forms of the \(n\)-th degree I, 185.}
\idxitem{— of the cubic equation I, 219.}
\idxitem{Transformation of the equation of the fifth degree I, 230.}
\idxitem{Transformed substitutions II, 156.}
\idxitem{Transitive and intransitive groups I, 482, 506.}
\idxitem{— permutation group of index 6 on six letters I, 628.}
\idxitem{Transposed substitutions II, 156.}
\idxitem{Transpositions I, 65, 492.}
\idxitem{Trigonometric solution of pure equations I, 111.}
\idxitem{— of cubic equations I, 349.}
\idxitem{Trinomial equations I, 352.}
\idxitem{Triple equations II, 342.}
\idxitem{Triple systems II, 310.}
\idxitem{Tschirnhausen transformation I, 170, 210.}
\idxitem{— of the cubic equation I, 219.}
\idxitem{Invertible periods I, 391.}
\idxitem{Unknown I, 20.}
\idxitem{Indeterminate equations I, 365.}
\idxitem{Infinite roots of an equation I, 136.}
\idxitem{Infinity of integral rational functions I, 104.}
\idxitem{Unicursal curves II, 416.}
\idxitem{Subdeterminants I, 72.}
\idxitem{— higher I, 77.}
\idxitem{— complementary I, 79.}
\idxitem{Uncountable sets II, 748.}
\idxitem{Original functions I, 25.}
\idxitem{Variable in the theory of fields II, 499.}
\idxitem{Ratios I, 10.}
\idxitem{Related numbers and functionals II, 631.}
\idxitem{Complete residue system I, 359; II, 541.}
\idxitem{Complete systems of bitangents II, 367.}
\idxitem{Volume II, 561, 690.}
\idxitem{Sign change of integral functions I, 107.}
\idxitem{Inflection points II, 326.}
\idxitem{Inflection tangents II, 326.}
\idxitem{— of a curve of third order II, 331.}
\idxitem{Wilson theorem I, 428.}
\idxitem{Division of the angle I, 551, 609.}
\idxitem{Duplication of the cube I, 609.}
\idxitem{Roots of equations I, 101.}
\idxitem{— of equations of odd degree I, 109.}
\idxitem{— of pure equations I, 109.}
\idxitem{— of metacyclic equations I, 638.}
\idxitem{— rational I, 403.}
\idxitem{Number I, 1, 12.}
\idxitem{— of the terms of a homogeneous function I, 58.}
\idxitem{— of permutations I, 64.}
\idxitem{Number classes modulo a composite modulus II, 60.}
\idxitem{Number sequences I, 15.}
\idxitem{Number field I, 449; II, 527.}
\idxitem{Sign change and sign sequence I, 262.}
\idxitem{Reducible and irreducible functions of several variables I, 164.}
\idxitem{Decomposition of integral functions into linear factors I, 102.}
\idxitem{— of groups into cosets I, 502; II, 8.}
\idxitem{— according to two divisors II, 120.}
\idxitem{Composite Abelian fields II, 649.}
\idxitem{Composition of linear substitutions I, 372; II, 153.}
\idxitem{— of substitutions of a normal field I, 470.}
\idxitem{— of permutations I, 473.}
\idxitem{Two-sided numbers I, 390.}
\end{multicols}
\clearpage

\section*{Corrections to the First Volume.}
\begin{longtable}{p{0.26\textwidth}p{0.66\textwidth}}
Page 93, line 6 from below & read $B_n^{(0)}$ instead of $B^{(0)}$.\\
117, formula (10) & read $\sqrt{\frac{p}{2}+\sqrt{R}}$ instead of $\sqrt{-\frac{p}{2}+\sqrt{R}}$.\\
158, formula (14), right side & interchange $m$ and $n$.\\
172, formula (14) & read $n-1$ instead of $n+1$.\\
185, line 9 from above & read $\psi'$ instead of $\psi$.\\
190, line 11 from below & read $a'_1$ instead of $(n-1)a'_1$.\\
194, line 7 from below & read $+9$ instead of $-9$.\\
200, line 7 from below & read §. 61 instead of §. 67.\\
202, line 3 from below & read $u$ instead of $v_i$ in both formulae.\\
213, line 6 from below & read $\Phi(r,k)$ instead of $\Phi(r,k)$.\\
221, line 12 from below & read $\frac{2}{9}$ instead of $\frac{2}{3}$.\\
222, lines 4, 5 from below & read $+q_0-q_0$ instead of $-q_0+q_0$.\\
233, lines 4, 7 from above & read $\frac{1}{5}D$ instead of $5D$.\\
236, line 12 from above & read $5c^2$ instead of $-5c^2$.\\
251, line 10 from above & read $-\sqrt{3}$ instead of $\sqrt{3}$.\\
268, line 5 from below & read $\frac{1}{5}D$ instead of $D$.\\
282, formula (2) & read $f_5$ instead of $f$.\\
318, line 2 from below & read $F(z)$ instead of $F(x)$.\\
332, formula (5) & read $f(\beta)-$ instead of $f(\beta)+$.\\
334, line 6 from below & read $0.00164$ instead of $0.0164$.\\
344, line 6 from below, formula & read $a_n>\varphi(a_n)$ and $\varphi(a_n)$ instead of $a_m\varphi(a_m)$ and $\varphi(a_m)$.\\
357, line 3 from below & read $6$ instead of $5$.\\
378, 379, formulae (13), (14) & replace $a_1,b_1,c_1$ by $\pm a_1,\pm b_1,\pm c_1$, and choose the sign so that $q$ becomes positive.\\
402, lines 6, 7 from below & interchange $a_4=1$ and $a_3=2$.\\
---, line 4 from below & read $9083$ instead of $9183$.\\
419 & compare with the first addendum to the second volume.\\
423, line 15 from above & read $\frac{(n-1)(n-2)}{2}$ instead of $\frac{n-1}{2}$.\\
---, line 24 from above & read $\frac12\left(\frac{(n-1)(n-2)}{2}-\frac{n-1}{2}\right)=\frac{(n-1)(n-3)}{4}\equiv0\pmod2$.\\
423, line 3 from below & read $(-1)^{\frac{(n-1)(n-3)}{4}}=+1$ instead of $(-1)^{\frac{n-1}{2}}$.\\
424, formula (14) & read $P=i^{\frac{n-1}{2}}\,n^{\frac{n-3}{2}}\sqrt n$.\\
436, lines 3, 22 from above & read $\frac{m}{n}$ instead of $n$.\\
482, line 8 from below & interchange transitive and intransitive.\\
502, line 18 from above & read $n_1$ instead of $n$.\\
504, line 9 from below & read $k_{q-1}$ instead of $k_{\nu-1}$.\\
514, line 5 from below & read $\mathfrak g_1$ instead of $\mathfrak G_1$.\\
531, line 11 from above & read $y_1$ instead of $y'_1$.\\
540, formula (10) & read $\gamma e-1$ instead of $\gamma e_1-1$.\\
544, line 15 from above & read $mB_k=\lambda^{-1}(\varepsilon,a)^k(\varepsilon h,a)^k(\varepsilon h^2,a)^k\cdots$.\\
---, line 23 from above & read $e^2$ instead of $e$.\\
551, line 15 from above & read $a^k$ instead of $(ck$.\\
592, line 8 from below & read: which is fulfilled only for $b=0,a=\pm1$, or $a=0,b=\pm1$, or $a=b=\pm1$.\\
612, line 12 from above & read $1-a$ instead of $a-b$.\\
618, line 1 from below & read $c_1$ instead of $C_1$.\\
623, line 6 from below & read $a_6$ instead of $a_G$.\\
636, lines 18, 21 from above & read $-g^{n-1}+1$ instead of $g^{n-1}-1$.\\
639, lines 10, 27 from above & read $n-2,n-3$ instead of $n-1,n-2$.\\
641, line 1 from above & read (6) instead of (5).\\
646, line 11 from above & read rational in $R(t_0)$ instead of rational.\\
---, lines 14, 19, 22, 26 from above & read $R(t_0)$ instead of $R$.\\
---, lines 17, 18 from above & delete beginning with ``All quantities''.\\
---, lines 1, 2 from below & delete beginning with ``the others''.\\
653, line 4 from above & read $-A_2rq$ instead of $+A_2rq$.\\
---, line 5 from above & read $+A_2q'$ instead of $-A_2q'$.\\
\end{longtable}

\section*{Corrections to the Second Volume.}
\begin{longtable}{p{0.26\textwidth}p{0.66\textwidth}}
Page 15, line 19 from above & read $N$ instead of $P$.\\
16, lines 12, 13 from above & interchange $P$ and $Q$.\\
21, line 19 from above & read (17) instead of (13).\\
26, line 17 from above & read §. 6, 4 instead of §. 6, 2.\\
569, line 5 from below & in the exponent read $n-r$ instead of $\mu-r$.\\
592, lines 6, 3 from below & read $\gamma^{p^{-1}}$ instead of $\gamma^{p^{r-1}}$.\\
\end{longtable}
\end{document}
